% Réponse immunitaire — SVTEEHB 3ème — Cours signé OnBuch+
% Programme officiel MINESEC. Compiler avec : tectonic N06-reponse-immunitaire.tex
\newcommand{\DOCMATIERE}{SVTEEHB}
\newcommand{\DOCNIVEAU}{3ème}
\newcommand{\DOCMODULE}{SVTEEHB – classe de 3ème}
\newcommand{\DOCLECON}{N06}
\newcommand{\DOCTITRE}{Réponse immunitaire}
% OnBuch+ — préambule « cours signés » ENRICHI (compilé avec tectonic / XeLaTeX)
% Système visuel « Pop » d'OnBuch+ : fond crème (jamais blanc), bordures encre
% épaisses, ombres « dures » décalées, titres Archivo Black en capitales.
% Version MATHÉMATIQUES : graphes (pgfplots/tikz), boîtes pédagogiques
% (définition, propriété, méthode, attention, le savais-tu, exercice résolu,
% exercice, TP, bilan), page de garde et signature OnBuch+ ancrée en bas.
%
% Macros attendues dans main.tex AVANT \input{preamble} :
%   \DOCMATIERE  \DOCNIVEAU  \DOCMODULE  \DOCLECON (numéro)  \DOCTITRE
\documentclass[11pt]{article}

\usepackage{fontspec}
\usepackage[french]{babel}
\usepackage[a4paper,margin=2cm,top=3.4cm,bottom=2.8cm,headheight=1.4cm,footskip=1.3cm]{geometry}
\usepackage{amsmath,amssymb,amsfonts}
\usepackage{mathtools} % \xmapsto, \coloneqq... souvent utilisés par les modèles
\usepackage{cancel} % simplifications barrées (bilans de forces)
% \abs{z} (module, valeur absolue) et \norm{u} (norme), utilisés par les modèles
\providecommand{\abs}{}\let\abs\relax\DeclarePairedDelimiter{\abs}{\lvert}{\rvert}
\providecommand{\norm}{}\let\norm\relax\DeclarePairedDelimiter{\norm}{\lVert}{\rVert}
\usepackage[table]{xcolor} % [table] : fournit aussi \rowcolors (alternance de lignes)
\usepackage{pagecolor}
\usepackage{tikz}
\usetikzlibrary{arrows.meta,positioning,calc,shapes.geometric,shapes.misc,shapes.arrows,decorations.pathmorphing,decorations.pathreplacing,decorations.markings,patterns,shadows,mindmap,trees,fit,backgrounds,matrix,intersections,angles,quotes}
\usepackage{pgfplots}
\usepackage[version=4]{mhchem} % équations nucléaires \ce{^{235}_{92}U}
\usepackage[european,siunitx]{circuitikz} % schémas électriques
\pgfplotsset{compat=1.18}
\usepgfplotslibrary{fillbetween,groupplots} % aires entre courbes (calcul intégral)
\usepackage{enumitem}
\usepackage{fancyhdr}
% [most] charge toutes les bibliothèques tcolorbox (listings, minted, raster,
% documentation...) et gonfle inutilement la mémoire TeX sur un document long
% (~300 boîtes) : on ne charge que ce qui est réellement utilisé.
\usepackage[skins,breakable]{tcolorbox}
\usepackage{graphicx}
\usepackage{colortbl}
\usepackage{booktabs}
\usepackage{tabularx}
\usepackage{multirow}
\usepackage{diagbox} % case d'en-tête diagonale des tableaux à double entrée
% Colonnes extensibles centrée / gauche / droite (C, L, R), souvent utilisées par les modèles
\newcolumntype{C}{>{\centering\arraybackslash}X}
\newcolumntype{L}{>{\raggedright\arraybackslash}X}
\newcolumntype{R}{>{\raggedleft\arraybackslash}X}
\usepackage{array}
\usepackage{multicol}
\usepackage{siunitx}
% parse-numbers=false : accepte aussi \SI{2,5 \times 10^{-3}}{...}
\sisetup{locale=FR,output-decimal-marker={,},group-minimum-digits=4,parse-numbers=false}
\DeclareSIUnit\ohm{\ensuremath{\symup{\Omega}}} % U+2126 absent de la police
\DeclareSIUnit\division{div} % réglages d'oscilloscope (ms/div, V/div)
\DeclareSIUnit\tr{tr} % tours (vitesse de rotation, ex. \SI{1500}{\tr\per\minute})
\DeclareSIUnit\molar{M} % concentration molaire (ex. \SI{3}{\molar}, \SI{1}{\micro\molar})
\DeclareSIUnit\an{an} % année (le modèle confond parfois avec \year, un compteur TeX) — ex. \SI{5}{\kilo\meter\per\an}
\DeclareSIUnit\FCFA{FCFA} % franc CFA (ex. \SI{500}{\FCFA\per\kilo\gram})
\DeclareSIUnit\degreeBrix{\textdegree Brix} % teneur en sucre d'un sirop/jus (réfractométrie)
\DeclareSIUnit\month{mois} % le modèle utilise parfois \month, un compteur TeX, comme unité de durée
\DeclareSIUnit\microliter{\micro\liter} % le modèle écrit parfois \microliter en un seul token
\DeclareSIUnit\million{millions} % dénombrement (ex. \SI{15}{\million\per\mL} spermatozoïdes)
\usepackage{qrcode}
% \degree : commande fréquemment utilisée par les modèles pour un angle ou une
% température en dehors de \SI{}{} (non fournie par les paquets chargés).
\providecommand{\degree}{^{\circ}}
% \qed (fin de démonstration), utilisé par les modèles sans amsthm
\providecommand{\qed}{\hfill$\square$}
% \barre : double barre des valeurs interdites dans un tableau de variations (tkz-tab)
\providecommand{\barre}{\ensuremath{\|}}
% environnement proof (amsthm non chargé) utilisé par les modèles
\ifdefined\proof\else\newenvironment{proof}[1][Démonstration]{\par\noindent\textit{#1.}\ }{\qed\par}\fi
\usepackage{lastpage}
\usepackage{hyperref}
\hypersetup{hidelinks}

% ============================================================
% Palette « Pop » OnBuch+ — mêmes hex que la classe P (plus_ui.dart)
% ============================================================
\definecolor{poporange}{HTML}{F59321}
\definecolor{popdark}{HTML}{E07A0C}
\definecolor{popcream}{HTML}{FFF6E8}
\definecolor{popink}{HTML}{1C1714}
\definecolor{poppurple}{HTML}{7A5AE0}
\definecolor{popgreen}{HTML}{2E9E6A}
\definecolor{poppink}{HTML}{E0457B}
\definecolor{popblue}{HTML}{2F7DE1}
\definecolor{popgold}{HTML}{C98A12}
\definecolor{popmuted}{HTML}{8A7F76}
\definecolor{popline}{HTML}{EADFD2}
\definecolor{popgray}{HTML}{8A7F76}
\colorlet{popgrey}{popgray}
% Alias tolérants (le modèle nomme parfois les couleurs différemment)
\colorlet{popwhite}{white}
\colorlet{popblack}{popink}
\colorlet{popred}{poppink}
\colorlet{popyellow}{popgold}
% Teintes claires (fonds de tableaux/encadrés) — le modèle nomme parfois
% les couleurs avec un suffixe L (ex. popblueL) sans les avoir définies.
\colorlet{poporangeL}{poporange!20}
\colorlet{popdarkL}{popdark!20}
\colorlet{poppurpleL}{poppurple!20}
\colorlet{popgreenL}{popgreen!20}
\colorlet{poppinkL}{poppink!20}
\colorlet{popblueL}{popblue!20}
\colorlet{popgoldL}{popgold!20}
\colorlet{popredL}{popred!20}
\colorlet{popgrayL}{popgray!20}
% Teintes claires pour fonds de tableaux / zones
\colorlet{popblueL}{popblue!10!white}
\colorlet{poporangeL}{poporange!14!white}
\colorlet{popgreenL}{popgreen!12!white}
\colorlet{poppurpleL}{poppurple!12!white}
\colorlet{poppinkL}{poppink!12!white}
\colorlet{popgoldL}{popgold!14!white}

\pagecolor{popcream}

% ============================================================
% Typographie
% ============================================================
\defaultfontfeatures{Ligatures=TeX}
\setmainfont{PlusJakartaSans-Medium.ttf}[Path=../../../fonts/,BoldFont=PlusJakartaSans-Bold.ttf]
% Maths en sans-serif (Fira Math) avec chiffres et lettres droites de Plus
% Jakarta, pour que les formules se fondent dans le texte.
\usepackage[math-style=ISO,bold-style=ISO]{unicode-math}
\setmathfont{FiraMath-Regular.otf}[Path=../../../fonts/]
\setmathfont{PlusJakartaSans-Medium.ttf}[Path=../../../fonts/,range={up/num,up/{Latin,latin},\mathup}]
\usepackage{icomma} % virgule décimale sans espace en mode math
% Caractères mathématiques Unicode absents des polices de texte (Plus Jakarta,
% Archivo Black) : rendus par leur équivalent mathématique, en texte comme en
% formule — sinon ils s'affichent en carré vide (ex. « Arithmétique dans ℤ »).
\usepackage{newunicodechar}
\newunicodechar{ℝ}{\ensuremath{\mathbb{R}}}
\newunicodechar{ℤ}{\ensuremath{\mathbb{Z}}}
\newunicodechar{ℂ}{\ensuremath{\mathbb{C}}}
\newunicodechar{ℕ}{\ensuremath{\mathbb{N}}}
\newunicodechar{ℚ}{\ensuremath{\mathbb{Q}}}
\newunicodechar{𝔻}{\ensuremath{\mathbb{D}}}
\newunicodechar{α}{\ensuremath{\alpha}}
\newunicodechar{β}{\ensuremath{\beta}}
\newunicodechar{γ}{\ensuremath{\gamma}}
\newunicodechar{δ}{\ensuremath{\delta}}
\newunicodechar{ε}{\ensuremath{\varepsilon}}
\newunicodechar{θ}{\ensuremath{\theta}}
\newunicodechar{λ}{\ensuremath{\lambda}}
\newunicodechar{μ}{\ensuremath{\mu}}
\newunicodechar{σ}{\ensuremath{\sigma}}
\newunicodechar{τ}{\ensuremath{\tau}}
\newunicodechar{φ}{\ensuremath{\varphi}}
\newunicodechar{ψ}{\ensuremath{\psi}}
\newunicodechar{ω}{\ensuremath{\omega}}
\newunicodechar{Σ}{\ensuremath{\Sigma}}
% majuscules produites par \MakeUppercase dans les titres (α-aminés -> Α)
\newunicodechar{Α}{\ensuremath{\alpha}}
\newunicodechar{Β}{\ensuremath{\beta}}
\newunicodechar{Γ}{\ensuremath{\Gamma}}
\newunicodechar{Δ}{\ensuremath{\Delta}}
\newunicodechar{Ε}{\ensuremath{\varepsilon}}
\newunicodechar{Θ}{\ensuremath{\Theta}}
\newunicodechar{Λ}{\ensuremath{\Lambda}}
\newunicodechar{Μ}{\ensuremath{\mu}}
\newunicodechar{Τ}{\ensuremath{\tau}}
\newunicodechar{Φ}{\ensuremath{\Phi}}
\newunicodechar{Ψ}{\ensuremath{\Psi}}
\newunicodechar{Ω}{\ensuremath{\Omega}}
\newunicodechar{↦}{\ensuremath{\mapsto}}
\newunicodechar{∈}{\ensuremath{\in}}
\newunicodechar{∉}{\ensuremath{\notin}}
\newunicodechar{∧}{\ensuremath{\wedge}}
\newunicodechar{⇔}{\ensuremath{\Leftrightarrow}}
\newunicodechar{⟺}{\ensuremath{\Longleftrightarrow}}
\newunicodechar{⇒}{\ensuremath{\Rightarrow}}
\newunicodechar{⟹}{\ensuremath{\Longrightarrow}}
\newunicodechar{∘}{\ensuremath{\circ}}
\newunicodechar{∪}{\ensuremath{\cup}}
\newunicodechar{∩}{\ensuremath{\cap}}
\newunicodechar{≡}{\ensuremath{\equiv}}
\newunicodechar{⊂}{\ensuremath{\subset}}
\newunicodechar{⊥}{\ensuremath{\perp}}
\newunicodechar{∃}{\ensuremath{\exists}}
\newunicodechar{∀}{\ensuremath{\forall}}
\newunicodechar{∛}{\ensuremath{\sqrt[3]{\,}}}
\newunicodechar{∜}{\ensuremath{\sqrt[4]{\,}}}
\newunicodechar{⃗}{\ensuremath{{}^{\to}}}
\newunicodechar{ⁿ}{\ifmmode^{n}\else\textsuperscript{\ensuremath{n}}\fi}
\newunicodechar{ˣ}{\ifmmode^{x}\else\textsuperscript{\ensuremath{x}}\fi}
\newunicodechar{ᵇ}{\ifmmode^{b}\else\textsuperscript{\ensuremath{b}}\fi}
\newunicodechar{ʳ}{\ifmmode^{r}\else\textsuperscript{\ensuremath{r}}\fi}
\newunicodechar{ᵘ}{\ifmmode^{u}\else\textsuperscript{\ensuremath{u}}\fi}
\newunicodechar{ʸ}{\ifmmode^{y}\else\textsuperscript{\ensuremath{y}}\fi}
\newunicodechar{ᵃ}{\ifmmode^{a}\else\textsuperscript{\ensuremath{a}}\fi}
\newunicodechar{ᵉ}{\ifmmode^{e}\else\textsuperscript{\ensuremath{e}}\fi}
\newunicodechar{ᵏ}{\ifmmode^{k}\else\textsuperscript{\ensuremath{k}}\fi}
\newunicodechar{ᵖ}{\ifmmode^{p}\else\textsuperscript{\ensuremath{p}}\fi}
\newunicodechar{ᴺ}{\ifmmode^{N}\else\textsuperscript{\ensuremath{N}}\fi}
\newunicodechar{ᶜ}{\ifmmode^{c}\else\textsuperscript{\ensuremath{c}}\fi}
\newunicodechar{ʰ}{\ifmmode^{h}\else\textsuperscript{\ensuremath{h}}\fi}
\newunicodechar{ᵠ}{\ifmmode^{q}\else\textsuperscript{\ensuremath{q}}\fi}
\newunicodechar{ₙ}{\ifmmode_{n}\else\textsubscript{\ensuremath{n}}\fi}
\newunicodechar{ₐ}{\ifmmode_{a}\else\textsubscript{\ensuremath{a}}\fi}
\newunicodechar{ᵢ}{\ifmmode_{i}\else\textsubscript{\ensuremath{i}}\fi}
\newunicodechar{ᵣ}{\ifmmode_{r}\else\textsubscript{\ensuremath{r}}\fi}
\newunicodechar{ₖ}{\ifmmode_{k}\else\textsubscript{\ensuremath{k}}\fi}
\newunicodechar{ᵦ}{\ifmmode_{\beta}\else\textsubscript{\ensuremath{\beta}}\fi}
\newunicodechar{ₓ}{\ifmmode_{x}\else\textsubscript{\ensuremath{x}}\fi}
\newfontfamily\popdisplay{ArchivoBlack-Regular.ttf}[Path=../../../fonts/]
\newfontfamily\poplogo{SpaceGrotesk-Bold.ttf}[Path=../../../fonts/]
\newfontfamily\popsemi{PlusJakartaSans-SemiBold.ttf}[Path=../../../fonts/]
\newfontfamily\popxbold{PlusJakartaSans-ExtraBold.ttf}[Path=../../../fonts/]
\newfontfamily\popxboldls{PlusJakartaSans-ExtraBold.ttf}[Path=../../../fonts/,LetterSpace=3.0]

\setlist[enumerate]{leftmargin=1.7em,itemsep=5pt,topsep=4pt}
\setlist[itemize]{leftmargin=1.7em,itemsep=4pt,topsep=4pt,label={\color{poporange}\textbullet}}
\setlength{\parindent}{0pt}
\setlength{\parskip}{5pt}
\renewcommand{\arraystretch}{1.25}

% pgfplots : style Pop par défaut
\pgfplotsset{
  every axis/.append style={axis line style={popink,line width=1pt},tick style={popink},
    grid style={popline,line width=0.5pt},label style={font=\small},tick label style={font=\footnotesize}},
  popcurve/.style={poporange,line width=1.8pt,no markers,smooth},
}
% Même style disponible dans un \draw TikZ simple (hors pgfplots)
\tikzset{popcurve/.style={poporange,line width=1.8pt}}
\tikzset{popgrid/.style={popline,very thin}}
\colorlet{popcurve}{poporange} % le nom du style est parfois utilisé comme couleur

% ============================================================
% Pill / squiggle
% ============================================================
\newcommand{\poppill}[3]{%
  \tikz[baseline=-0.55ex]\node[draw=popink,line width=1.2pt,rounded corners=2.6pt,%
    fill=#2,inner xsep=6.5pt,inner ysep=3pt]{%
    \popxboldls\fontsize{7}{7}\selectfont\color{#3}\MakeUppercase{#1}};%
}
\newcommand{\popsquiggle}[1]{%
  \tikz[baseline=-0.4ex]{
    \draw[#1,line width=2.6pt,line cap=round]
      plot [smooth] coordinates {(0,0.09) (0.34,0.01) (0.68,0.21) (1.02,0.01) (1.36,0.10)};
  }%
}
% Mot-clé surligné (marqueur orange sous le texte)
\newcommand{\cle}[1]{{\popxbold\color{popdark}#1}}

% ============================================================
% En-tête / pied
% ============================================================
\pagestyle{fancy}
\fancyhf{}
\renewcommand{\headrulewidth}{0pt}
\renewcommand{\footrulewidth}{0pt}
\lhead{\raisebox{-0.15ex}{\poplogo\fontsize{13}{13}\selectfont\color{popink}OnBuch\color{poporange}+}}
\rhead{\poppill{\DOCMATIERE\ \textbullet\ \DOCNIVEAU\ \textbullet\ Leçon \DOCLECON}{white}{popink}}
\lfoot{\popxbold\fontsize{7.6}{7.6}\selectfont\color{popmuted}\MakeUppercase{Cours signé OnBuch+}\ \textbullet\ {\popsemi\DOCTITRE}}
\rfoot{\tikz[baseline=-0.6ex]\node[rounded corners=8.5pt,fill=popink,minimum height=17pt,minimum width=17pt,inner xsep=5pt,inner sep=0]{\popxbold\fontsize{8}{8}\selectfont\color{popcream}\thepage\,/\,\pageref*{LastPage}};}

% ============================================================
% Titres
% ============================================================
\newcommand{\chaptitre}[2]{%
  \begin{tcolorbox}[enhanced,colback=poporange,colframe=popink,boxrule=2.6pt,arc=11pt,
      drop shadow=popink,left=15pt,right=15pt,top=12pt,bottom=11pt,before skip=3pt,after skip=18pt]
    \poppill{#2}{popink}{popcream}\par\vspace{9pt}
    {\raggedright\hyphenpenalty=10000\exhyphenpenalty=10000\popdisplay\fontsize{22}{24}\selectfont\color{popcream}\MakeUppercase{#1}\par}
  \end{tcolorbox}
}
% Section numérotée : pastille orange avec le numéro + titre Archivo
\newcounter{popsec}
\newcommand{\coursec}[1]{%
  \stepcounter{popsec}\setcounter{popsub}{0}%
  \par\vspace{16pt}\noindent
  \begin{minipage}{\linewidth}
  \tikz[baseline=-0.7ex]\node[draw=popink,line width=1.6pt,fill=poporange,rounded corners=4pt,minimum size=22pt,inner sep=0]{\popdisplay\fontsize{12}{12}\selectfont\color{popcream}\thepopsec};%
  \hspace{8pt}{\popdisplay\fontsize{15}{17}\selectfont\color{popink}\MakeUppercase{#1}}\par
  \vspace{4pt}\noindent{\color{poporange}\rule{36pt}{3.6pt}}
  \end{minipage}\par\nopagebreak\vspace{8pt}
}
\newcounter{popsub}
\newcommand{\courssub}[1]{%
  \stepcounter{popsub}%
  \par\vspace{9pt}\noindent{\popxbold\fontsize{11.5}{13}\selectfont\color{popink}{\color{poporange}\thepopsec.\thepopsub}\hspace{6pt}#1}\par\nopagebreak\vspace{4pt}
}

% ============================================================
% Boîtes pédagogiques — même recette : fond blanc, bordure épaisse colorée,
% ombre dure encre, badge plein. Titre optionnel [#1].
% ============================================================
\tcbset{popbase/.style={breakable,enhanced,colback=white,boxrule=2.3pt,arc=9pt,
  shadow={3pt}{-3pt}{0pt}{popink},left=14pt,right=14pt,top=11pt,bottom=11pt,before skip=11pt,after skip=15pt,
  fontupper=\popsemi\fontsize{10.6}{14.5}\selectfont\color{popink},
  fontlower=\popsemi\fontsize{10.6}{14.5}\selectfont\color{popink}}}
% Coche dessinée (le glyphe ✓ n'existe pas dans les polices du document)
\newcommand{\popcheck}[1]{\tikz[baseline=-0.5ex]\draw[#1,line width=1.6pt,line cap=round,line join=round](-0.13,0.0)--(-0.04,-0.09)--(0.13,0.1);}
\providecommand{\coche}{\popcheck{popgreen}}
\providecommand{\resultat}[1]{\par\smallskip\noindent\fcolorbox{popgreen}{popgreenL}{\parbox{\dimexpr\linewidth-2\fboxsep-2\fboxrule\relax}{\textbf{Résultat.} #1}}\par\smallskip}
\newcommand{\popboxhead}[3]{\poppill{#1}{#2}{white}\ifx\relax#3\relax\else\hspace{6pt}{\popxbold\fontsize{10.6}{12}\selectfont\color{popink}#3}\fi\par\vspace{7pt}}

\newtcolorbox{aretenir}[1][]{popbase,colframe=popgreen,before upper={\popboxhead{À retenir}{popgreen}{#1}}}
\newtcolorbox{exemplebox}[1][]{popbase,colframe=poporange,before upper={\popboxhead{Exemple}{poporange}{#1}}}
\newtcolorbox{definition}[1][]{popbase,colframe=popblue,before upper={\popboxhead{Définition}{popblue}{#1}}}
\newtcolorbox{propriete}[1][]{popbase,colframe=poppurple,before upper={\popboxhead{Propriété}{poppurple}{#1}}}
\newtcolorbox{methode}[1][]{popbase,colframe=popgold,before upper={\popboxhead{Méthode}{popgold}{#1}}}
\newtcolorbox{attention}[1][]{popbase,colframe=poppink,before upper={\popboxhead{Attention}{poppink}{#1}}}
\newtcolorbox{experience}[1][]{popbase,colframe=popblue,colback=popblueL,before upper={\popboxhead{Expérience}{popblue}{#1}}}
% « Le savais-tu ? » : carte violette pleine (contexte, culture, Cameroun)
\newtcolorbox{savaistu}[1][]{popbase,colback=poppurple,colframe=popink,
  fontupper=\popsemi\fontsize{10.4}{14.2}\selectfont\color{white},
  before upper={\poppill{Le savais-tu\,?}{white}{poppurple}\ifx\relax#1\relax\else\hspace{6pt}{\popxbold\color{white}#1}\fi\par\vspace{7pt}}}
% Exercice résolu : énoncé en haut, \tcblower puis la solution sur fond crème
\newcounter{popexo}\newcounter{popexr}
\newtcolorbox{exoresolu}[1][]{popbase,colframe=poporange,colbacklower=popcream,
  segmentation style={popink,line width=1.2pt,dashed},
  before upper={\stepcounter{popexr}\popboxhead{Exercice résolu \thepopexr}{poporange}{#1}},
  before lower={\poppill{Solution}{popink}{popcream}\par\vspace{6pt}}}
% Exercice (non résolu en place) — la correction est dans la partie Corrigés
\newtcolorbox{exercice}[1][]{popbase,colframe=popink,
  before upper={\stepcounter{popexo}\popboxhead{Exercice \thepopexo}{popink}{#1}}}
% Corrigé
\newtcolorbox{corrige}[1][]{popbase,colframe=popgreen,colback=popgreenL,
  before upper={\popboxhead{Corrigé}{popgreen}{#1}}}
% Objectifs / prérequis / bilan
\newtcolorbox{objectifs}{popbase,colframe=popink,colback=poporangeL,
  before upper={\popboxhead{Objectifs de la leçon}{popink}{}}}
\newtcolorbox{prerequis}{popbase,colframe=popink,colback=popblueL,
  before upper={\popboxhead{Prérequis}{popblue}{}}}
\newtcolorbox{bilan}{popbase,colframe=popink,colback=white,boxrule=2.8pt,
  before upper={\popboxhead{Fiche bilan}{poporange}{L'essentiel à connaître pour l'examen}}}
% Figure encadrée avec légende
\newtcolorbox{popfigure}[1][]{enhanced,breakable=false,colback=white,colframe=popline,boxrule=1.4pt,arc=8pt,
  left=10pt,right=10pt,top=10pt,bottom=8pt,before skip=10pt,after skip=12pt,halign=center,
  lower separated=false,halign lower=center,
  fontlower=\popsemi\fontsize{9}{11.5}\selectfont\color{popmuted},
  #1}
\newcommand{\legende}[1]{\par\vspace{6pt}{\popsemi\fontsize{9}{11.5}\selectfont\color{popmuted}#1}}

% ============================================================
% Signature OnBuch+
% ============================================================
\newcommand{\poplogoinline}[1]{{\poplogo\fontsize{#1}{#1}\selectfont\color{popink}OnBuch\color{poporange}+}}
\newcommand{\signatureonbuch}{%
  \begin{tcolorbox}[enhanced,colback=popink,colframe=popink,boxrule=2.2pt,arc=10pt,
      drop shadow=poporange,left=14pt,right=14pt,top=10pt,bottom=10pt,before skip=0pt,after skip=0pt]
    \tikz[baseline=-0.7ex]\node[circle,fill=popgreen,draw=popcream,line width=1.3pt,minimum size=22pt,inner sep=0]{\popcheck{white}};%
    \hspace{9pt}\begin{minipage}[c]{0.62\linewidth}
      {\popxbold\fontsize{12}{13}\selectfont\color{popcream}Signé\ \textbullet\ {\poplogo OnBuch\color{poporange}+}}\par\vspace{2pt}
      {\raggedright\popsemi\fontsize{8}{10.5}\selectfont\color{popline}\MakeUppercase{Équipe pédagogique OnBuch+}\\\MakeUppercase{Conforme au programme officiel MINESEC}\par}
    \end{minipage}\hfill
    {\poplogo\fontsize{22}{22}\selectfont\color{popcream}OnBuch\color{poporange}+}
  \end{tcolorbox}%
}

% ============================================================
% Page de garde
% #1 titre · #2 matière · #3 niveau · #4 module · #5 n° leçon · #6 durée ·
% #7 description / objectifs (texte ou itemize)
% ============================================================
\newcommand{\pagegarde}[7]{%
  \thispagestyle{empty}
  \newgeometry{a4paper,margin=1.7cm,top=1.6cm,bottom=1.4cm}%
  \begin{tikzpicture}[remember picture,overlay]
    \fill[poporange!18] ([xshift=-3.2cm,yshift=-3.6cm]current page.north east) circle (4.6cm);
    \fill[poppurple!14] ([xshift=2.4cm,yshift=7.6cm]current page.south west) circle (3.8cm);
  \end{tikzpicture}%
  {\poplogo\fontsize{30}{30}\selectfont\color{popink}OnBuch\color{poporange}+}\hfill
  \raisebox{0.6ex}{\poppill{Cours signé \textbullet\ Premium}{popink}{popcream}}\par
  \vspace{1pt}\popsquiggle{poppurple}\par
  \vspace{8pt}
  \begin{tcolorbox}[enhanced,colback=poporange,colframe=popink,boxrule=2.8pt,arc=13pt,
      drop shadow=popink,left=18pt,right=18pt,top=12pt,bottom=13pt,after skip=10pt]
    \poppill{#2 \textbullet\ #3}{popink}{popcream}\hspace{5pt}\poppill{#4}{white}{popink}\par\vspace{8pt}
    {\popdisplay\fontsize{14}{14}\selectfont\color{popink}LEÇON #5}\par\vspace{4pt}
    {\raggedright\hyphenpenalty=10000\exhyphenpenalty=10000\popdisplay\fontsize{25}{27}\selectfont\color{popcream}\MakeUppercase{#1}\par}
    \vspace{8pt}
    {\popxbold\fontsize{9.5}{9.5}\selectfont\color{popink}\MakeUppercase{Durée conseillée : #6}}
  \end{tcolorbox}
  \begin{tcolorbox}[enhanced,colback=white,colframe=popink,boxrule=2.2pt,arc=10pt,
      drop shadow=popink,left=16pt,right=16pt,top=10pt,bottom=10pt,after skip=8pt]
    \poppill{Dans cette leçon}{poppurple}{white}\par\vspace{5pt}
    \popsemi\fontsize{9.3}{12.6}\selectfont\color{popink}#7
  \end{tcolorbox}
  \noindent\begin{minipage}[t]{0.487\linewidth}
    \begin{tcolorbox}[enhanced,colback=popgreen,colframe=popink,boxrule=2.2pt,arc=10pt,
        drop shadow=popink,left=11pt,right=11pt,top=10pt,bottom=10pt,height=3.1cm,valign=center]
      \tcbox[colback=white,colframe=popink,boxrule=1.3pt,arc=5pt,boxsep=0pt,nobeforeafter,
        left=3pt,right=3pt,top=3pt,bottom=3pt,valign=center]{\qrcode[height=1.9cm]{https://play.google.com/store/apps/details?id=com.luvvix.onbuch&hl=fr}}%
      \hspace{8pt}\begin{minipage}[c]{0.5\linewidth}
        {\popdisplay\fontsize{11}{12.5}\selectfont\color{white}\MakeUppercase{Télécharge OnBuch}}\par\vspace{4pt}
        {\raggedright\popsemi\fontsize{8.4}{11}\selectfont\color{white}Gratuit \textbullet\ Google Play\\Tuteur IA, annales, concours, résultats\par}
      \end{minipage}
    \end{tcolorbox}
  \end{minipage}\hfill
  \begin{minipage}[t]{0.487\linewidth}
    \begin{tcolorbox}[enhanced,colback=poppurple,colframe=popink,boxrule=2.2pt,arc=10pt,
        drop shadow=popink,left=11pt,right=11pt,top=10pt,bottom=10pt,height=3.1cm,valign=center]
      \tikz[baseline=-0.7ex]\node[circle,fill=white,draw=popink,line width=1.3pt,minimum size=1.6cm,inner sep=0]{\popdisplay\fontsize{24}{24}\selectfont\color{poppurple}+};%
      \hspace{8pt}\begin{minipage}[c]{0.6\linewidth}
        {\popdisplay\fontsize{11}{12.5}\selectfont\color{white}\MakeUppercase{Passe à OnBuch+}}\par\vspace{4pt}
        {\raggedright\popsemi\fontsize{8.4}{11}\selectfont\color{white}Tous les cours signés, Tuteur IA illimité, fiches PDF — onglet OnBuch+ de l'app\par}
      \end{minipage}
    \end{tcolorbox}
  \end{minipage}
  \vspace{8pt}
  \popcontenu
  \vfill
  \signatureonbuch
  \restoregeometry
  \newpage
}

% Rangée « Ce cours contient » (page de garde)
\newcommand{\poptile}[3]{% #1 couleur #2 gros texte #3 légende
  \begin{tcolorbox}[enhanced,colback=#1,colframe=popink,boxrule=2pt,arc=9pt,shadow={2.5pt}{-2.5pt}{0pt}{popink},
      left=6pt,right=6pt,top=9pt,bottom=9pt,halign=center,valign=center,height=1.75cm,width=\linewidth,nobeforeafter]
    {\popdisplay\fontsize{12.5}{13}\selectfont\color{white}\MakeUppercase{#2}}\par\vspace{3pt}
    {\centering\popsemi\fontsize{7.8}{9.5}\selectfont\color{white}#3\par}
  \end{tcolorbox}}
\newcommand{\popcontenu}{%
  \par\noindent{\popxboldls\fontsize{7.5}{7.5}\selectfont\color{popmuted}CE COURS SIGNÉ CONTIENT}\par\vspace{7pt}
  \noindent\begin{minipage}[t]{0.235\linewidth}\poptile{popblue}{Cours}{complet, illustré, pas à pas}\end{minipage}\hfill
  \begin{minipage}[t]{0.235\linewidth}\poptile{poporange}{Méthodes}{et exercices résolus}\end{minipage}\hfill
  \begin{minipage}[t]{0.235\linewidth}\poptile{poppink}{Exercices}{type examen + corrigés}\end{minipage}\hfill
  \begin{minipage}[t]{0.235\linewidth}\poptile{popgreen}{Fiche bilan}{l'essentiel à retenir}\end{minipage}\par}

% Sommaire
\newtcolorbox{sommaire}{enhanced,colback=white,colframe=popink,boxrule=2.2pt,arc=10pt,
  drop shadow=popink,left=18pt,right=18pt,top=13pt,bottom=13pt,before skip=6pt,after skip=20pt,
  before upper={\poppill{Sommaire}{poppurple}{white}\par\vspace{9pt}\popsemi\fontsize{10.6}{14.5}\selectfont\color{popink}}}

% Signature de fin de cours — poussée tout en bas de la dernière page
\newcommand{\findecours}{%
  \par\vfill\nopagebreak
  \begin{center}\popsquiggle{poporange}\end{center}\vspace{4pt}
  \signatureonbuch
}
\AtBeginDocument{\def\llbracket{\mathopen{[\mkern-3.5mu[}}\def\rrbracket{\mathclose{]\mkern-3.5mu]}}} % intervalles d'entiers (glyphe ⟦ absent de la police)
% Glyphes absents de Fira Math : redéfinis à partir de glyphes disponibles
\AtBeginDocument{%
  \def\setminus{\mathbin{\backslash}}\let\smallsetminus\setminus
  \def\vdots{\mathord{\vbox{\baselineskip3.2pt\lineskiplimit0pt\kern4pt\hbox{.}\hbox{.}\hbox{.}}}}%
  \def\ddots{\mathinner{\mkern1mu\raise6.5pt\hbox{.}\mkern2mu\raise3.6pt\hbox{.}\mkern2mu\raise0.7pt\hbox{.}\mkern1mu}}%
  \def\triangle{\mathord{\scalebox{0.9}{$\symup{\Delta}$}}}%
  \def\nabla{\mathord{\raisebox{\depth}{\scalebox{1}[-1]{$\symup{\Delta}$}}}}%
  \def\checkmark{\popcheck{popdark}}%
  \def\star{\mathbin{\ast}}\def\bigstar{\mathord{\ast}}%
  \def\diamond{\mathbin{\scalebox{0.6}{$\Diamond$}}}%
  \def\bigcup{\mathop{\mathchoice{\raisebox{-0.3em}{\scalebox{1.7}{$\cup$}}}{\scalebox{1.25}{$\cup$}}{\cup}{\cup}}\displaylimits}%
  \def\bigcap{\mathop{\mathchoice{\raisebox{-0.3em}{\scalebox{1.7}{$\cap$}}}{\scalebox{1.25}{$\cap$}}{\cap}{\cap}}\displaylimits}%
  \def\lesseqqgtr{\mathrel{\lessgtr}}\def\bigoplus{\mathop{\oplus}\displaylimits}%
}
\providecommand{\ssi}{si et seulement si} % abréviation fréquente
\AtBeginDocument{\providecommand{\wideparen}{\overparen}} % arc de cercle

%% FIGFIX-BEGIN v5 — correctifs globaux des figures (docs/onbuch-generation/tools/figfix)
\usepackage{adjustbox}\usepackage{etoolbox}\tcbuselibrary{raster}
% toute figure TikZ/pgfplots plus large que la ligne est réduite à la largeur disponible
\BeforeBeginEnvironment{tikzpicture}{\begin{adjustbox}{max width=\linewidth}}
\AfterEndEnvironment{tikzpicture}{\end{adjustbox}}
% légendes pgfplots lisibles par-dessus les courbes
\pgfplotsset{every axis/.append style={legend style={fill=white,fill opacity=0.92,text opacity=1,draw=black!15,inner sep=3pt}}}
% étiquettes d'arêtes (syntaxe edge["..."]) avec fond blanc
\tikzset{every edge quotes/.append style={fill=white,inner sep=1.2pt}}
%% FIGFIX-END
\begin{document}

\pagegarde{Réponse immunitaire}{SVTEEHB}{3ème}{SVTEEHB – classe de 3ème}{N06}{4 séances (fiche de progression harmonisée)}{Cette leçon explore les deux lignes de défense de l'organisme : l'immunité innée (non spécifique), rapide et identique face à tout pathogène, et l'immunité adaptative (spécifique), lente mais ciblée et mémorisée. Elle s'appuie sur des faits expérimentaux historiques (Metchnikoff, Pasteur) et l'analyse de documents microscopiques pour construire les mécanismes biologiques.
\vspace{4pt}
\begin{multicols}{2}\small
\begin{itemize}
\item À la fin de cette leçon, je sais décrire les barrières physiques et chimiques (peau, muqueuses) constitutives de la première ligne de défense non spécifique.
\item Je sais expliquer le déroulement de la réaction inflammatoire (signes cliniques, acteurs cellulaires et moléculaires) à partir de l'expérience de Metchnikoff sur la phagocytose.
\item Je sais interpréter des images microscopiques ou des schémas pour identifier les étapes de la phagocytose et le rôle des macrophages et polynucléaires neutrophiles.
\item Je sais distinguer l'immunité non spécifique (innée) de l'immunité spécifique (adaptative) par leur rapidité, leur spécificité et leur mémoire.
\item Je sais identifier les deux grands types de lymphocytes (B et T) et décrire leur rôle respectif dans l'immunité humorale (anticorps) et l'immunité cellulaire.
\item Je sais exploiter une courbe de cinétique de la réponse immunitaire (primaire vs secondaire) pour mettre en évidence la mémoire immunologique.
\end{itemize}\end{multicols}}

\begin{sommaire}
\begin{multicols}{2}
\begin{enumerate}
\item 1. Les barrières de première ligne : peau et muqueuses
\item 2. La réaction inflammatoire et la phagocytose : la deuxième ligne innée
\item 3. L'immunité spécifique (adaptative) : les lymphocytes, acteurs de la spécificité et de la mémoire
\item 4. Coopération, régulation et applications concrètes : Vaccination et Sérothérapie
\item 5. Activités d'intégration : Synthèse et Résolution de problèmes
\item Activité d'intégration
\item Méthodes et exercices résolus
\item Exercices
\item Corrigés des exercices
\item Fiche bilan
\end{enumerate}
\end{multicols}
\end{sommaire}

\begin{objectifs}
\begin{itemize}
\item À la fin de cette leçon, je sais décrire les barrières physiques et chimiques (peau, muqueuses) constitutives de la première ligne de défense non spécifique.
\item Je sais expliquer le déroulement de la réaction inflammatoire (signes cliniques, acteurs cellulaires et moléculaires) à partir de l'expérience de Metchnikoff sur la phagocytose.
\item Je sais interpréter des images microscopiques ou des schémas pour identifier les étapes de la phagocytose et le rôle des macrophages et polynucléaires neutrophiles.
\item Je sais distinguer l'immunité non spécifique (innée) de l'immunité spécifique (adaptative) par leur rapidité, leur spécificité et leur mémoire.
\item Je sais identifier les deux grands types de lymphocytes (B et T) et décrire leur rôle respectif dans l'immunité humorale (anticorps) et l'immunité cellulaire.
\item Je sais exploiter une courbe de cinétique de la réponse immunitaire (primaire vs secondaire) pour mettre en évidence la mémoire immunologique.
\item Je sais appliquer ces connaissances pour justifier l'intérêt de la vaccination et des règles d'hygiène (désinfection, lavage des mains) dans la prévention des infections au Cameroun.
\item Je sais rédiger une conclusion argumentée mobilisant des résultats expérimentaux pour expliquer comment l'organisme assure son intégrité.
\end{itemize}
\end{objectifs}

\begin{prerequis}
\begin{itemize}
\item Structure cellulaire (noyau, cytoplasme, membrane) et notion de cellule eucaryote (4ème).
\item Notion de micro-organismes (bactéries, virus) et de milieu de culture (4ème).
\item Fonctionnement du système sanguin : globules rouges, globules blancs, plasma (début 3ème).
\item Notion de molécule organique : protéines (anticorps), lipides (membranes).
\item Méthode d'analyse de documents scientifiques (observation, hypothèse, conclusion).
\end{itemize}
\end{prerequis}

\newpage

\coursec{Les barrières de première ligne : peau et muqueuses}

Au quartier Mokolo à Yaoundé, Aminatou se blesse au pied avec un clou rouillé en jouant dans la cour. Trois jours plus tard, son pied est rouge, chaud, gonflé et douloureux : c'est l'\cle{inflammation}. Pourtant, son corps n'a jamais « appris » à connaître ce clou ni les microbes qu'il transportait. Chaque jour, sans que nous nous en rendions compte, des milliards de bactéries, de virus et de poussières tentent d'entrer dans notre organisme — dans l'air que nous respirons, l'eau que nous buvons, les aliments que nous mangeons. Et pourtant, nous ne sommes pas malades en permanence !

\textbf{Problématique :} comment l'organisme se défend-il \emph{immédiatement} contre n'importe quel agresseur, sans l'avoir jamais rencontré ? Quelles sont ces premières barrières qui empêchent les microbes d'entrer, et que se passe-t-il lorsqu'elles sont franchies ?

\begin{definition}[L'immunité innée ou non spécifique]
L'\cle{immunité innée} (ou immunité \cle{non spécifique}) est l'ensemble des défenses de l'organisme présentes dès la naissance. Elle est :
\begin{itemize}
\item \textbf{immédiate} : elle agit en quelques minutes à quelques heures ;
\item \textbf{sans mémoire} : elle ne se souvient pas de l'agresseur rencontré ;
\item \textbf{identique quel que soit l'agent} : la réaction est la même face à une bactérie, un virus ou une écharde.
\end{itemize}
La peau et les muqueuses en constituent la \textbf{première ligne de défense}.
\end{definition}

\courssub{La peau : une muraille vivante à triple protection}

\textbf{Intuition.} Observe ta main : tant que la peau est intacte, aucun microbe ne la traverse. C'est une véritable muraille. Mais cette muraille est plus subtile qu'un simple mur de briques : elle combine trois types de protection.

\begin{definition}[Les trois types de barrières]
\begin{itemize}
\item \textbf{Barrière mécanique (physique)} : obstacle matériel qui empêche le passage (couche de cellules mortes kératinisées).
\item \textbf{Barrière chimique} : substances qui tuent ou ralentissent les microbes (acidité, enzymes).
\item \textbf{Barrière biologique} : la \cle{flore commensale} (microbiote), ensemble de microbes inoffensifs vivant sur nous et empêchant les pathogènes de s'installer par \textbf{compétition} (effet barrière).
\end{itemize}
\end{definition}

\textbf{1. La barrière mécanique : l'épiderme kératinisé.} La peau est formée de trois couches superposées : l'\cle{épiderme} (en surface), le \cle{derme} et l'\cle{hypoderme} (en profondeur). La surface de l'épiderme est constituée de cellules mortes, aplaties, remplies de \cle{kératine} (une protéine dure et imperméable) : c'est la \textbf{couche cornée}. Ces cellules se détachent continuellement (desquamation) et sont remplacées par les cellules profondes : les microbes qui s'y déposent sont éliminés avec elles, comme sur un tapis roulant.

\textbf{2. La barrière chimique : le film hydrolipidique.} Les glandes sudoripares produisent la \textbf{sueur} et les glandes sébacées le \textbf{sébum}. Leur mélange forme à la surface de la peau le \cle{film hydrolipidique}, dont le pH est acide, voisin de \num{5,5}. Or la plupart des bactéries pathogènes se développent mal en milieu acide : le film hydrolipidique freine donc leur multiplication.

\textbf{3. La barrière biologique : la flore cutanée commensale.} Notre peau n'est pas stérile : elle héberge des millions de bactéries inoffensives qui occupent la place et consomment les nutriments. Un pathogène qui arrive ne trouve ni place ni nourriture : c'est l'\textbf{effet barrière} de la flore commensale.

\begin{exemplebox}[Quelques ordres de grandeur]
\begin{itemize}
\item pH de la peau : environ \num{5,5} (acide) ; pH du suc gastrique : environ \num{1,5} (très acide).
\item Le corps humain abrite environ $10^{14}$ bactéries commensales, soit davantage que ses propres cellules (environ $10^{13}$) : en nombre de cellules, nous sommes plus bactériens qu'humains !
\end{itemize}
\end{exemplebox}

\begin{popfigure}
\begin{center}
\begin{tikzpicture}[scale=1.0]
% Hypoderme
\fill[popgoldL] (0,0) rectangle (10,1.2);
\draw[popgold!60, thick] (0,1.2) .. controls (2,1.5) and (4,0.9) .. (6,1.2) .. controls (8,1.5) and (9,1.0) .. (10,1.2);
% Derme
\fill[poppinkL] (0,1.2) .. controls (2,1.5) and (4,0.9) .. (6,1.2) .. controls (8,1.5) and (9,1.0) .. (10,1.2)
 -- (10,3.2) .. controls (8,2.9) and (6,3.5) .. (4,3.2) .. controls (2,2.9) and (1,3.4) .. (0,3.2) -- cycle;
\draw[poppink!70, thick] (0,3.2) .. controls (1,3.4) and (2,2.9) .. (4,3.2) .. controls (6,3.5) and (8,2.9) .. (10,3.2);
% Épiderme
\fill[poporangeL] (0,3.2) .. controls (1,3.4) and (2,2.9) .. (4,3.2) .. controls (6,3.5) and (8,2.9) .. (10,3.2)
 -- (10,4.1) -- (0,4.1) -- cycle;
% Couche cornée
\fill[poporange!60] (0,4.1) rectangle (10,4.7);
\foreach \x in {0.4,1.2,2.0,2.8,3.6,4.4,5.2,6.0,6.8,7.6,8.4,9.2} {\draw[popdark!40, thin] (\x,4.15) rectangle (\x+0.6,4.62);}
% Film hydrolipidique
\draw[popblue, very thick] (0,4.85) -- (10,4.85);
% Glande sudoripare
\draw[popblue!70, thick] (2.5,1.3) .. controls (2.2,1.9) and (2.8,2.5) .. (2.5,3.1) .. controls (2.3,3.6) and (2.6,4.0) .. (2.5,4.85);
\draw[popblue!70, thick, decorate, decoration={coil, amplitude=1.5pt, segment length=3pt}] (2.5,0.5) -- (2.5,1.3);
% Glande sébacée + poil
\draw[popdark, thick] (7,0.9) -- (7.4,5.3);
\fill[popgreen!50] (6.4,2.6) ellipse (0.45 and 0.28);
% Labels
\node[right, font=\small] at (10.2,4.85) {\textbf{Film hydrolipidique} (pH $\approx 5{,}5$)};
\node[right, font=\small] at (10.2,4.4) {\textbf{Couche cornée} (kératine, cellules mortes)};
\node[right, font=\small] at (10.2,3.6) {\textbf{Épiderme} (cellules vivantes)};
\node[right, font=\small] at (10.2,2.2) {\textbf{Derme} (vaisseaux, glandes)};
\node[right, font=\small] at (10.2,0.6) {\textbf{Hypoderme} (graisse)};
\node[left, font=\small, text=popblue!80] at (2.3,0.9) {glande sudoripare};
\node[left, font=\small, text=popgreen!60!black] at (6.3,2.6) {glande sébacée};
\node[font=\small, text=popdark] at (7.9,5.2) {poil};
% Bactéries bloquées
\foreach \x in {1.0,3.8,5.6,8.6} {\fill[poppurple] (\x,5.15) circle (0.09);}
\draw[->, poppurple, thick] (5.6,5.6) -- (5.6,5.3);
\node[font=\small, text=poppurple] at (5.6,5.85) {microbes bloqués à la surface};
\draw[->, popink, very thick] (0.6,5.5) -- (0.6,4.95);
\end{tikzpicture}
\end{center}
\legende{Coupe schématique de la peau. De l'extérieur vers l'intérieur : film hydrolipidique acide (barrière chimique), couche cornée kératinisée (barrière mécanique), épiderme, derme (glandes sudoripares et sébacées) et hypoderme. Les microbes déposés à la surface ne franchissent pas la couche cornée et sont éliminés par desquamation.}
\end{popfigure}

\begin{methode}[Observer une coupe histologique de peau au microscope]
\begin{enumerate}
\item Repérer d'abord la couche la plus superficielle, foncée et sans noyaux : c'est la \textbf{couche cornée} de l'épiderme (cellules mortes kératinisées).
\item Juste en dessous, identifier l'\textbf{épiderme vivant} : cellules serrées, pourvues de noyaux.
\item Dans le \textbf{derme}, plus clair et plus lâche, chercher les \textbf{glandes sudoripares} (petits tubes enroulés en pelote) et les \textbf{glandes sébacées} (petits sacs accolés aux racines des poils).
\item Conclure : la peau est à la fois une barrière physique (couche cornée) et une usine chimique (sueur, sébum).
\end{enumerate}
\end{methode}

\courssub{Les muqueuses : des portes surveillées}

\textbf{Intuition.} L'organisme doit bien communiquer avec l'extérieur : respirer, manger, éliminer. À ces endroits, la peau laisse place aux \cle{muqueuses}, des parois fines et humides qui tapissent les voies respiratoires, digestives et urino-génitales. Plus fragiles que la peau, elles possèdent leurs propres systèmes de défense.

\textbf{1. La muqueuse respiratoire : le tapis roulant du mucus.} Les voies respiratoires (nez, trachée, bronches) sont tapissées d'un \textbf{épithélium cilié} : des cellules portant des \cle{cils vibratiles}, petits battants mobiles. Entre elles, les \textbf{cellules caliciformes} sécrètent le \cle{mucus}, un liquide gluant qui piège poussières et microbes comme du papier tue-mouches. Les cils battent en permanence vers le haut et chassent ce mucus chargé d'impuretés vers la gorge, où il est avalé puis détruit dans l'estomac. C'est l'\textbf{escalator mucociliaire}.

\begin{popfigure}
\begin{center}
\begin{tikzpicture}[scale=1.05]
% Couche de cellules
\foreach \x in {0,1.5,3,4.5,6,7.5,9} {
  \draw[fill=popgreenL, draw=popgreen!60!black] (\x,0) rectangle (\x+1.4,1.6);
  \fill[popgreen!60!black] (\x+0.7,0.5) circle (0.18);
}
% Cellule caliciforme
\draw[fill=popgoldL, draw=popgold!80!black] (4.5,0) .. controls (4.3,0.8) and (4.4,1.3) .. (4.7,1.6) -- (5.7,1.6) .. controls (6.0,1.3) and (6.1,0.8) .. (5.9,0) -- cycle;
\fill[popgold!80!black] (5.2,0.35) circle (0.18);
% Cils
\foreach \x in {0.2,0.5,0.8,1.1,1.4,1.7,2.0,2.3,2.6,2.9,3.2,3.5,3.8,4.1,4.4,4.7,5.0,5.3,5.6,5.9,6.2,6.5,6.8,7.1,7.4,7.7,8.0,8.3,8.6,8.9,9.2} {
  \draw[popblue, thick] (\x,1.6) .. controls (\x+0.15,2.1) and (\x-0.15,2.3) .. (\x+0.25,2.6);
}
% Mucus
\fill[popblue!25, opacity=0.85] (0,2.35) .. controls (2,2.7) and (4,2.2) .. (6,2.55) .. controls (8,2.8) and (9,2.3) .. (10.4,2.5) -- (10.4,3.2) .. controls (8,3.0) and (5,3.4) .. (2,3.1) .. controls (1,3.0) and (0.5,3.25) .. (0,3.15) -- cycle;
% Bactéries piégées
\foreach \p in {(1.2,2.75),(3.4,2.9),(5.6,2.7),(7.8,2.95),(9.2,2.7)} {\fill[poppurple] \p circle (0.1);}
% Flèche sens du transport
\draw[->, popink, very thick] (1.5,3.7) -- (8.5,3.7);
\node[font=\small, text=popink] at (5,4.0) {sens du transport du mucus (vers la gorge)};
% Labels
\node[right, font=\small] at (10.6,2.7) {\textbf{mucus} (piège les microbes)};
\node[right, font=\small] at (10.6,2.0) {\textbf{cils vibratiles}};
\node[right, font=\small] at (10.6,0.8) {\textbf{cellules de l'épithélium}};
\node[font=\small, text=popgold!70!black] at (5.2,-0.5) {cellule caliciforme (sécrète le mucus)};
\draw[->, poppurple] (9.2,2.7) -- (9.9,2.35);
\node[font=\small, text=poppurple] at (8.3,3.35) {bactéries piégées};
\end{tikzpicture}
\end{center}
\legende{Schéma de la muqueuse respiratoire. Les cellules caliciformes sécrètent le mucus qui piège poussières et bactéries ; les cils vibratiles battent et poussent ce mucus contaminé vers la gorge (escalator mucociliaire), où il sera avalé et détruit par l'acidité gastrique.}
\end{popfigure}

\textbf{2. L'acidité gastrique.} L'estomac sécrète un suc très acide, de pH compris entre 1 et 2. Cette acidité tue la quasi-totalité des microbes avalés avec les aliments ou avec le mucus respiratoire. C'est pourquoi, lors d'épidémies de choléra, c'est l'eau contaminée bue en grande quantité (qui dilue l'acidité) qui permet au bacille de passer.

\textbf{3. Le lysozyme : une enzyme qui dissout les bactéries.} Les larmes, la salive et le mucus contiennent le \cle{lysozyme}, une enzyme capable de détruire la paroi de nombreuses bactéries : il les « dissout » littéralement. C'est pourquoi les yeux, pourtant exposés en permanence, sont rarement infectés.

\begin{experience}[Mise en évidence du pouvoir antibactérien des larmes (lysozyme)]
\textbf{Matériel :} boîtes de Pétri contenant une gélose ensemencée avec la bactérie inoffensive \emph{Micrococcus luteus}, larmes recueillies (ou solution de lysozyme), sérum physiologique (témoin), étuve à $37\ ^\circ\mathrm{C}$.\par
\textbf{Protocole :}
\begin{enumerate}
\item On réalise deux petits puits dans la gélose de chaque boîte.
\item Puits A : on dépose quelques gouttes de larmes. Puits B : on dépose du sérum physiologique (témoin).
\item On incube 24 à 48 h à $37\ ^\circ\mathrm{C}$ et on observe.
\end{enumerate}
\textbf{Observations :} la gélose ensemencée devient trouble partout (culture de la bactérie), \textbf{sauf autour du puits A} où apparaît une zone claire, sans bactéries : la \textbf{zone d'inhibition}. Autour du puits B, la culture est normale.\par
\textbf{Interprétation :} une substance présente dans les larmes a diffusé dans la gélose et a détruit (lysé) les bactéries autour du puits A. Cette substance est le \textbf{lysozyme}. Le témoin B prouve que l'effet n'est pas dû au liquide lui-même.\par
\textbf{Conclusion :} les larmes possèdent un pouvoir antibactérien : elles constituent une barrière chimique de l'immunité innée.
\end{experience}

\begin{popfigure}
\begin{center}
\begin{tikzpicture}[scale=1.0]
% Boîte de Pétri
\draw[popdark, thick, fill=popcream] (0,0) circle (3);
\draw[popdark!50, thick] (0,0) circle (3.15);
% Culture trouble (points)
\foreach \a in {0,20,40,60,80,100,120,140,160,180,200,220,240,260,280,300,320,340} {\fill[popgreen!40] ({2.6*cos(\a)},{2.6*sin(\a)}) circle (0.05);}
\foreach \a in {10,45,80,115,150,185,220,255,290,325} {\fill[popgreen!40] ({1.7*cos(\a)},{1.7*sin(\a)}) circle (0.05);}
% Zone d'inhibition autour du puits A
\fill[popcream] (-1.2,0.6) circle (0.95);
\draw[poporange, dashed, thick] (-1.2,0.6) circle (0.95);
\fill[popblue!60] (-1.2,0.6) circle (0.28);
% Puits B témoin
\fill[popmuted!50] (1.4,-0.9) circle (0.28);
% Labels
\draw[->, popdark] (-1.2,0.6) -- (-3.6,2.2); \node[font=\small, left] at (-3.6,2.2) {puits A : larmes (lysozyme)};
\draw[->, poporange] (-1.9,1.3) -- (-3.8,0.9); \node[font=\small, left, text=poporange] at (-3.8,0.9) {zone d'inhibition (bactéries lysées)};
\draw[->, popdark] (1.4,-0.9) -- (3.9,-1.8); \node[font=\small, right] at (3.9,-1.8) {puits B : sérum physiologique (témoin)};
\draw[->, popgreen!50!black] (2.3,1.3) -- (4.0,1.9); \node[font=\small, right, text=popgreen!50!black] at (4.0,1.9) {culture de \emph{M. luteus}};
\end{tikzpicture}
\end{center}
\legende{Boîte de Pétri ensemencée avec \emph{Micrococcus luteus} après 48 h à $37\ ^\circ\mathrm{C}$. Autour du puits contenant les larmes (A), une zone claire d'inhibition témoigne de la lyse des bactéries par le lysozyme ; autour du puits témoin (B), la culture est intacte.}
\end{popfigure}

\begin{savaistu}[La découverte du lysozyme par Fleming (1922)]
En 1922, le médecin écossais Alexander Fleming, enrhumé, laissa tomber une goutte de son mucus nasal sur une culture de bactéries. Quelques jours plus tard, les bactéries avaient disparu à cet endroit ! Il venait de découvrir le \textbf{lysozyme}, présent aussi dans ses larmes, et montra qu'il détruit la paroi bactérienne. Cette découverte préparait celle, en 1928, de la pénicilline, le premier antibiotique — qui lui valut le prix Nobel.
\end{savaistu}

\courssub{Quand la barrière est franchie : les portes d'entrée}

La première ligne de défense est très efficace... tant qu'elle est \textbf{intacte}. Toute effraction de la peau ou des muqueuses ouvre une \cle{porte d'entrée} aux pathogènes :
\begin{itemize}
\item une \textbf{plaie} (comme celle d'Aminatou) permet aux bactéries du sol de pénétrer : le bacille du \textbf{tétanos}, présent dans la terre et sur les objets souillés, en est l'exemple redoutable ;
\item une \textbf{brûlure} étendue détruit la couche cornée sur une large surface : les infections deviennent alors le principal danger ;
\item une \textbf{intubation} ou un cathéter traverse la muqueuse respiratoire ou urinaire et introduit des microbes directement à l'intérieur ;
\item si les microbes atteignent le sang et s'y multiplient, on parle de \textbf{septicémie}, une infection généralisée grave.
\end{itemize}
C'est précisément lorsque cette première ligne est franchie que se déclenche la \textbf{deuxième ligne} de l'immunité innée : la réaction inflammatoire et la phagocytose, que nous étudierons dans la section suivante.

\begin{attention}[Erreur fréquente : confondre flore commensale et microbes pathogènes]
La flore commensale n'est \textbf{pas} un ensemble d'ennemis ! Ce sont des microbes \emph{utiles} qui vivent normalement sur la peau et sur les muqueuses (notamment intestinales) et nous \emph{protègent} en empêchant les pathogènes de s'installer. Notre corps n'est donc pas stérile, et c'est une chance : vouloir « tout désinfecter » détruit cette flore protectrice et fragilise l'organisme. Ne jamais écrire « les microbes de la peau sont dangereux » sans préciser lesquels.
\end{attention}

\begin{exoresolu}[Analyse d'une situation concrète]
\textbf{Énoncé.} Après sa blessure au pied, Aminatou se rend au centre de santé. L'infirmier nettoie la plaie, la désinfecte et vérifie sa vaccination contre le tétanos.\par
1) Quelle barrière de l'immunité innée a été franchie par le clou ? Justifier.\par
2) Pourquoi le clou, même propre en apparence, est-il dangereux ?\par
3) Expliquer pourquoi le nettoyage et la désinfection de la plaie sont urgents.
\tcblower
\textbf{Solution détaillée.}\par
1) Le clou a franchi la \textbf{première ligne de défense} : la peau. Il a déchiré la couche cornée kératinisée (barrière mécanique) et le film hydrolipidique acide (barrière chimique), créant une porte d'entrée directe vers les tissus profonds.\par
2) Même propre en apparence, un objet comme un clou porte à sa surface des micro-organismes de l'environnement (terre, poussière), dont peut-être le bacille du tétanos. La rouille indique un objet resté au sol, donc contaminé. La blessure introduit ces microbes directement dans l'organisme, à l'abri des barrières naturelles.\par
3) Le nettoyage élimine mécaniquement les microbes et les saletés introduits dans la plaie ; le désinfectant détruit chimiquement ceux qui restent. Agir vite empêche les microbes de se multiplier et de gagner le sang (risque de septicémie) ou de provoquer une infection grave comme le tétanos. Ces gestes \emph{remplacent} temporairement la barrière naturelle détruite.
\end{exoresolu}

\begin{aretenir}[L'essentiel de la section]
\begin{itemize}
\item L'\textbf{immunité innée} est immédiate, sans mémoire et identique face à tout agresseur.
\item Sa \textbf{première ligne} est constituée de la \textbf{peau} (couche cornée kératinisée, film hydrolipidique de pH $\approx 5{,}5$, flore commensale) et des \textbf{muqueuses} (mucus et cils vibratiles, acidité gastrique de pH 1 à 2, lysozyme des larmes et de la salive).
\item On distingue trois types de barrières : \textbf{mécanique}, \textbf{chimique} et \textbf{biologique} (flore commensale, effet barrière par compétition).
\item Toute effraction (plaie, brûlure, intubation) est une porte d'entrée : risque d'infection locale (tétanos) ou générale (septicémie).
\end{itemize}
\end{aretenir}

\coursec{La réaction inflammatoire et la phagocytose : la deuxième ligne innée}

Dans la section précédente, nous avons vu que la \cle{peau} et les \cle{muqueuses} constituent de redoutables barrières de première ligne. Mais que se passe-t-il lorsqu'un microbe parvient à franchir ce rempart ? Imagine une épine qui pénètre dans ton doigt en cueillant des mangues : quelques heures plus tard, la zone devient rouge, chaude, gonflée et douloureuse. Ces signes familiers ne sont pas le fait du microbe lui-même : ce sont les manifestations visibles d'une riposte organisée de ton organisme, la \cle{réaction inflammatoire}, au cœur de laquelle agissent des cellules « mangeuses », les \cle{phagocytes}. C'est cette deuxième ligne de défense, rapide et non spécifique, que nous allons étudier.

\courssub{Les signes de l'inflammation : une alarme visible}

\begin{definition}[Réaction inflammatoire]
La \cle{réaction inflammatoire} est l'ensemble des réactions locales déclenchées par les tissus en réponse à une agression (infection, blessure, brûlure). Elle vise à attirer les cellules de défense vers le foyer de l'agression, à détruire les agents agresseurs et à préparer la réparation du tissu.
\end{definition}

Dès le \textsc{i}\ier{} siècle après J.-C., le médecin romain \textbf{Celse} décrivait les quatre signes cardinaux de l'inflammation, encore enseignés aujourd'hui :

\begin{itemize}
\item \textbf{Rubor} (rougeur) : la zone lésée devient rouge ;
\item \textbf{Calor} (chaleur) : la zone est plus chaude que le reste du corps ;
\item \textbf{Tumor} (gonflement) : la zone enfle, formant un œdème ;
\item \textbf{Dolor} (douleur) : la zone est douloureuse, spontanément ou au toucher.
\end{itemize}

On ajoute souvent un cinquième signe, décrit plus tard : la \textbf{functio laesa} (perte de fonction) : un doigt enflammé se plie difficilement, une cheville enflammée supporte mal la marche.

\begin{attention}[Inflammation n'est pas infection]
Ne confonds jamais ces deux termes au BEPC ! L'\textbf{infection} est la présence et la multiplication d'un microbe pathogène dans l'organisme. L'\textbf{inflammation} est la \emph{réponse de défense} de l'organisme à une agression. Il peut y avoir inflammation sans infection (coup, brûlure, entorse) et infection sans inflammation visible (début silencieux de certaines maladies).
\end{attention}

\courssub{Le mécanisme vasculaire : ouvrir les vannes et les portes}

Comment expliquer ces quatre signes ? Par deux phénomènes qui se produisent dans les petits vaisseaux sanguins de la zone agressée.

\textbf{1. La vasodilatation artériolaire.} Les cellules lésées et les cellules de défense déjà présentes libèrent des substances chimiques de signalisation (médiateurs de l'inflammation, comme l'histamine). Sous leur action, les petites artères (\cle{artérioles}) de la zone se dilatent : le débit sanguin local augmente fortement. Cet afflux de sang chaud explique la \textbf{rougeur} et la \textbf{chaleur}.

\textbf{2. L'augmentation de la perméabilité capillaire et veineuse.} Les parois des petits vaisseaux deviennent « poreuses » : leurs cellules s'écartent légèrement. Du liquide riche en protéines (protéines du \cle{complément}, anticorps) s'échappe alors du sang vers les tissus : c'est l'\cle{exsudat}. Son accumulation provoque le \textbf{gonflement} (œdème), qui comprime les terminaisons nerveuses et participe à la \textbf{douleur}.

\textbf{3. La diapédèse.} Les globules blancs (\cle{leucocytes}) qui circulent dans le sang traversent activement la paroi des vaisseaux, en se faufilant entre les cellules : c'est la \cle{diapédèse}. Ils gagnent ainsi le tissu agressé.

\textbf{4. Le chimiotactisme.} Une fois dans le tissu, comment les leucocytes trouvent-ils le foyer de l'agression ? Ils sont guidés par des « pistes odorantes » chimiques : des molécules appelées \cle{cytokines} et des fragments du complément, libérées au niveau du foyer, forment un gradient de concentration que les leucocytes remontent, comme un chien suit une piste. Cette attraction dirigée est le \cle{chimiotactisme}.

\begin{popfigure}
\begin{tikzpicture}[scale=1.0, every node/.style={font=\small}]
% Vaisseau sanguin
\fill[popgoldL] (0,0) rectangle (12,2.2);
\fill[popgold!60] (0,0.25) rectangle (12,1.95);
\draw[popdark, thick] (0,0) -- (12,0);
\draw[popdark, thick] (0,2.2) -- (12,2.2);
\node[popdark] at (1.4,1.1) {\textbf{Sang}};
% Endothélium : cellules écartées à droite
\foreach \x in {0.3,1.1,1.9,2.7,3.5,4.3,5.1,5.9,6.7} {
  \draw[popblue, fill=popblueL] (\x,2.2) ellipse (0.38 and 0.16);
}
\draw[popblue, fill=popblueL] (7.7,2.2) ellipse (0.38 and 0.16);
\draw[popblue, fill=popblueL] (9.6,2.2) ellipse (0.38 and 0.16);
\draw[popblue, fill=popblueL] (11.4,2.2) ellipse (0.38 and 0.16);
\node[popblue, anchor=west] at (0.2,2.75) {\textbf{Endothélium} (cellules écartées)};
\draw[->, popblue] (4.2,2.7) -- (8.6,2.42);
% Leucocyte en diapédèse
\draw[popgreen, very thick, fill=popgreenL] (8.65,2.6) circle (0.42);
\draw[popgreen, thick, fill=popgreen] (8.65,2.6) circle (0.16);
\node[popgreen!60!black, anchor=south] at (8.65,3.15) {Diapédèse d'un leucocyte};
% Exsudat
\fill[popblue!20] (7.2,3.4) rectangle (12,4.6);
\node[popblue!70!black] at (9.6,4.0) {Exsudat (liquide + protéines)};
% Foyer d'infection
\foreach \p in {(2,5.6),(2.6,5.2),(1.6,5.1),(2.3,5.9),(3,5.7)} {
  \draw[popink, fill=popink] \p circle (0.14);
}
\node[popink, anchor=south west] at (1.2,6.0) {\textbf{Foyer : microbes}};
% Chimiotactisme : gradient
\foreach \i/\r in {1/0.9,2/1.7,3/2.5} {
  \draw[popink!40, dashed] (2.2,5.5) circle (\r);
}
\draw[->, poporange, very thick] (8.3,3.6) .. controls (6.5,4.2) .. (4.2,5.0);
\node[poporange!80!black, anchor=west] at (5.4,4.6) {Chimiotactisme};
% Vasodilatation
\draw[<->, poppink, very thick] (0.4,-0.35) -- (3.2,-0.35);
\node[poppink!80!black, anchor=north] at (1.8,-0.45) {Vaisseau dilaté (vasodilatation)};
\end{tikzpicture}
\legende{Schéma de la réaction inflammatoire locale. Le vaisseau se dilate (afflux de sang : rougeur et chaleur), sa paroi devient perméable (exsudat : gonflement), les leucocytes traversent la paroi (diapédèse) puis sont guidés vers le foyer microbien par des substances chimiques (chimiotactisme).}
\end{popfigure}

\courssub{L'expérience fondatrice de Metchnikoff (1882)}

Comment a-t-on découvert que des cellules de notre organisme « mangent » les microbes ? Par une expérience simple et géniale du zoologiste russe \textbf{Élie Metchnikoff}.

\begin{experience}[L'épine de rosier et l'étoile de mer (Metchnikoff, 1882)]
\textbf{Matériel :} larves transparentes d'étoile de mer, épine de rosier, microscope optique.

\textbf{Protocole :} Metchnikoff introduit une fine épine de rosier dans le corps d'une larve d'étoile de mer (animal transparent, donc observable en direct), puis observe au microscope pendant plusieurs heures. Il réalise des expériences voisines en introduisant des grains de colorant chez des daphnies (petits crustacés d'eau douce, eux aussi transparents).

\textbf{Observations :} autour de l'épine, des cellules mobiles de la larve accourent, entourent le corps étranger et l'englobent progressivement dans leur cytoplasme. Les grains de colorant introduits chez la daphnie se retrouvent \emph{à l'intérieur} de ces cellules mobiles.

\textbf{Interprétation :} ces cellules mobiles ingèrent activement les corps étrangers. Metchnikoff les baptise \cle{phagocytes} (du grec \emph{phagein}, manger, et \emph{kutos}, cellule) et propose que cette « alimentation cellulaire » soit en réalité un mécanisme de \textbf{défense} de l'organisme : c'est la naissance de l'\textbf{immunité cellulaire}.
\end{experience}

\begin{methode}[Interpréter l'expérience de Metchnikoff (démarche d'investigation)]
\begin{enumerate}
\item \textbf{Observation :} des cellules mobiles s'accumulent autour du corps étranger introduit et finissent par l'englober.
\item \textbf{Hypothèse :} ces cellules « mangent » les corps étrangers pour les éliminer.
\item \textbf{Vérification :} le colorant introduit se retrouve à l'intérieur des cellules mobiles, ce qui prouve l'ingestion.
\item \textbf{Conclusion :} l'organisme possède des cellules de défense, les phagocytes, capables d'ingérer et de détruire les corps étrangers : c'est la phagocytose.
\end{enumerate}
\end{methode}

\begin{savaistu}[Un prix Nobel... et du lait caillé]
Élie Metchnikoff partagea le prix Nobel de médecine en \textbf{1908} avec Paul Ehrlich, défenseur de l'immunité « humorale » (par les liquides du corps et les anticorps). Convaincu que les microbes de l'intestin accéléraient le vieillissement, Metchnikoff buvait chaque jour du lait fermenté riche en \emph{Lactobacillus} : il est ainsi considéré comme l'un des pères de l'idée des « yaourts santé » !
\end{savaistu}

\courssub{La phagocytose : ingestion et destruction des microbes}

\begin{definition}[Phagocytose]
La \cle{phagocytose} est l'ingestion puis la digestion intracellulaire de particules de taille supérieure à \SI{0,5}{\micro\meter} (microbes, débris cellulaires) par des cellules spécialisées appelées phagocytes : les \textbf{polynucléaires neutrophiles} (PNN) et les \textbf{macrophages}.
\end{definition}

La phagocytose se déroule en \textbf{cinq étapes} :

\begin{enumerate}
\item \textbf{Adhérence :} le phagocyte fixe le microbe à sa surface grâce à des récepteurs qui reconnaissent les corps étrangers (reconnaissance facilitée si le microbe est enrobé d'anticorps ou de protéines du complément : c'est l'\cle{opsonisation}).
\item \textbf{Engouffrement :} la membrane du phagocyte émet des prolongements, les \cle{pseudopodes}, qui entourent le microbe et l'enferment dans une vésicule interne, le \cle{phagosome}.
\item \textbf{Fusion :} des vésicules remplies d'enzymes digestives, les \cle{lysosomes}, fusionnent avec le phagosome pour former le \cle{phagolysosome}.
\item \textbf{Digestion :} les enzymes (lysozyme, protéases) et des substances très oxydantes détruisent le microbe à l'intérieur du phagolysosome.
\item \textbf{Rejet :} les résidus indigestibles sont expulsés hors de la cellule.
\end{enumerate}

\begin{popfigure}
\begin{tikzpicture}[scale=0.85, every node/.style={font=\small}]
% Étape 1 : adhérence
\draw[popgreen, very thick, fill=popgreenL] (0,0) circle (0.9);
\draw[popgreen, fill=popgreen] (0,0) circle (0.3);
\draw[popink, fill=popink] (1.5,0.3) circle (0.18);
\draw[popink, fill=popink] (1.7,-0.15) circle (0.18);
\node[align=center] at (0.6,-1.5) {\textbf{1. Adhérence}};
% Étape 2 : engouffrement
\begin{scope}[xshift=3.8cm]
\draw[popgreen, very thick, fill=popgreenL] (0,0) circle (0.9);
\draw[popgreen, fill=popgreen] (-0.2,-0.2) circle (0.3);
\draw[popgreen, very thick] (0.75,0.5) .. controls (1.4,0.7) and (1.7,0.4) .. (1.75,0);
\draw[popgreen, very thick] (0.9,-0.35) .. controls (1.5,-0.6) and (1.8,-0.3) .. (1.75,0);
\draw[popink, fill=popink] (1.45,0.05) circle (0.18);
\node[align=center] at (0.6,-1.5) {\textbf{2. Engouffrement}\\ (pseudopodes)};
\end{scope}
% Étape 3 : phagosome + lysosomes
\begin{scope}[xshift=7.6cm]
\draw[popgreen, very thick, fill=popgreenL] (0,0) circle (0.9);
\draw[popgreen, fill=popgreen] (-0.35,-0.35) circle (0.28);
\draw[popblue, thick, fill=popblueL] (0.35,0.25) circle (0.32);
\draw[popink, fill=popink] (0.35,0.25) circle (0.15);
\draw[poporange, fill=poporangeL, thick] (-0.45,0.45) circle (0.16);
\draw[poporange, fill=poporangeL, thick] (-0.1,0.6) circle (0.16);
\draw[->, poporange, thick] (-0.3,0.5) -- (0.1,0.35);
\node[align=center] at (0.2,-1.5) {\textbf{3. Fusion}\\ phagosome + lysosomes};
\end{scope}
% Étape 4 : digestion
\begin{scope}[xshift=11.4cm]
\draw[popgreen, very thick, fill=popgreenL] (0,0) circle (0.9);
\draw[popgreen, fill=popgreen] (-0.35,-0.35) circle (0.28);
\draw[popblue, thick, fill=popblueL] (0.3,0.2) circle (0.34);
\foreach \a in {30,90,150,210,270,330} {
  \draw[popink] (0.3,0.2) ++(\a:0.1) -- ++(\a:0.16);
}
\node[align=center] at (0.2,-1.5) {\textbf{4. Digestion}\\ (enzymes)};
\end{scope}
% Étape 5 : rejet
\begin{scope}[xshift=15.2cm]
\draw[popgreen, very thick, fill=popgreenL] (0,0) circle (0.9);
\draw[popgreen, fill=popgreen] (0,-0.2) circle (0.28);
\draw[popink!60, fill=popink!30] (1.25,0.35) circle (0.08);
\draw[popink!60, fill=popink!30] (1.4,0.1) circle (0.06);
\draw[->, popdark, thick] (0.8,0.3) -- (1.15,0.3);
\node[align=center] at (0.3,-1.5) {\textbf{5. Rejet}\\ des résidus};
\end{scope}
% Flèches entre étapes
\foreach \x in {2.0,5.8,9.6,13.4} {
  \draw[->, poppurple, very thick] (\x,0) -- (\x+0.7,0);
}
\end{tikzpicture}
\legende{Les cinq étapes de la phagocytose (d'après les descriptions de Metchnikoff). En vert : le phagocyte et son noyau ; en noir : le microbe ; en orange : les lysosomes ; en bleu : le phagosome puis le phagolysosome.}
\end{popfigure}

\begin{exemplebox}[Des chiffres pour mesurer l'efficacité]
Un \textbf{macrophage} peut phagocyter environ \textbf{100 bactéries} au cours de sa vie ; un \textbf{polynucléaire neutrophile}, environ \textbf{20}. Lors de la digestion, le phagocyte déclenche un véritable « coup de tonnerre » chimique, le \cle{burst oxydatif} : il produit en quelques minutes des substances très agressives pour les microbes, dont le peroxyde d'hydrogène (\ce{H2O2}, l'eau oxygénée) et l'hypochlorite (\ce{HOCl}), le principe actif de... l'eau de Javel ! Ton organisme fabrique donc sa propre eau de Javel microscopique pour désinfecter de l'intérieur.
\end{exemplebox}

\courssub{Les phagocytes : neutrophiles et macrophages, deux profils complémentaires}

\begin{propriete}[Premiers arrivés et résidents]
\begin{itemize}
\item Les \textbf{polynucléaires neutrophiles (PNN)} sont les « premiers arrivés » : ils affluent sur le foyer en quelques heures (pic vers 6 à 12 h après l'infection). Leur durée de vie est courte (quelques jours) : ils meurent en combattant, et leur accumulation avec les microbes et les débris forme le \cle{pus}.
\item Les \textbf{macrophages} sont des cellules « résidentes » des tissus, à longue durée de vie. Ils arrivent ensuite, phagocytent les microbes mais aussi les PNN morts et les débris : ce sont les \textbf{nettoyeurs} qui préparent la guérison. Ils peuvent aussi exposer des fragments du microbe digéré à leur surface, ce qui alerte les lymphocytes et fait le pont avec l'immunité spécifique (section suivante).
\end{itemize}
\end{propriete}

La cinétique d'arrivée des PNN peut se suivre en comptant ces cellules dans le sang et au niveau du foyer après une infection expérimentale :

\begin{popfigure}
\begin{center}
\begin{tikzpicture}
\begin{axis}[
  width=0.85\linewidth, height=6.2cm,
  xlabel={Temps après l'infection (h)},
  ylabel={Nombre de PNN au foyer (unités arbitraires)},
  xmin=0, xmax=48, ymin=0, ymax=110,
  xtick={0,6,12,24,36,48},
  grid=both, grid style={popline!40},
  axis line style={popdark},
]
\addplot[popcurve, domain=0:48, samples=200] {100*exp(-((x-9)^2)/90)};
\addplot[mark=*, poporange, thick] coordinates {(9,100)};
\node[poporange!80!black, anchor=south] at (axis cs:9,100) {Pic vers 6--12 h};
\end{axis}
\end{tikzpicture}
\end{center}
\legende{Courbe cinétique de l'afflux des polynucléaires neutrophiles (PNN) au niveau du foyer infecté en fonction du temps. Le pic est atteint entre 6 et 12 heures après le début de l'infection, puis l'afflux diminue quand les macrophages prennent le relais.}
\end{popfigure}

\begin{attention}[Le pus n'est pas un échec]
Le pus d'un bouton ou d'une plaie est souvent perçu comme un signe de gravité. En réalité, il témoigne d'un combat intense : c'est un mélange de PNN morts, de microbes détruits et de débris de tissus. Sa présence signe le travail de la deuxième ligne de défense. En revanche, une plaie qui s'aggrave, s'étend ou s'accompagne de fièvre doit être montrée à un agent de santé.
\end{attention}

\courssub{Le complément : des protéines du sang auxiliaires de la défense}

Le sang contient un ensemble de protéines, le \cle{système du complément}, qui fonctionne en cascade (chaque protéine activée active la suivante, comme des dominos). Activé au contact des microbes, le complément aide la défense de trois façons, à retenir :

\begin{itemize}
\item il \textbf{opsonise} : il enrobe les microbes et les rend plus « appétissants » pour les phagocytes, ce qui facilite l'adhérence ;
\item il \textbf{attire} : certains de ses fragments servent de signal de chimiotactisme pour guider les leucocytes vers le foyer ;
\item il \textbf{détruit} : certaines protéines du complément s'assemblent pour percer la membrane des microbes et les tuer directement.
\end{itemize}

\courssub{Synthèse et exercice résolu}

\begin{aretenir}[L'essentiel de la section]
\begin{itemize}
\item La réaction inflammatoire est la réponse locale de défense à toute agression ; ses signes sont : rougeur, chaleur, gonflement, douleur (et perte de fonction).
\item Elle résulte de la vasodilatation, de l'augmentation de perméabilité des vaisseaux (exsudat), de la diapédèse et du chimiotactisme des leucocytes.
\item La phagocytose (découverte par Metchnikoff, 1882) est l'ingestion et la digestion des microbes par les phagocytes, en 5 étapes : adhérence, engouffrement, fusion avec les lysosomes, digestion, rejet.
\item Les PNN arrivent vite et meurent en formant le pus ; les macrophages nettoient et alertent les lymphocytes.
\item Le complément opsonise, attire et détruit les microbes.
\end{itemize}
\end{aretenir}

\begin{exoresolu}[Analyse d'une situation concrète]
\textbf{Énoncé.} En jouant au football, Junior se blesse au genou sur un terrain caillouteux. Des microbes pénètrent dans la plaie. Le lendemain, le genou est rouge, chaud, gonflé et douloureux ; deux jours plus tard, un peu de pus apparaît, puis la plaie s'assainit et guérit.
\begin{enumerate}
\item Nomme le phénomène responsable des quatre signes observés le lendemain.
\item Explique l'origine de la rougeur et du gonflement.
\item Que traduit l'apparition du pus ?
\end{enumerate}
\tcblower
\textbf{Solution détaillée.}
\begin{enumerate}
\item Il s'agit de la \textbf{réaction inflammatoire}, réponse de défense non spécifique déclenchée par la lésion et la pénétration des microbes.
\item La \textbf{rougeur} (accompagnée de chaleur) provient de la \textbf{vasodilatation} des petits vaisseaux de la zone : le débit sanguin local augmente. Le \textbf{gonflement} provient de l'augmentation de la \textbf{perméabilité} des vaisseaux : du liquide et des protéines s'échappent du sang et s'accumulent dans le tissu, formant l'\textbf{exsudat} (œdème).
\item Le pus traduit le travail des \textbf{polynucléaires neutrophiles}, premiers phagocytes arrivés sur le foyer : c'est un mélange de PNN morts après avoir phagocyté, de microbes détruits et de débris de tissus. Son apparition puis la guérison montrent que la phagocytose a éliminé les microbes et que les macrophages ont nettoyé la zone.
\end{enumerate}
\end{exoresolu}

\begin{exercice}[Pour t'entraîner]
\begin{enumerate}
\item Définis la phagocytose et cite ses cinq étapes dans l'ordre.
\item Raconte l'expérience de Metchnikoff (protocole, observation, conclusion) en cinq lignes maximum.
\item Pourquoi dit-on que la réaction inflammatoire est « non spécifique » ?
\end{enumerate}
\end{exercice}

\begin{corrige}[Exercice : phagocytose et inflammation]
\begin{enumerate}
\item La phagocytose est l'ingestion puis la digestion intracellulaire de particules (microbes, débris) par des cellules spécialisées, les phagocytes. Étapes : (1) adhérence du microbe à la surface du phagocyte ; (2) engouffrement par des pseudopodes formant le phagosome ; (3) fusion du phagosome avec les lysosomes (phagolysosome) ; (4) digestion du microbe par les enzymes ; (5) rejet des résidus hors de la cellule.
\item Metchnikoff introduit une épine de rosier dans une larve d'étoile de mer transparente et observe au microscope : des cellules mobiles accourent autour de l'épine et l'englobent. Il conclut que l'organisme possède des cellules capables d'ingérer les corps étrangers pour le défendre : les phagocytes.
\item Parce qu'elle se déclenche de la même façon (vasodilatation, exsudat, diapédèse, phagocytose) quelle que soit la nature de l'agresseur : elle ne « reconnaît » pas un microbe particulier, contrairement à l'immunité spécifique des lymphocytes.
\end{enumerate}
\end{corrige}

La deuxième ligne de défense est donc rapide, puissante et efficace contre la plupart des agressions. Mais certains microbes y échappent ou récidivent. L'organisme dispose alors d'une troisième ligne, plus lente mais d'une précision remarquable et dotée de mémoire : l'immunité spécifique des \cle{lymphocytes}, que nous étudierons dans la section suivante.

\coursec{Les barrières de première ligne : peau et muqueuses}

\courssub{La peau : une barrière physique, chimique et microbiologique}

La peau est la première interface entre l'organisme et l'environnement extérieur. Elle assure une protection \textbf{non spécifique} (identique quel que soit l'agresseur) et \textbf{immédiate}.

\begin{popfigure}
\begin{tikzpicture}[
    node distance=0.8cm and 1cm,
    layer/.style={draw, thick, rounded corners, fill=popcream, minimum width=5cm, minimum height=1.2cm, align=center, font=\small},
    arrow/.style={->, thick, draw=popdark, shorten >=2pt, shorten <=2pt},
    labelbox/.style={draw, fill=popblueL, rounded corners, font=\scriptsize, align=center, minimum width=2.5cm}
]
% Layers
\node[layer] (epi) {\textbf{Épiderme} : Cellules kératinisées (cornées), jonctions serrées, desquamation continue};
\node[layer, below=of epi] (der) {\textbf{Derme} : Vaisseaux, nerfs, follicules pileux, glandes sébacées/sudoripares};
\node[layer, below=of der] (hyp) {\textbf{Hypoderme} : Tissu adipeux, isolation, réserve énergétique};
% Annotations
\node[labelbox, left=of epi, xshift=-1cm, align=center] (l1){Barrière \textbf{physique}\\(Kératine, desquamation)};
\node[labelbox, left=of der, xshift=-1cm, align=center] (l2){Barrière \textbf{chimique}\\(Sébum pH $\approx$ 5,5 ; Suueur ; Lysozyme)};
\node[labelbox, left=of hyp, xshift=-1cm, align=center] (l3){Barrière \textbf{microbiologique}\\(Flore commensale \emph{Staphylococcus epidermidis})};
\draw[arrow] (l1) -- (epi.west);
\draw[arrow] (l2) -- (der.west);
\draw[arrow] (l3) -- (hyp.west);
\end{tikzpicture}
\legende{Structure cutanée et ses trois fonctions barrières. La desquamation élimine les microbes fixés. Le film hydrolipidique (sébum + sueur) crée un pH acide inhibiteur. La flore résidente empêche la colonisation par des pathogènes (effet de compétition).}
\end{popfigure}

\begin{propriete}[Les trois verrous cutanés]
\begin{itemize}
    \item \textbf{Physique :} Kératinisation (cornification) des couches superficielles $\to$ étanchéité, résistance mécanique. Desquamation quotidienne $\to$ élimination passive des microbes adhérents.
    \item \textbf{Chimique :} \textbf{Film hydrolipidique} (pH $\approx$ 5,5) acide $\to$ inhibe la plupart des bactéries pathogènes (néutrophiles). \textbf{Lysozyme} (dans sueur, larmes) $\to$ hydrolyse parois bactériennes (peptidoglycane). \textbf{Défensines} (peptides antimicrobiens) sécrétées par kératinocytes.
    \item \textbf{Microbiologique :} \textbf{Flore cutanée résidente} (\emph{Staphylococcus epidermidis}, \emph{Cutibacterium acnes}, champignons \emph{Malassezia}) $\to$ compétition pour nutriments/espace, production de bactériocines, stimulation de l'immunité innée locale.
\end{itemize}
\end{propriete}

\courssub{Les muqueuses : protection des portes d'entrée naturelles}

Les muqueuses tapissent les cavités ouvertes sur l'extérieur (respiratoire, digestive, uro-génitale, oculaire). Elles combinent barrière physique, nettoyage actif et armes chimiques.

\begin{center}
\begin{tabularx}{\linewidth}{|l|X|X|}\hline
\textbf{Muqueuse} & \textbf{Mécanismes physiques / Mécaniques} & \textbf{Armes chimiques / Microbiologiques} \\ \hline
\textbf{Respiratoire} & Épiderme cilié (mucociliaire) : mucus piège particules/microbes $\xrightarrow{\text{battements ciliaires}}$ expectoration / déglutition. Éternuement, toux (réflexes d'expulsion). & Lysozyme, lactoferrine, défensines, IgA sécrétoires. Flore nasopharyngée. \\ \hline
\textbf{Digestive} & Péristaltisme (charriage), desquamation entérocytes, mucus gastrique épais. & \textbf{Acidité gastrique (pH 1,5--3,5)} : barrière majeure. Bile, enzymes pancréatiques. Flore intestinale dense (effet barrière, vit K, stimulation immunité). \\ \hline
\textbf{Uro-génitale} & Flux urinaire (charriage), renouvellement épithélial. & pH vaginal acide (flore de Döderlein : \emph{Lactobacillus} $\to$ lactate/H$_2$O$_2$). IgA sécrétoires. \\ \hline
\textbf{Œil (Conjonctive)} & Clignement (nettoyage mécanique), film lacrymal (renouvellement). & \textbf{Lysozyme} (forte concentration), lactoferrine, IgA. \\ \hline
\end{tabularx}
\end{center}

\begin{attention}[Rupture des barrières = Porte d'entrée]
Toute lésion (plaie, brûlure, cathéter, intubation, rapport sexuel non protégé, piqûre d'insecte \emph{Anopheles} pour paludisme) brise la continuité épithéliale et permet l'invasion microbienne. L'hygiène (lavage mains, asepsie soins) et la vaccination (tétanos) sont les réponses préventives.
\end{attention}

\coursec{La réaction inflammatoire et la phagocytose : la deuxième ligne innée}

\courssub{La réaction inflammatoire : alerte et recrutement}

Lorsqu'un agent franchit les barrières (ex. : éclat de bois contaminé dans le doigt), les tissus lésés et les cellules sentinelles (macrophages résidents, mastocytes) libèrent des \textbf{médiateurs chimiques de l'inflammation} (histamine, prostaglandines, cytokines IL-1, TNF-$\alpha$, chémokines).

\begin{popfigure}
\begin{tikzpicture}[
    node distance=1cm and 1.5cm,
    stepbox/.style={draw, thick, rounded corners, fill=poporangeL, minimum width=3.5cm, minimum height=1.3cm, align=center, font=\small},
    arrow/.style={->, thick, draw=popdark},
    sign/.style={font=\small, text=popdark}
]
\node[stepbox, align=center] (s1){1. Aggression / Infection\\Libération médiateurs\\(Histamine, Cytokines)};
\node[stepbox, right=of s1, align=center] (s2){2. Vasodilatation artérioles\\+ Perméabilité veines\\$\to$ \textbf{Rubor, Calor, Tumor}};
\node[stepbox, right=of s2, align=center] (s3){3. Diapédèse : Sortie leucocytes\\ (Neutrophiles $\to$ Macrophages)\\Guidés par chémokines (\textbf{Chimiotaxie})};
\node[stepbox, right=of s3, align=center] (s4){4. Phagocytose + Réparation\\Résolution ou Abcès / Chronicité};
\draw[arrow] (s1) -- (s2);
\draw[arrow] (s2) -- (s3);
\draw[arrow] (s3) -- (s4);
% Cardinal signs
\node[sign, above=0.2cm of s2.north] {Rubor (Rougeur)};
\node[sign, above=0.8cm of s2.north] {Calor (Chaleur)};
\node[sign, above=1.4cm of s2.north] {Tumor (Œdème/Gonflement)};
\node[sign, below=0.8cm of s2.south] {Dolor (Douleur : compression + médiateurs)};
\node[sign, below=1.4cm of s2.south] {Functio laesa (Perte fonction)};
\end{tikzpicture}
\legende{Les 4 étapes de la réaction inflammatoire aiguë et les 5 signes cardinaux (Celse, 1er siècle ; Galien). L'inflammation isole le foyer, apporte les effecteurs (leucocytes, protéines plasmatiques : complément, anticorps) et initie la réparation.}
\end{popfigure}

\begin{definition}[Médiateurs clés de l'inflammation]
\begin{itemize}
    \item \textbf{Histamine / Sérotinine (Mastocytes, Basophiles) :} Vasodilatation immédiate, augmentation perméabilité vasculaire (œdème).
    \item \textbf{Prostaglandines / Leucotriènes :} Potentialisent douleur, fièvre, chimiotaxie.
    \item \textbf{Cytokines (IL-1, TNF-$\alpha$, IL-6) :} Systémiques (fièvre, synthèse protéines de phase aiguë au foie : CRP) et locales (activation endothélium, chimiotaxie).
    \item \textbf{Chémokines (IL-8 / CXCL8) :} Gradient chimiotactique spécifique pour neutrophiles puis monocytes.
\end{itemize}
\end{definition}

\courssub{La phagocytose : l'ingestion et la digestion intracellulaire}

Les \textbf{phagocytes professionnels} sont les \textbf{polynucléaires neutrophiles} (premiers arrivés, vie courte, pus) et les \textbf{macrophages} (issus des monocytes, longue vie, nettoyage, présentation antigène).

\begin{experience}[Observation de la phagocytose au microscope (TP classique)]
\textbf{Matériel :} Lame de sang périphérique ou culture de macrophages (ligne RAW 264.7 ou péritonéaux souris), bactéries \emph{E. coli} ou levures \emph{Saccharomyces} (vivantes ou latex billes), coloration May-Grünwald Giemsa.
\textbf{Protocole :}
\begin{enumerate}
    \item Mettre en contact phagocytes et proies (ratio 1:10) à 37°C.
    \item Prélèvements à $t=0, 15, 30, 60$ min $\to$ étalement, fixation, coloration.
    \item Observation $\times 1000$ (immersion).
\end{enumerate}
\textbf{Observations attendues :}
\begin{itemize}
    \item $t=0$ : Proies extracellulaires.
    \item $t=15$ min : Formation de \textbf{pseudopodes} $\to$ englobement $\to$ \textbf{phagosome} (vacuole intracellulaire).
    \item $t=30$ min : Fusion phagosome + granules lysosomaux $\to$ \textbf{phagolysosome}. Coloration change (digestion).
    \item $t=60$ min : Résidus indigestes (corps résiduels) $\to$ exocytose.
\end{itemize}
\textbf{Interprétation :} La phagocytose est un processus actif, en 4 phases : \textbf{Chimiotaxie $\to$ Adhérence/Reconnaissance $\to$ Ingestion (Phagosome) $\to$ Digestion (Phagolysosome)}.
\end{experience}

\begin{propriete}[Le phagolysosome : une arme chimique]
La fusion apporte des enzymes hydrolytiques (pH acide $\approx$ 4,5) et lance l'\textbf{explosion respiratoire} (NADPH oxydase $\to$ $O_2^-$ $\to$ $H_2O_2$, HOCl via myélopéroxydase, $^{\bullet}OH$). Ces \textbf{espèces réactives de l'oxygène (ERO)} tuent la plupart des bactéries. \emph{Mycobacterium tuberculosis} résiste en bloquant la fusion phagosome-lysosome.
\end{propriete}

\begin{aretenir}[Immunité innée (Non spécifique) : Résumé]
\begin{itemize}
    \item \textbf{Acteurs :} Barrières (peau/muqueuses), Phagocytes (Neutrophiles, Macrophages), NK, Protéines (Complément, Protéines de phase aiguë, Cytokines).
    \item \textbf{Caractères :} Rapide (min/heures), Identique pour tout agresseur (pas de spécificité fine), \textbf{Pas de mémoire} (réponse identique au 2ème contact), Ne distingue pas Soi/Non-soi finement.
    \item \textbf{Rôle clé :} Contrôler l'infection en attendant (et en activant) l'immunité spécifique.
\end{itemize}
\end{aretenir}

\coursec{L'immunité spécifique (adaptative) : les lymphocytes, acteurs de la spécificité et de la mémoire}

\courssub{Pourquoi une troisième ligne de défense ?}

L'immunité innée est parfois insuffisante : microbes intracellulaires (virus, \emph{Mycobacterium}), bactéries encapsulées (résistantes à la phagocytose), charge infectieuse massive. L'organisme a besoin d'une riposte \textbf{sur mesure} : \textbf{reconnaissance précise} de l'ennemi + \textbf{mémoire} pour réagir plus vite et plus fort ultérieurement. C'est le rôle des \textbf{lymphocytes}.

\begin{definition}[Antigène (Ag) et Anticorps (Ac)]
\begin{itemize}
    \item \textbf{Antigène :} Toute substance \textbf{étrangère} (protéine, polysaccharide de surface microbe, toxine) capable de \textbf{déclencher} une réponse immunitaire spécifique et d'être \textbf{reconnue} par des lymphocytes ou des anticorps.
    \item \textbf{Anticorps (Immunoglobuline) :} Protéine en Y, produite par les \textbf{plasmocytes} (lymphocytes B différenciés), capable de se lier \textbf{spécifiquement} à l'antigène qui a déclenché sa production (clé dans serrure).
\end{itemize}
\end{definition}

\courssub{Les deux familles de lymphocytes : B et T}

Tous les lymphocytes naissent dans la \textbf{moelle osseuse} mais mûrissent en deux lieux différents.

\begin{popfigure}
\begin{tikzpicture}[
    level distance=1.6cm,
    sibling distance=3cm,
    every node/.style={draw, rounded corners, fill=popcream, thick, minimum width=2.5cm, minimum height=0.9cm, font=\small, align=center},
    edge from parent/.style={thick, draw=popdark, ->},
    level 1/.style={sibling distance=4.5cm}
]
\node[align=center]{Cellule Souche\\ Hématopoïétique\\(Moelle osseuse)}
    child { node[align=center]{Lignée Myéloïde\\(Innée : Neutrophiles,\\Macrophages, Plaquettes,\\Globules rouges)} }
    child { node {Lignée Lymphoïde}
        child { node[align=center]{\textbf{Lymphocytes B}\\Maturation : \textbf{Moelle osseuse}\\Récepteur : BCR (IgM membranaire)\\Rôle : Immunité \textbf{Humorale}\\(Anticorps)} }
        child { node[align=center]{\textbf{Lymphocytes T}\\Maturation : \textbf{Thymus}\\Récepteur : TCR\\Rôle : Immunité \textbf{Cellulaire}\\(Destruction cellules infectées,\\Aide aux B)} }
    };
\end{tikzpicture}
\legende{Origine commune et divergence des deux lignées lymphocytaires. B = Bursa équivalent (Moelle osseuse chez l'homme). T = Thymus.}
\end{popfigure}

\begin{propriete}[Différences majeures entre Lymphocytes B et T (Niveau 3ème)]
\begin{center}
\begin{tabular}{|l|c|c|}\hline
\textbf{Caractère} & \textbf{Lymphocyte B} & \textbf{Lymphocyte T} \\ \hline
\textbf{Lieu de maturation} & Moelle osseuse & Thymus \\ \hline
\textbf{Récepteur de surface} & BCR = IgM membranaire & TCR (ne voit Ag que présenté) \\ \hline
\textbf{Reconnaissance de l'Ag} & \textbf{Directe} (Ag natif, conformation 3D) & \textbf{Indirecte} (Peptide Ag sur CMH d'une autre cellule) \\ \hline
\textbf{Effecteur principal} & \textbf{Plasmocyte} $\to$ sécrète \textbf{Anticorps solubles} & \textbf{LT Cytotoxique (CD8)} $\to$ tue cellule cible\\ & & \textbf{LT Auxiliaire (CD4)} $\to$ aide (cytokines) \\ \hline
\textbf{Type d'immunité} & \textbf{Humorale} (Ac dans lymphe/sang) & \textbf{Cellulaire} (contact cellule-cellule) \\ \hline
\textbf{Cibles privilégiées} & Bactéries extracellulaires, Virus libres, Toxines & Cellules infectées (virus, bactéries intracellulaires), Cancéreuses, Greffons \\ \hline
\end{tabular}
\end{center}
\end{propriete}

\courssub{Activation et sélection clonale : le principe de la spécificité}

Chaque lymphocyte possède \textbf{un récepteur unique} (spécifique d'un seul antigène). Le répertoire est immense ($\sim 10^7$ spécificités différentes). L'antigène \textbf{sélectionne} le lymphocyte qui le reconnaît.

\begin{methode}[Théorie de la sélection clonale (Burnet, 1957) -- Étapes]
\begin{enumerate}
    \item \textbf{Rencontre :} Un lymphocyte naïf (jamais vu l'Ag) rencontre \textbf{son} antigène spécifique dans un organe lymphoïde (ganglion, rate, amygdales).
    \item \textbf{Activation (Signal 1 + Signal 2) :} Reconnaissance de l'Ag + Signaux de co-stimulation (cytokines, contact cellule-cellule). \emph{Pour les B : aide souvent nécessaire des LT CD4+ (Coopération T-B).}
    \item \textbf{Expansion clonale :} Divisions mitotiques rapides $\to$ \textbf{millions de clones identiques} (même récepteur) en 3--5 jours.
    \item \textbf{Différenciation :} Deux destins :
    \begin{itemize}
        \item \textbf{Cellules effectrices} (vie courte) : Plasmocytes (usines à Ac) ou LT Cytotoxiques (tueurs). Exécutent la réponse \textbf{immédiate}.
        \item \textbf{Cellules mémoire} (vie longue, années) : Restent en veille. Réagissent \textbf{plus vite et plus fort} au 2ème contact (base de la vaccination).
    \end{itemize}
    \item \textbf{Élimination de l'Ag} $\to$ Apoptose de 90--95\% des effecteurs (contraction). Persistance du pool mémoire.
\end{enumerate}
\end{methode}

\begin{popfigure}
\begin{tikzpicture}[
    node distance=1.2cm and 1.5cm,
    cell/.style={draw, thick, rounded corners, minimum width=2.8cm, minimum height=1.1cm, align=center, font=\small},
    arrow/.style={->, thick, draw=popdark},
    div/.style={draw=popblue, thick, dashed, ->, font=\scriptsize}
]
\node[cell, fill=popblueL, align=center] (Naif){Lymphocyte B ou T \textbf{Naïf}\\(1 cellule, récepteur unique)};
\node[cell, fill=poporangeL, below=of Naif, align=center] (Act){\textbf{ACTIVATION}\\(Rencontre Ag + Signaux)};
\node[cell, fill=popgreenL, below=of Act, xshift=-2.5cm, align=center] (Eff){\textbf{CELLULES EFFECTRICES}\\(Millions de clones, vie courte)\\\textbf{B $\to$ Plasmocytes (Ac)}\\ \textbf{T $\to$ Cytotoxiques / Aide}};
\node[cell, fill=poppinkL, below=of Act, xshift=2.5cm, align=center] (Mem){\textbf{CELLULES MÉMOIRE}\\(Millions de clones, vie longue)\\\textbf{Veille : Réponse secondaire rapide}};
\draw[arrow] (Naif) -- node[left] {Rencontre Ag} (Act);
\draw[div] (Act) -- node[left, align=center]{Mitoses\\Expansion clonale} (Eff);
\draw[div] (Act) -- node[right] {Différenciation} (Mem);
\end{tikzpicture}
\legende{Schéma de la sélection clonale. Un seul lymphocyte spécifique est sélectionné par l'antigène, prolifère massivement (expansion clonale) et se différencie en effecteurs (réponse primaire) et en cellules mémoire (base de la réponse secondaire).}
\end{popfigure}

\courssub{Les deux bras de la réponse spécifique : Humorale et Cellulaire}

\begin{popfigure}
\begin{tikzpicture}[
    node distance=1cm and 1.2cm,
    title/.style={draw, thick, fill=popblue, text=white, rounded corners, minimum width=4.5cm, minimum height=0.8cm, font=\bfseries\small, align=center},
    box/.style={draw, thick, rounded corners, minimum width=3.5cm, minimum height=1.1cm, align=center, font=\small},
    arrow/.style={->, thick, draw=popdark},
]
% Titles
\node[title] (TH) at (0,4) {IMMUNITÉ HUMORALE (Médiée par B)};
\node[title] (TC) at (9,4) {IMMUNITÉ CELLULAIRE (Médiée par T)};
% Humoral
\node[box, fill=popblueL, below=of TH] (Bn) {Lymphocyte B naïf};
\node[box, fill=popblueL, below=of Bn] (ActB) {Activation (Ag + Aide LT CD4+)};
\node[box, fill=popblueL, below=of ActB, align=center] (Plasma){PLASMOCYTE\\Sécrétion massive d'Anticorps};
\node[box, fill=popblueL, below=of Plasma] (MemB) {Lymphocyte B Mémoire};
% Cellular
\node[box, fill=popgreenL, below=of TC] (Tn) {Lymphocyte T naïf (CD4 \& CD8)};
\node[box, fill=popgreenL, below=of Tn] (ActT) {Activation (CPA + MHC + Co-stim)};
\node[box, fill=popgreenL, below=of ActT, xshift=-1.5cm, align=center] (Th){LT Auxiliaire (CD4+)\\Cytokines : Aide B, Active Macrophages, Active CD8};
\node[box, fill=popgreenL, below=of ActT, xshift=1.5cm, align=center] (Tc){LT Cytotoxique (CD8+)\\Perforine/Granzymes $\to$ Lyse cible};
\node[box, fill=popgreenL, below=of Th] (MemT) {Lymphocytes T Mémoire (CD4+ \& CD8+)};
% Targets
\node[box, fill=poporangeL, right=of Plasma, xshift=1cm, align=center] (CibH){CIBLES :\\Bactéries extracellulaires\\Virus libres, Toxines};
\node[box, fill=poporangeL, right=of Tc, xshift=1cm, align=center] (CibC){CIBLES :\\Cellules infectées (virus, bactéries intracellulaires)\\Cellules tumorales, Greffons};
% Mechanisms Antibodies
\node[box, fill=poppinkL, below=of CibH, yshift=-0.5cm, align=center] (MechH){ACTIONS DES ANTICORPS :\\$\bullet$ \textbf{Neutralisation} (blocage entrée cellule)\\$\bullet$ \textbf{Opsonisation} (marquage $\to$ phagocytose facilitée)\\$\bullet$ \textbf{Activation Complément} (lyse, inflammation)\\$\bullet$ \textbf{ADCC} (via cellules NK)};
\draw[arrow] (Bn) -- (ActB) -- (Plasma) -- (MemB);
\draw[arrow] (Tn) -- (ActT);
\draw[arrow] (ActT) -| (Th);
\draw[arrow] (ActT) -| (Tc);
\draw[arrow] (Th) -- (MemT);
\draw[arrow] (Tc) -- (MemT);
\draw[arrow] (Plasma) -- (CibH);
\draw[arrow] (Tc) -- (CibC);
\draw[arrow] (Th) -- node[above, sloped, font=\tiny] {Cytokines (IFN-$\gamma$)} (CibC);
\end{tikzpicture}
\legende{Les deux bras indissociables de l'immunité adaptative. L'immunité humorale (Ac) nettoie l'espace extracellulaire. L'immunité cellulaire (LT CD8) élimine les cellules hôtes compromises. Les LT CD4+ (Auxiliaires) sont les chefs d'orchestre : ils aident les B (humoral) et activent macrophages/CD8 (cellulaire).}
\end{popfigure}

\courssub{Cinétique de la réponse immunitaire : Preuve de la mémoire}

L'expérience classique (injection Ag, dosage Ac sériques) montre deux profils radicalement différents.

\begin{popfigure}
\begin{tikzpicture}
\begin{axis}[
    width=\linewidth, height=8cm,
    xlabel={Temps (jours)}, ylabel={Titre d'anticorps sériques (échelle logarithmique)},
    xmin=0, xmax=60, ymin=0, ymax=7,
    xtick={0,7,14,21,28,35,42,49,56},
    ytick={0,1,2,3,4,5,6,7},
    grid=major, grid style={dashed, gray!30},
    legend style={at={(0.02,0.98)}, anchor=north west, fill=white, fill opacity=0.9, draw=popdark, font=\small},
    axis line style={thick},
    tick label style={font=\small},
    label style={font=\small},
    clip=false
]
% Primary Response
\addplot[popcurve, popblue, very thick, domain=0:35, samples=200] 
    ({x}, { 
        (x < 7) ? 0 : 
        (x < 14) ? 2.5*((x-7)/7) : 
        (x < 21) ? 2.5 + 1.5*((x-14)/7) : 
        (x < 28) ? 4.0 - 0.5*((x-21)/7) :
        3.5*exp(-(x-28)/10)
    });
\addlegendentry{Réponse Primaire (1$^{\text{er}}$ contact)}
% Secondary Response (Booster at day 35)
\addplot[popcurve, poporange, very thick, domain=35:60, samples=200] 
    ({x}, { 
        (x < 38) ? 3.5*exp(-(x-28)/10) : 
        (x < 42) ? 0.5 + 5.5*((x-38)/4) : 
        (x < 49) ? 6.0 - 0.5*((x-42)/7) :
        5.5*exp(-(x-49)/15)
    });
\addlegendentry{Réponse Secondaire (Rappel J35)}
% Annotations
\node[popblue, font=\small, align=left] at (10, 1.5) {Phase de\\latence $\approx$ 7--10 j};
\node[popblue, font=\small, align=left] at (17, 4.5) {Pic : IgM puis IgG};
\node[popblue, font=\small, align=left] at (30, 2.5) {Déclin progressif};
\node[poporange, font=\small, align=left] at (38, 2.5) {Injection Rappel};
\node[poporange, font=\small, align=left] at (42, 6.2) {Pas de latence\\Pic IgG $\times$ 10--100\\Affinité $\uparrow\uparrow$};
\end{axis}
\end{tikzpicture}
\legende{Cinétique comparative Primaire vs Secondaire. \textbf{Primaire :} Latence (activation/expansion), titre modéré, IgM dominants au début. \textbf{Secondaire :} Pas de latence (cellules mémoire prêtes), titre massif, IgG dominants et haute affinité. Preuve fonctionnelle de la \textbf{mémoire immunologique}.}
\end{popfigure}

\begin{methode}[Interpréter une courbe de réponse immunitaire (Sérotherapie / Vaccination)]
\begin{enumerate}
    \item \textbf{Repérer les phases :} Latence (taux basal), Montée exponentielle, Plateau/Pic, Déclin.
    \item \textbf{Comparer Primaire et Secondaire :} Absence de latence, pente plus raide, pic 10 à 100 fois plus haut, IgG précoces et majoritaires (vs IgM en primaire).
    \item \textbf{Conclure :} La réponse secondaire prouve l'existence de \textbf{cellules mémoire} (B et T) générées lors du premier contact.
    \item \textbf{Application vaccinale :} Primo-vaccination (souvent 2 doses) = réponse primaire amplifiée. Rappels (boosters) = réponses secondaires pour maintenir taux d'Ac $>$ seuil protecteur.
\end{enumerate}
\end{methode}

\begin{exemplebox}[Différence Vaccination vs Sérothérapie]
\begin{center}
\begin{tabular}{|l|c|c|}\hline
\textbf{Caractère} & \textbf{VACCINATION (Immunité Active)} & \textbf{SÉROTHÉRAPIE (Immunité Passive)} \\ \hline
\textbf{Substance injectée} & Antigène (microbe atténué, inactivé, fragment, ARNm) & Anticorps tout faits (sérum hyperimmune) \\ \hline
\textbf{Origine des Ac} & \textbf{Produits par l'organisme} (Plasmocytes du vacciné) & \textbf{Produits par un autre organisme} (Cheval, Humain, Souris humanisée) \\ \hline
\textbf{Délai de protection} & Lent (semaines : latence + expansion) & \textbf{Immédiat} (Ac déjà présents) \\ \hline
\textbf{Durée de protection} & \textbf{Longue} (années, vie : cellules mémoire) & \textbf{Courte} (semaines/mois : catabolisme Ac, pas de mémoire) \\ \hline
\textbf{Risque} & Faible (réactions locales, fièvre) & Réactions allergiques (maladie du sérum), transmission agents (rare) \\ \hline
\textbf{Indication principale} & \textbf{Prévention} (Calendrier vaccinal EPI Cameroun) & \textbf{Urgence curative} (Tétanos non vacciné, morsure serpent, rage post-exposition) \\ \hline
\end{tabular}
\end{center}
\end{exemplebox}

\coursec{Coopération, Régulation et Applications concrètes : Vaccination et Sérothérapie}

\courssub{La coopération T-B : condition de l'immunité humorale de qualité}

La plupart des antigènes (protéines) sont \textbf{thymo-dépendants} : le lymphocyte B a \textbf{besoin de l'aide d'un LT CD4+ activé} pour produire des IgG, générer de la mémoire et maturer l'affinité.
\begin{itemize}
    \item \textbf{Étape 1 :} Le B internalise l'Ag via son BCR, le dégrade, présente des peptides sur son \textbf{CMH de classe II}.
    \item \textbf{Étape 2 :} Un LT CD4+ spécifique (déjà activé par une CPA) reconnaît le complexe Peptide-CMH-II sur le B.
    \item \textbf{Étape 3 :} Contact \textbf{CD40L (sur T) -- CD40 (sur B)} + Cytokines (IL-4, IL-21) $\to$ \textbf{Activation complète du B} (Switch de classe IgM$\to$IgG/IgA, Maturation d'affinité, Mémoire, Plasmocytes vie longue).
\end{itemize}
\emph{Sans aide T (Ag thymo-indépendants : polysaccharides capsulaires purs) $\to$ Réponse faible, quasi seulement IgM, pas de mémoire, pas de switch. D'où vaccins conjugués (polysaccharide + protéine porteuse) pour recruter l'aide T (ex. \emph{Haemophilus influenzae} b, Pneumocoque, Méningocoque).}

\courssub{Le calendrier vaccinal au Cameroun (Programme Élargi de Vaccination - PEV)}

\begin{center}
\begin{tabularx}{\linewidth}{|l|X|X|}\hline
\textbf{Maladie ciblée} & \textbf{Type de vaccin (Exemples)} & \textbf{Calendrier PEV Cameroun (Simplifié)} \\ \hline
Tuberculose & BCG (Atténué vivant) & À la naissance (intradermique) \\ \hline
Polio & VPO (Atténué oral) / VPI (Inactivé injectable) & VPO : Naissance, 6, 10, 14 sem. VPI : 14 sem. \\ \hline
Diphtérie, Tétanos, Coqueluche, Hépatite B, Hib & Pentavalent (Inactivé / Sous-unités / Conjugué) & 6, 10, 14 semaines (IM) \\ \hline
Rougeole / Rubéole & RR (Atténué vivant) & 9 mois, 15--18 mois \\ \hline
Fièvre jaune & VVAA (Atténué vivant) & 9 mois \\ \hline
Méningite A (Ceinture africaine) & MenAfriVac (Conjugué) & Campagnes de masse 1--29 ans + routine 9 mois \\ \hline
VPH (Cancer col utérus) & VPH (VLP - Particules virales) & Filles 9--14 ans (2 doses) \\ \hline
COVID-19 & ARNm / Vecteur viral / Inactivé & Selon recommandations OMS/MSP actuelles \\ \hline
\end{tabularx}
\end{center}

\begin{savaistu}[Stratégie "Zéro Palu" et Vaccin RTS,S/AS01 (Mosquirix)]
Le Cameroun fait partie des pays pilotes (avec Ghana, Kenya, Malawi) pour le déploiement du vaccin \textbf{RTS,S/AS01} contre \emph{Plasmodium falciparum} (paludisme). C'est un vaccin \textbf{sous-unité recombinant} (protéine CSP du sporozoïte fusionnée à l'antigène de surface HBsAg de l'hépatite B, adjuvant AS01). Il cible le stade pré-érythrocytaire. Efficacité $\approx$ 30--40\% sur 4 ans pour les formes graves. Il \textbf{ne remplace pas} les MILDA (moustiquaires) ni la chimioprévention (CPS), mais s'y ajoute. Déploiement progressif dans les zones à transmission modérée/élevée depuis 2024.
\end{savaistu}

\courssub{La sérothérapie : immunité passive d'urgence}

Injection d'immunoglobulines polyvalentes ou spécifiques (sérums hyperimmunes).
\begin{itemize}
    \item \textbf{Sérum antitetanique (SAT) / Tétanos :} Urgence plaie souillée chez non vacciné ou vaccination incomplète. Neutralise la toxine \emph{libre} (pas celle déjà fixée sur neurones). Demi-vie $\approx$ 2--3 semaines.
    \item \textbf{Sérum antivenimeux (SAV) :} Morsure serpent venimeux (Vipère, Cobra, Mamba au Cameroun). Polyvalent (Afrique) ou monovalent. Administration IV stricte en milieu hospitalier (surveillance choc anaphylactique).
    \item \textbf{Immunoglobulines anti-Rage :} Associées au vaccin rage post-exposition (catégorie III : morsure, griffure, léchage muqueuse). Infiltration autour plaie + IM site distant. Neutralisent virus localement avant que vaccin n'agisse.
\end{itemize}

\begin{attention}[Pièges BEPC : Vaccination vs Sérothérapie]
\begin{itemize}
    \item Ne pas confondre \textbf{Antigène} (injecté dans vaccin) et \textbf{Anticorps} (injectés en sérothérapie).
    \item La vaccination \textbf{ne soigne pas} une infection en cours (trop lent), elle \textbf{prévient}. La sérothérapie \textbf{soigne/previent immédiatement} mais ne protège pas durablement.
    \item "Rappel vaccinal" = stimulation mémoire (réponse secondaire). "Réinjection sérum" = nouvelle dose passive (pas de mémoire).
\end{itemize}
\end{attention}

\coursec{Activités d'intégration : Synthèse et Résolution de problèmes}

\courssub{Tableau de synthèse : Immunité Innée vs Adaptative}

\begin{center}
\begin{tabularx}{\linewidth}{|l|X|X|}\hline
\textbf{Critère} & \textbf{Immunité Innée (Non Spécifique)} & \textbf{Immunité Adaptative (Spécifique)} \\ \hline
\textbf{Lignes de défense} & 1ère (Barrières) \& 2ème (Inflammation, Phagocytose, NK, Complément) & 3ème (Lymphocytes B \& T) \\ \hline
\textbf{Spécificité} & Faible (PAMPs : motifs conservés ex. LPS, flagelline, ARN viral) & \textbf{Élevée} (Épitope unique, diversité $>10^7$) \\ \hline
\textbf{Récepteurs} & Germline (codés dans génome, identiques chez tous) : TLR, NLR, etc. & \textbf{Somatiques} (générés par recombinaison V(D)J, uniques par clone) \\ \hline
\textbf{Vitesse} & Immédiate à Rapide (min -- heures) & Lente (jours -- semaines pour primaire) \\ \hline
\textbf{Mémoire} & \textbf{Aucune} (réponse identique à chaque fois) & \textbf{OUI} (Cellules mémoire : réponse secondaire rapide/forte) \\ \hline
\textbf{Effecteurs} & Phagocytes, NK, Complément, Cytokines, Protéines phase aiguë & \textbf{Anticorps (B/Plasmocytes)}, \textbf{LT Cytotoxiques (CD8)}, \textbf{LT Auxiliaires (CD4)} \\ \hline
\textbf{Transmissibilité} & Non (sauf transmission passive mère-enfant IgG placenta / IgA lait) & Non (sauf transmission passive Ac) \\ \hline
\end{tabularx}
\end{center}

\courssub{Exercices résolus type BEPC}

\begin{exoresolu}[Analyse d'une courbe de vaccination antitétanique]
\textbf{Énoncé :} Un adolescent reçoit une primo-vaccination tétanos (2 injections IM à 1 mois d'intervalle : J0 et J30), puis un rappel à 1 an (J365). On dose les antitoxines tétaniques (IgG) sériques (UI/mL).
\begin{center}
\begin{tabular}{|c|c|c|c|c|c|c|c|c|}\hline
Jour & J0 & J14 & J30 (Inj 2) & J45 & J90 & J365 (Rappel) & J379 & J400 \\ \hline
Taux (UI/mL) & 0 & 0,05 & 0,1 & 5 & 2 & 1,5 & 120 & 80 \\ \hline
\end{tabular}
\end{center}
\textbf{Questions :}
\begin{enumerate}
    \item Identifier les périodes correspondant à la réponse primaire et à la réponse secondaire.
    \item Expliquer les différences de cinétique et de titre observées entre le pic à J45 et le pic à J379.
    \item Interpréter la présence d'un taux résiduel (1,5 UI/mL) à J365, juste avant le rappel.
\end{enumerate}

\tcblower
\textbf{Solution détaillée :}
\begin{enumerate}
    \item \textbf{Réponse Primaire :} Période J0 $\to$ J90 environ. Pic à \textbf{J45 (5 UI/mL)}. La 2$^{\text{ème}}$ injection à J30 agit comme un rappel de la primo-vaccination (amplification réponse primaire).
    \item \textbf{Réponse Secondaire :} Période J365 $\to$ J400. Pic à \textbf{J379 (120 UI/mL)}.
    \item \textbf{Comparaison J45 vs J379 :}
    \begin{itemize}
        \item \textbf{Latence :} J45 : $\approx$ 15 jours post-injection (temps activation naïve, expansion clonale). J379 : $\approx$ 14 jours post-rappel \textbf{MAS pas de latence réelle} (réponse détectable dès J370 env.), car cellules mémoire déjà présentes et pré-activées.
        \item \textbf{Titre :} J379 (120) est $\times 24$ plus haut que J45 (5). Activation massive clones mémoire $\to$ plus de plasmocytes, sécrétion plus intense.
        \item \textbf{Qualité Ac :} J45 : mélange IgM/IgG, affinité moyenne. J379 : \textbf{Presque uniquement IgG de haute affinité} (maturation d'affinité achevée dans centres germinatifs lors primo).
    \end{itemize}
    \item \textbf{Taux résiduel J365 (1,5 UI/mL) :} Correspond aux \textbf{anticorps à longue durée de vie} (IgG, demi-vie $\approx$ 21--28 j) produits par des \textbf{plasmocytes à vie longue} installés dans la moelle osseuse (niches de survie : IL-6, APRIL, CXCL12). Ils assurent une \textbf{protection humorale immédiate} (taux $>$ seuil protecteur 0,1 UI/mL OMS) en attendant la montée du rappel. Preuve de persistance de la mémoire humorale.
\end{enumerate}
\end{exoresolu}

\begin{exoresolu}[Cas clinique : Plaie souillée et statut vaccinal]
\textbf{Énoncé :} Un enfant de 7 ans, non vacciné contre le tétanos, se blesse gravement au pied avec un clou rouillé dans un champ (sol contaminé \emph{Clostridium tetani}). Il arrive au centre de santé 6h après.
\textbf{Questions :}
\begin{enumerate}
    \item Quelle est la conduite à tenir \textbf{immédiate} (curative/préventive d'urgence) ? Justifier.
    \item Quelle conduite tenir \textbf{ultérieurement} pour une protection durable ? Justifier.
\end{enumerate}

\tcblower
\textbf{Solution :}
\begin{enumerate}
    \item \textbf{Sérothérapie d'urgence : Injection de Sérum Antitetanique (SAT) par voie IM (ou IV si forme sévère), dose selon notice (ex. 1500--3000 UI).}
    \textbf{Justification :} L'enfant n'a \textbf{aucune mémoire immunitaire} anti-tétanos (non vacciné). L'incubation du tétanos est courte (3--21 j, moy. 10 j). La vaccination active (anatoxine) mettrait $\approx$ 2 semaines pour produire des Ac protecteurs $\to$ \textbf{trop lent}. Le SAT apporte des \textbf{Ac tout faits (immunité passive)} qui neutralisent \textbf{immédiatement} la toxine tétanique \emph{libre} dans l'interstitium/sang avant qu'elle n'atteigne le SNC. \emph{Note : Ne neutralise pas la toxine déjà fixée sur les neurones.}
    \item \textbf{Primo-vaccination complète :} Injection simultanée (site différent, autre membre) de la \textbf{1$^{\text{ère}}$ dose d'anatoxine tétanique (souvent DTC/HepB/Hib)}. Puis respecter le calendrier : 2$^{\text{ème}}$ dose à 1 mois (J30), 3$^{\text{ème}}$ dose à 6 mois (J180), puis rappels (1 an, puis tous les 10--20 ans).
    \textbf{Justification :} La sérothérapie ne dure que 2--3 semaines (catabolisme Ac, pas de mémoire). Seule la \textbf{vaccination active} induit une \textbf{mémoire immunologique} (cellules B et T mémoire) pour une protection \textbf{durable (années)}. L'association SAT + Vaccin (sites différents) est la conduite standard OMS pour plaie à risque chez non vacciné.
\end{enumerate}
\end{exoresolu}

\courssub{Exercices d'application (À traiter par l'élève)}

\begin{exercice}[Synthèse : Le VIH/Sida et l'effondrement de l'immunité]
Le VIH (Virus de l'Immunodéficience Humaine) cible et détruit préférentiellement les \textbf{lymphocytes T CD4+} (auxiliaires).
\begin{enumerate}
    \item Rappeler les deux rôles principaux des LT CD4+ (Th) dans la coordination de la réponse immunitaire.
    \item En déduire pourquoi la chute du taux de CD4+ ($< 200/\text{mm}^3$ au stade SIDA) entraîne \textbf{simultanément} :
    \begin{itemize}
        \item Une défaillance de l'immunité \textbf{cellulaire} (ex. : tuberculose disséminée, pneumocystose).
        \item Une défaillance de l'immunité \textbf{humorale} (ex. : infections bactériennes encapsulées récidivantes, mauvaise réponse aux vaccins).
    \end{itemize}
    \item Expliquer le paradoxe suivant : les patients SIDA avancés ont souvent un \textbf{taux d'IgG total très élevé} (hypergammaglobulinémie) mais sont \textbf{très sensibles aux nouvelles infections}.
\end{enumerate}
\end{exercice}

\begin{corrige}[Exercice : VIH/Sida]
\begin{enumerate}
    \item \textbf{Rôles des LT CD4+ (Chefs d'orchestre) :}
    \begin{itemize}
        \item \textbf{Aide à l'immunité cellulaire :} Sécrètent IFN-$\gamma$ $\to$ activent les macrophages (tuerie \emph{Mycobacterium}, \emph{Leishmania}) ; aident la différenciation/activation des LT CD8+ cytotoxiques (contrôle virus, tumeurs).
        \item \textbf{Aide à l'immunité humorale :} (Via sous-type Tfh -- Folliculaire) Contact CD40L-CD40 + Cytokines (IL-4, IL-21) $\to$ \textbf{Indispensables} pour : Switch de classe (IgM $\to$ IgG/IgA), Maturation d'affinité, Génération B mémoire, Plasmocytes vie longue.
    \end{itemize}
    \item \textbf{Conséquences de la déplétion CD4+ :}
    \begin{itemize}
        \item \textbf{Immunité cellulaire :} Pas d'IFN-$\gamma$ $\to$ Macrophages non activés $\to$ prolifération intracellulaire \emph{M. tuberculosis} (TB disséminée), \emph{Pneumocystis jirovecii} (pneumonie). Pas d'aide CD8 $\to$ réactivation virus (CMV, Herpès, VZV).
        \item \textbf{Immunité humorale :} Pas d'aide Tfh $\to$ Blocage au stade IgM, pas de switch, pas d'affinité, pas de mémoire B $\to$ Incapacité à faire des Ac protecteurs contre \emph{nouveaux} Ag (pneumocoque, méningocoque, \emph{H. influenzae}) ni répondre aux vaccins (inactivés ou atténués vivants contre-indiqués).
    \end{itemize}
    \item \textbf{Hypergammaglobulinémie polyclonale :} Stimulation \textbf{polyclonalité} (non spécifique) des B résiduels par le VIH lui-même (gp120, ARN viral via TLR7/8) + inflammation chronique (TNF-$\alpha$, IL-6, BAFF). $\to$ Activation de \textbf{multiples clones B non spécifiques} $\to$ Production massive d'IgG \textbf{"bruit de fond"} (auto-Ac, Ac anti-VIH non neutralisants, Ac anti-anciens microbes). Mais \textbf{absence de réponse spécifique neuve} (pas de centre germinatif fonctionnel sans Tfh). D'où : "Beaucoup d'anticorps, mais pas les bons, pas de mémoire nouvelle".
\end{enumerate}
\end{corrige}

\begin{exercice}[Application : Choix thérapeutique]
Un agriculteur de 45 ans, dont la vaccination tétanos n'est pas à jour (dernière dose il y a 15 ans), se blesse profondément à la main avec une machette souillée de terre (risque tétanos élevé). Il consulte 4h après.
\begin{enumerate}
    \item Faut-il lui administrer du Sérum Antitetanique (SAT) ? Justifier en comparant la cinétique de l'immunité passive et active chez ce sujet qui a une mémoire immunitaire résiduelle.
    \item Faut-il lui faire un rappel vaccinal (anatoxine) ? Justifier.
    \item Préciser les conditions d'administration (sites, voie) si les deux sont indiqués.
\end{enumerate}
\end{exercice}

\begin{corrige}[Exercice : Choix thérapeutique]
\begin{enumerate}
    \item \textbf{OUI, SAT indiqué.} Bien qu'il ait été vacciné (mémoire existante), le délai depuis le dernier rappel (15 ans) est trop long : le taux d'Ac circulants est probablement \textbf{sous le seuil protecteur (0,1 UI/mL)}. La réponse de rappel (secondaire) prendrait $\approx$ 3--5 jours pour atteindre un taux protecteur. L'incubation tétanos peut être courte (quelques jours). Le SAT assure une \textbf{neutralisation immédiate} de la toxine libérée \emph{maintenant} (immunité passive), comblant le "trou" de protection avant que la réponse anamnestique ne monte.
    \item \textbf{OUI, Rappel vaccinal (Anatoxine) INDISPENSABLE.} Le SAT ne dure que 2--3 semaines. Seule la réactivation de la mémoire (réponse secondaire) par le rappel vaccinal garantit une \textbf{protection durable} (remontée taux Ac $>$ seuil protecteur en quelques jours, puis maintien par plasmocytes vie longue et mémoire réactivée).
    \item \textbf{Administration :} \textbf{Sites différents, membres différents.} Ex: SAT dans fesse droite (voie IM profonde), Anatoxine dans deltoïde gauche (voie IM). \textbf{JAMAIS même seringue, JAMAIS même site} (risque neutralisation de l'anatoxine par les Ac du SAT, inflammation locale majeure). Surveillance allergie (SAT d'origine équine souvent).
\end{enumerate}
\end{corrige}

\begin{exercice}[Analyse de document : Phagocytose et évasion bactérienne]
On observe au microscope des macrophages en culture mis en présence de deux souches de \emph{Staphylococcus aureus} :
\begin{itemize}
    \item \textbf{Souche A :} Non encapsulée.
    \item \textbf{Souche B :} Encapsulée (capsule polysaccharidique épaisse).
\end{itemize}
Au bout de 30 min, on compte le nombre de bactéries internalisées par 100 macrophages.
\begin{center}
\begin{tabular}{|c|c|}\hline
\textbf{Souche} & \textbf{Nombre moyen de bactéries / 100 macrophages} \\ \hline
A (Non encapsulée) & 320 \\ \hline
B (Encapsulée) & 12 \\ \hline
\end{tabular}
\end{center}
\begin{enumerate}
    \item Comparer les résultats et en déduire le rôle de la capsule dans la virulence de \emph{S. aureus}.
    \item Expliquer le mécanisme moléculaire par lequel la capsule empêche la phagocytose (notion d'opsonisation).
    \item Comment l'immunité spécifique (anticorps) permet-elle de lever ce blocage ? Nommer ce phénomène.
\end{enumerate}
\end{exercice}

\begin{corrige}[Exercice : Phagocytose et capsule]
\begin{enumerate}
    \item \textbf{Résultat :} La souche non encapsulée (A) est massivement phagocytée (320/100 macro). La souche encapsulée (B) l'est très peu (12/100 macro). \textbf{Conclusion :} La capsule polysaccharidique est un \textbf{facteur de virulence majeur} qui protège la bactérie contre la phagocytose (deuxième ligne innée).
    \item \textbf{Mécanisme :} La capsule masque les ligands de surface bactériens (PAMPs : peptidoglycane, acide teichoïque) reconnus par les récepteurs de reconnaissance (PRR) des macrophages (ex. récepteurs scavenger, TLR2). Elle crée une barrière physique stérique et une charge négative repoussant la membrane du phagocyte. \textbf{Résultat :} Pas de reconnaissance $\to$ Pas d'adhérence $\to$ Pas d'ingestion.
    \item \textbf{Levée du blocage par l'immunité spécifique :} Les \textbf{anticorps anti-capsule} (produits par les plasmocytes suite activation B avec aide T \emph{si vaccin conjugué}, ou réponse TI-2 faible) se lient à la capsule. Leur fragment Fc est reconnu par les \textbf{récepteurs Fc$\gamma$R (CD64, CD32, CD16)} du macrophage. \textbf{Phénomène : OPSONISATION.} L'anticorps fait le "pont" entre l'antigène (capsule) et le phagocyte $\to$ Phagocytose efficace restaurée.
\end{enumerate}
\end{corrige}

\begin{aretenir}[Bilan de l'Unité : Réponse Immunitaire (3ème)]
\begin{enumerate}
    \item \textbf{1ère ligne (Barrières) :} Peau (kératine, pH acide, flore) + Muqueuses (mucus, cils, pH extrêmes, lysozyme, flore). \textbf{Non spécifique, immédiate, pas de mémoire.}
    \item \textbf{2ème ligne (Innée cellulaire/humorale) :} Inflammation (vasodilatation, œdème, recrutement) + Phagocytose (Neutrophiles, Macrophages : chimiotaxie, ingestion, digestion phagolysosomale) + NK + Complément. \textbf{Non spécifique, rapide (heures), pas de mémoire.}
    \item \textbf{3ème ligne (Adaptative / Spécifique) :} Lymphocytes B (Humoral : Ac) et T (Cellulaire : CD8 tueurs, CD4 chefs d'orchestre). \textbf{Spécifique, lente (jours), MÉMOIRE.}
    \item \textbf{Sélection clonale :} Ag sélectionne clone spécifique $\to$ Expansion $\to$ Effecteurs (court terme) + Mémoire (long terme).
    \item \textbf{Vaccination (Active) :} Ag $\to$ Immunité propre + Mémoire $\to$ Prévention durable. \textbf{Sérothérapie (Passive) :} Ac $\to$ Immunité empruntée $\to$ Urgence curative, protection temporaire.
    \item \textbf{VIH/Sida :} Destruction CD4+ $\to$ Effondrement coordination (Cellulaire + Humorale) $\to$ Infections opportunistes.
\end{enumerate}
\end{aretenir}

\coursec{Coopération, régulation et applications concrètes : Vaccination et Sérothérapie}

Dans la section précédente, nous avons identifié les acteurs de l'immunité spécifique : les lymphocytes B (producteurs d'anticorps) et les lymphocytes T (cytotoxiques CD8+ et auxiliaires CD4+). Nous avons vu que chaque lymphocyte ne reconnaît qu'un seul épitope. Mais comment ces cellules communiquent-elles pour monter une réponse efficace, adaptée à la nature de l'agresseur ? Comment la réponse s'arrête-t-elle ? Et comment l'homme exploite-t-il ces mécanismes pour se protéger (vaccination) ou traiter en urgence (sérothérapie) ?

\courssub{La coopération lymphocyte T auxiliaire -- Lymphocyte B : clé de la réponse humorale de qualité}

\emph{Observation :} Les personnes dont les lymphocytes T CD4+ sont détruits (par le VIH au stade SIDA) produisent peu d'anticorps de classe IgG, IgA, IgE, mais conservent une production d'IgM. Elles sont très sensibles aux bactéries encapsulées (\emph{Streptococcus pneumoniae}, \emph{Haemophilus influenzae}) qui nécessitent des anticorps IgG opsonisants pour être éliminées. Cela montre que les lymphocytes T CD4+ sont indispensables pour une réponse anticorps de \textbf{haute qualité} (autres classes qu'IgM, mémoire).

\begin{experience}[Mise en évidence de la coopération T-B (expérience classique simplifiée)]
\textbf{Protocole :} On met en culture des lymphocytes B seuls, des lymphocytes T CD4+ seuls, ou un mélange des deux, en présence d'un antigène protéique (ex. toxine tétanique). On mesure les anticorps produits après quelques jours.
\textbf{Résultats :}
\begin{itemize}
    \item B seuls + Ag : Faible production d'IgM, pas d'IgG, pas de mémoire.
    \item T seuls + Ag : Pas d'anticorps.
    \item B + T CD4+ + Ag : Forte production d'IgM \textbf{et} d'IgG (et IgA), apparition de lymphocytes B mémoire.
\end{itemize}
\textbf{Interprétation :} La réponse anticorps de classe IgG, IgA, IgE (dite \textbf{T-dépendante}) nécessite l'aide des lymphocytes T CD4+ activés. Le lymphocyte B doit présenter l'antigène au lymphocyte T CD4+.
\textbf{Conclusion :} La coopération T-B est un dialogue cellulaire indispensable pour générer une immunité humorale efficace (commutation de classe, maturation de l'affinité, mémoire).
\end{experience}

\begin{popfigure}
\begin{tikzpicture}[
    node distance=1.5cm and 2cm,
    cell/.style={draw, rounded corners, fill=popcream, thick, minimum width=2.8cm, minimum height=1.6cm, align=center},
    mol/.style={draw, ellipse, fill=popblueL, thick, minimum width=1.6cm, minimum height=0.9cm, align=center},
    arrow/.style={-Stealth, thick, popdark},
    cytokine/.style={-Stealth, dashed, poppink, thick},
    label/.style={font=\small, text width=3.5cm, align=center}
]
% Lymphocyte B
\node[cell, align=center] (LB){Lymphocyte B\\ Présentateur d'Ag\\ (MHC II + Peptide)};
\node[mol, below=0.4cm of LB] (BCR) {BCR (Récepteur Ag)};

% Lymphocyte T CD4+
\node[cell, right=of LB, align=center] (LT){Lymphocyte T CD4+\\ (Auxiliaire / Th2)\\ TCR + CD4};
\node[mol, above=0.4cm of LT, align=center] (Cyt){Cytokines\\ (ex. IL-4, IL-21)};

% Interaction
\draw[arrow] (LB.east) -- node[above, label, align=center]{Reconnaissance\\ Peptide-MHC II par TCR} (LT.west);
\draw[cytokine] (Cyt.south) -- node[right, label, align=center]{Signaux chimiques\\ (Cytokines)} (LT.north);

% Outcomes
\node[draw, rounded corners, fill=popgreenL, thick, minimum width=7cm, minimum height=2.2cm, below=1.2cm of LB.south west, anchor=north west, align=left] (Outcomes) {
    \textbf{Conséquences dans le Lymphocyte B :}\\
    1. \textbf{Prolifération} (expansion clonale)\\
    2. \textbf{Commutation de classe} : IgM $\rightarrow$ IgG, IgA, IgE\\
    3. \textbf{Augmentation de l'affinité} des anticorps\\
    4. \textbf{Différenciation} : Plasmocytes (usine à Ac) + \textbf{Lymphocytes B Mémoire}
};

\draw[arrow, poporange] (LB.south) -- (Outcomes.north west);
\draw[arrow, poporange] (LT.south) -- (Outcomes.north east);
\end{tikzpicture}
\legende{Schéma de la coopération T-B (antigène T-dépendant). Le lymphocyte B internalise l'antigène via son BCR, le présente sur MHC II. Le lymphocyte T CD4+ spécifique reconnaît le complexe et sécrète des cytokines (IL-4, etc.). Cela déclenche la commutation de classe, la maturation de l'affinité et la formation de cellules mémoire.}
\end{popfigure}

\begin{definition}[Antigène T-dépendant vs T-indépendant]
Un \textbf{antigène T-dépendant} (protéines) induit une réponse nécessitant l'aide des lymphocytes T CD4+ : commutation de classe (IgG, IgA, IgE), maturation de l'affinité, mémoire immunologique forte.
Un \textbf{antigène T-indépendant} (polysaccharides capsulaires bactériens, LPS) stimule directement les lymphocytes B : réponse rapide mais \textbf{principalement IgM}, \textbf{pas de commutation de classe}, \textbf{pas de mémoire} durable, \textbf{peu efficace chez l'enfant de moins de 2 ans}.
\end{definition}

\begin{attention}[Point clé : Les vaccins conjugués]
Les vaccins contre \emph{Haemophilus influenzae} b (Hib), le pneumocoque et le méningocoque (utilisés au PEV Cameroun) sont des \textbf{vaccins conjugués}. Ils couplent le polysaccharide (T-indépendant) à une protéine porteuse (toxine tétanique...). Cela transforme la réponse en réponse \textbf{T-dépendante} : production d'IgG, mémoire, efficacité \textbf{chez le nourrisson}. C'est une application directe de la coopération T-B.
\end{attention}

\courssub{Régulation de la réponse immunitaire : les freins physiologiques}

Une réponse immunitaire non contrôlée détruit les tissus sains (maladies auto-immunes, allergies, tempêtes cytokiniques). L'organisme possède des mécanismes pour \textbf{stopper la réponse} une fois le pathogène éliminé et pour \textbf{tolérer le "soi"}.

\begin{propriete}[Principaux mécanismes de régulation]
\begin{itemize}
    \item \textbf{Lymphocytes T régulateurs (Treg) :} Sous-population de lymphocytes T CD4+ (marqueurs CD25, FoxP3). Ils sécrètent des cytokines inhibitrices (\textbf{IL-10}, \textbf{TGF-$\beta$}) et freinent l'activation des lymphocytes T effecteurs et des cellules présentatrices d'antigène. Essentiels pour la tolérance aux antigènes du "soi" et l'arrêt de la réponse après guérison.
    \item \textbf{Molécules de contrôle (Checkpoints) :} Récepteurs inhibiteurs sur les lymphocytes T activés (ex. \textbf{CTLA-4}, \textbf{PD-1}). Leur liaison à des ligands sur les cellules cibles ou présentatrices envoie un signal "STOP" (frein précoce pour CTLA-4 dans le ganglion, frein tardif pour PD-1 dans les tissus).
\end{itemize}
\end{propriete}

\begin{savaistu}[Application médicale : Immunothérapie du cancer (Nobel 2018)]
Certains cancers "éteignent" les lymphocytes T en exprimant PD-L1 (qui se lie à PD-1). Des médicaments (anticorps monoclonaux anti-PD-1 ou anti-CTLA-4) bloquent ces freins ("inhibiteurs de checkpoints"), réactivant l'immunité anti-tumorale. Effet secondaire : risque de maladies auto-immunes (levée de la tolérance).
\end{savaistu}

\courssub{La Vaccination : mimer l'infection pour générer la mémoire}

\begin{definition}[Vaccin et Adjuvant]
Un \textbf{vaccin} est une préparation antigénique qui induit une \textbf{immunité active protectrice} (lymphocytes B/T mémoire + anticorps circulants) \textbf{sans provoquer la maladie}.
Un \textbf{adjuvant} (ex. sels d'aluminium "Alum") est une substance ajoutée aux vaccins inactivés/sous-unitaires pour \textbf{stimuler l'immunité innée} (signal danger, inflammation, maturation des cellules présentatrices) $\rightarrow$ meilleure activation des lymphocytes T. Sans adjuvant, ces vaccins sont faiblement immunogènes.
\end{definition}

\begin{popfigure}
\begin{center}
\begin{tabularx}{\linewidth}{|l|X|X|X|X|}
\hline
\rowcolor{popblueL}
\textbf{Type de vaccin} & \textbf{Principe / Antigène} & \textbf{Avantages} & \textbf{Inconvénients / Précautions} & \textbf{Exemples PEV Cameroun} \\
\hline
\textbf{Vivants atténués} & Souche pathogène mutée, réplication limitée, non virulente. & Immunité forte (humorale + cellulaire), durable, 1-2 doses. Mimétisme infection naturelle. & Risque théorique de réversion (rare), \textbf{contre-indiqué grossesse/immunodéprimés}, chaîne du froid stricte. & \textbf{BCG}, \textbf{Rougeole}, \textbf{Fièvre jaune}, \textbf{Rotavirus} (oral), \textbf{Polio OPV}. \\
\hline
\textbf{Inactivés (tués)} & Pathogène entier tué (formol, chaleur). Pas de réplication. & Sûr (pas de réversion), stable, utilisable grossesse/immunodéprimés. & Immunité plus faible (souvent humorale seule), \textbf{adjuvant obligatoire}, rappels fréquents. & \textbf{Rage}, \textbf{Polio IPV} (injectable), \textbf{Grippe}. \\
\hline
\textbf{Sous-unitaires / Recombinants / Conjugués} & Protéine purifiée (HBsAg), VLP (HPV), ou Polysaccharide + Protéine (Pneumo, Hib, Méningo). & Très sûr, bien défini. Conjugués = réponse T-dépendante chez nourrisson. & Faiblement immunogène $\rightarrow$ \textbf{adjuvant + rappels}. & \textbf{Hépatite B}, \textbf{HPV}, \textbf{Pneumocoque}, \textbf{Hib}, \textbf{Méningocoque}. \\
\hline
\textbf{Toxoïdes} & Toxine bactérienne inactivée (formol) $\rightarrow$ perd toxicité, garde immunogénicité. & Cible l'effet pathogène (toxine). & Nécessite adjuvant (Alum) + rappels (tous les 10-20 ans). & \textbf{Tétanos}, \textbf{Diphtérie} (dans Pentavalent). \\
\hline
\end{tabularx}
\end{center}
\legende{Tableau comparatif des types de vaccins utilisés au PEV Cameroun. La stratégie combine vivants atténués (BCG, Rougeole, FJ, Rota, OPV), inactivés (IPV, Rage), sous-unitaires/conjugués (HepB, HPV, Pneumo, Hib, Méningo) et toxoïdes (Tétanos, Diphtérie).}
\end{popfigure}

\begin{popfigure}
\begin{tikzpicture}
\begin{axis}[
    width=\linewidth, height=7.5cm,
    xlabel={Temps (semaines)},
    ylabel={Taux d'anticorps (échelle logarithmique)},
    xmin=0, xmax=50,
    ymin=0, ymax=5.5,
    xtick={0,4,8,12,16,20,24,36,48},
    xticklabels={0, 1m, 2m, 3m, 4m, 5m, 6m, 9m, 12m},
    ytick={0,1,2,3,4,5},
    yticklabels={0, $10^1$, $10^2$, $10^3$, $10^4$, $10^5$},
    legend pos=north west,
    grid=major,
]
% Infection naturelle
\addplot[popcurve, popblue, very thick, domain=0:50, samples=200] 
    {2.5*exp(-(x-10)^2/30) + 0.5*exp(-(x-30)^2/100) + 0.3};
\addlegendentry{Infection naturelle}

% Vaccination 2 doses (Prime + Boost)
\addplot[popcurve, poporange, very thick, domain=0:50, samples=200] 
    {1.5*exp(-(x-4)^2/10) + 3.5*exp(-(x-12)^2/15) + 1.0*exp(-(x-30)^2/100) + 0.5};
\addlegendentry{Vaccination (Prime + Rappel)}

% Seuil protection
\addplot[dashed, popred, thick, domain=0:50] {1.5};
\addlegendentry{Seuil protecteur}

\node[popblue, anchor=south east] at (axis cs:10,3.5) {Pic primaire};
\node[poporange, anchor=south west] at (axis cs:12,4.5) {Pic secondaire (Rappel)};
\node[popmuted, anchor=north] at (axis cs:25,1.6) {Niveau mémoire protecteur};
\end{axis}
\end{tikzpicture}
\legende{Kinétique de la réponse anticorps : Infection naturelle vs Vaccination. La vaccination (prime-boost) mime une ré-infection : le rappel génère une réponse secondaire plus rapide, plus intense (affinité $\uparrow$, IgG $\uparrow$) et plus durable, maintenant le taux au-dessus du seuil protecteur.}
\end{popfigure}

\begin{methode}[Analyser un calendrier vaccinal (ex. PEV Cameroun)]
\begin{enumerate}
    \item \textbf{Identifier le vaccin} et son \textbf{type} (vivant atténué, inactivé, sous-unitaire/conjugué, toxoïde).
    \item \textbf{Noter l'âge} de la première dose (fenêtre de vulnérabilité vs anticorps maternels IgG transplacentaires).
    \item \textbf{Compter le nombre de doses} du primo-vaccinage (souvent 3 doses pour inactivés/sous-unitaires : 6, 10, 14 semaines).
    \item \textbf{Repérer les rappels (boosters)} : essentiels pour entretenir la mémoire (ex. Rougeole 9 mois + 15-18 mois ; DTC rappels 15-18 mois, 4-7 ans, 12-15 ans).
    \item \textbf{Vérifier les contre-indications} : Immunodépression $\rightarrow$ pas de vaccins vivants (BCG, Rougeole, Fièvre jaune, Rotavirus, OPV). Grossesse $\rightarrow$ pas de vaccins vivants.
\end{enumerate}
\end{methode}

\begin{exemplebox}[Calendrier PEV Cameroun 2024 (Extrait simplifié)]
\begin{center}
\begin{tabular}{|c|l|l|}
\hline
\textbf{Âge} & \textbf{Vaccins} & \textbf{Type / Notes} \\
\hline
À la naissance & BCG (ID), HepB (IM) -- dose 0 & Vivant atténué / Recombinant (dans 24h) \\
\hline
6 sem (1,5 mois) & Pentavalent (DTC-HepB-Hib) D1, Pneumo D1, Rota D1, Polio OPV D1 / IPV D1 & Toxoïdes+Conjugués / Conjugué / Vivant atténué oral / Vivant atténué + Inactivé \\
\hline
10 sem (2,5 mois) & Pentavalent D2, Pneumo D2, Rota D2, Polio OPV D2 & Rappels primo-vaccinage \\
\hline
14 sem (3,5 mois) & Pentavalent D3, Pneumo D3, Polio OPV D3 / IPV D2 & Fin primo-vaccinage \\
\hline
9 mois & Rougeole, Fièvre jaune, Vitamine A & Vivants atténués (couverture cible $>95\%$) \\
\hline
15-18 mois & Rougeole (D2), DTC (Rappel 1) & Rappels mémoire \\
\hline
9-14 ans (Filles) & HPV (2 doses, 6 mois écart) & VLP Recombinant (Prévention cancer col utérus) \\
\hline
\end{tabular}
\end{center}
\emph{ID = Intradermique, IM = Intramusculaire. La polio utilise une stratégie séquentielle IPV/OPV pour l'éradication.}
\end{exemplebox}

\begin{attention}[Idées reçues dangereuses]
\begin{itemize}
    \item \textbf{FAUX :} "Les vaccins fatiguent/surchargent le système immunitaire du bébé."
    \item \textbf{FAITS :} Un nourrisson rencontre $\sim 10^6$ à $10^{10}$ antigènes \emph{chaque jour} (commensaux, alimentation). Le calendrier vaccinal complet représente $\sim \mathbf{10^2}$ à $\mathbf{10^3}$ antigènes \emph{au total} sur plusieurs mois. C'est \textbf{négligeable} ($<0,001\%$).
    \item \textbf{FAUX :} "Les antibiotiques soignent la grippe / le rhume / le COVID."
    \item \textbf{FAITS :} Les antibiotiques ciblent les \textbf{bactéries}. Ils sont \textbf{inactifs sur les virus}. Usage injustifié $\rightarrow$ résistances (SAMR, ESBL) -- fléau majeur au Cameroun.
\end{itemize}
\end{attention}

\courssub{La Sérothérapie / Immunothérapie passive : l'urgence et l'immunodéprimé}

\begin{definition}[Immunité passive (Sérothérapie)]
Administration d'\textbf{anticorps préformés} (polyclonaux ou monoclonaux) pour une \textbf{protection immédiate}. \textbf{Pas d'activation du système immunitaire du receveur} $\rightarrow$ \textbf{pas de mémoire}, \textbf{durée de vie limitée} à la demi-vie des IgG injectées ($\sim \mathbf{21\text{--}28\text{ jours}}$).
\end{definition}

\begin{propriete}[Comparaison Immunité Active vs Passive]
\begin{center}
\begin{tabular}{|l|c|c|}
\hline
\rowcolor{popblueL}
\textbf{Critère} & \textbf{Immunité Active (Vaccin / Infection)} & \textbf{Immunité Passive (Sérum / Ig Maternelles / mAb)} \\
\hline
\textbf{Délai de protection} & Lent (10--14 j primo, 2--5 j rappel) & \textbf{Immédiat (heures)} \\
\hline
\textbf{Durée} & Longue (années -- vie entière) & Courte (semaines -- mois : $t_{1/2}$ IgG $\approx$ 21 j) \\
\hline
\textbf{Mémoire immunologique} & \textbf{OUI} & \textbf{NON} \\
\hline
\textbf{Risque maladie} & Faible à nul & Nul (produit purifié) \\
\hline
\textbf{Indications principales} & Prévention routine, éradication, contrôle épidémies & \textbf{Urgence} (tétanos, morsure serpent, rage post-expo), \textbf{Immunodéprimés}, Prophylaxie RSV, Maladies auto-immunes/cancers (mAb thérapeutiques) \\
\hline
\end{tabular}
\end{center}
\end{propriete}

\begin{exemplebox}[Applications concrètes au Cameroun]
\begin{itemize}
    \item \textbf{Sérum antivenimeux (SAV) :} Anticorps polyclonaux (chevaux/ovins immunisés contre venins de vipères, cobras, mambas). \textbf{Urgence vitale} en zone rurale (Nord, Adamaoua, Est). Administration IV précoce. Risque : choc anaphylactique (protéines hétérologues) $\rightarrow$ test cutané + surveillance.
    \item \textbf{Immunoglobulines antitétaniques (IGAT) :} Humaines (donneurs vaccinés). Indiquées : Blessure sale + vaccination incomplète/inconnue + $\ge$ 10 ans depuis dernier rappel. \textbf{Ne remplace pas le vaccin} (on fait les deux : sérothérapie immédiate + vaccination active en site distinct).
    \item \textbf{Anticorps monoclonaux (mAb) :} Palivizumab (anti-RSV, prématurés), Anti-TNF (polyarthrite, Crohn), Anti-PD-1 (cancers). Coût très élevé $\rightarrow$ accès limité, mais révolution thérapeutique.
\end{itemize}
\end{exemplebox}

\courssub{Immunité collective et hygiène : protéger la communauté}

\begin{popfigure}
\begin{tikzpicture}[
    scale=0.85,
    every node/.style={font=\small},
    person/.style={circle, draw=popdark, thick, minimum size=5.5mm},
    infected/.style={person, fill=popred!80},
    immune/.style={person, fill=popgreen!70},
    susceptible/.style={person, fill=poporange!50},
    arrow/.style={-Stealth, popred, thick},
    protect/.style={-Stealth, popgreen, thick, dashed}
]
% Population grid
\foreach \x in {0,1,2,3,4} {
    \foreach \y in {0,1,2,3} {
        \pgfmathsetmacro{\r}{random(0,100)}
        \ifnum\r<85
            \node[immune] (p-\x-\y) at (\x*1.3, -\y*1.3) {};
        \else
            \node[susceptible] (p-\x-\y) at (\x*1.3, -\y*1.3) {};
        \fi
    }
}
% Index case
\node[infected] at (2.6, -2.6) {};

% Blocked transmission arrows
\foreach \x/\y in {1/2, 3/2, 2/1, 2/3} {
    \draw[protect] (2.6, -2.6) -- (\x*1.3, -\y*1.3) node[midway, above, sloped, font=\tiny, popgreen] {Bloqué};
}

% Legend
\begin{scope}[shift={(7.0,-1.5)}]
\node[immune] (l1) at (0,0) {}; \node[right=2mm of l1] {Vacciné / Immunisé ($>90\%$)};
\node[susceptible] (l2) at (0,-0.8) {}; \node[right=2mm of l2] {Susceptible (non vacciné, échec, immunodéprimé)};
\node[infected] (l3) at (0,-1.6) {}; \node[right=2mm of l3] {Cas index (Malade)};
\draw[protect] (0,-2.5) -- (1.5,-2.5) node[right, font=\tiny] {Chaîne de transmission brisée};
\end{scope}

\node[align=center, font=\bfseries\large, popdark] at (3.2, 1.6) {Immunité Collective (Herd Immunity)};
\node[align=center, font=\small, popmuted] at (3.2, 1.1) {Seuil critique $P_c = 1 - 1/R_0$. Rougeole ($R_0\sim15$) $\rightarrow P_c > 93\%$. Diphtérie ($R_0\sim6$) $\rightarrow P_c > 83\%$.};
\end{tikzpicture}
\legende{Immunité collective. Lorsque la couverture vaccinale dépasse le seuil critique ($P_c = 1 - 1/R_0$), la circulation du pathogène est interrompue. Les non-vaccinés sont indirectement protégés. Au Cameroun, objectif PEV/OMS : $>90\%$ rougeole, $>80\%$ autres antigènes.}
\end{popfigure}

\begin{exemplebox}[Objectifs OMS / PEV Cameroun -- Résultats majeurs]
\begin{itemize}
    \item \textbf{Couverture cible OMS :} $\ge 90\%$ national, $\ge 80\%$ par district (équité).
    \item \textbf{Éradication Variole :} 1980 (seule maladie humaine éradiquée par vaccination).
    \item \textbf{Élimination Polio sauvage Afrique :} Certifiée août 2020. Vigilance sur poliovirus dérivés vaccinaux (cVDPV).
    \item \textbf{Tétanos maternel et néonatal (TMN) :} Éliminé comme problème de santé publique au Cameroun (2018) gr\^ace vaccination femmes enceintes + accouchements propres.
    \item \textbf{Fièvre jaune :} Campagnes de masse + routine 9 mois. Risque persistant zone forestière.
    \item \textbf{HPV :} Introduction 2020 (filles 9-14 ans, 2 doses). Objectif élimination cancer col utérus (90\% couv. vaccinale, 70\% dépistage, 90\% traitement).
\end{itemize}
\end{exemplebox}

\begin{methode}[Hygiène quotidienne : agir sur les barrières de 1ère ligne (Rappel Sections 1 \& 2)]
\begin{enumerate}
    \item \textbf{Lavage des mains à l'eau et au savon :} Mécanique (décollement biofilm/bactéries) + Chimique (savon = tensioactif $\rightarrow$ solubilise enveloppe lipidique virus \emph{enveloppés} : grippe, COVID, VIH, HBV, HCV ; déstabilise membranes bactériennes). \textbf{Durée :} 30-40 s. Moments clés : après toilettes, avant manger, après soins bébé, en rentrant.
    \item \textbf{Désinfection plaie :} Nettoyage abondant eau/savon ou sérum physiologique $\rightarrow$ élimination mécanique. Antiseptique (Bétadine\footnote{\textregistered PVP-Iode} / Eau oxygénée\footnote{$H_2O_2$} / Chlorhexidine) $\rightarrow$ chimio-antisepsie \emph{avant} inflammation. \textbf{Pas d'alcool 70\%/90\% sur plaie ouverte} (nécrose, douleur, retarde cicatrisation).
    \item \textbf{Hygiène respiratoire :} Tousser/éternuer dans coude ou mouchoir jetable $\rightarrow$ limite aérosols/gouttelettes (grippe, COVID, tuberculose, méningite). Masque si symptômes ou lieu clos mal ventilé.
\end{enumerate}
\end{methode}

\begin{exoresolu}[Analyse d'une situation vaccinale : Nourrisson de 4 mois, zone rurale Cameroun]
\textbf{Situation :} Nourrisson de 4 mois ($\sim$17 semaines) au CSI. Carnet : BCG + HepB0 à la naissance ; Pentavalent D1, Pneumo D1, Rota D1, OPV D1 à 6 sem. Le calendrier prévoit D2 à 10 sem, D3 à 14 sem. L'enfant a du retard.
\textbf{Questions :}
1. Quels vaccins administrer aujourd'hui ? Justifier.
2. Pourquoi le vaccin Rota a-t-il une limite d'âge stricte (1ère dose $<$ 15 sem, dernière dose $<$ 32 sem) ?
3. L'allaitement maternel exclusif influence-t-il la vaccination ?

\textbf{Solution détaillée :}
1. \textbf{Vaccins à administrer :} \textbf{Pentavalent D2 (ou D3 si D2 faite), Pneumo D2 (ou D3), Rota D2 (si D1 faite et âge $<$ 32 sem), OPV D2 (ou D3), IPV D2 (si non fait).} \emph{Principe :} On ne recommence jamais une série, on la \textbf{continue} (rattrapage). Intervalle minimum 4 semaines respecté. Pas de dose maximale pour rattrapage.
2. \textbf{Rotavirus :} Vaccin vivant atténué oral. Risque d'\textbf{invagination intestinale aiguë} (IAA) augmentée transitoirement (1-7 j post-vaccination) surtout après 1ère dose. L'IAA est pathologie chirurgicale du nourrisson (pic 5-9 mois). Pour minimiser risque attribuable au vaccin : 1ère dose $\le$ 14 sem 6 j (souvent 15 sem au Cameroun), dernière dose $\le$ 32 sem. \textbf{Au-delà : contre-indiqué.}
3. \textbf{Allaitement maternel :} IgA sécrétoires du lait protègent muqueuse intestinale. Elles \textbf{ne neutralisent pas} les vaccins injectables (Pentavalent, Pneumo, IPV, BCG, FJ, Rougeole, HPV). Pour vaccins \textbf{oraux vivants} (Rota, OPV), interférence \emph{théorique} possible (IgA anti-Rota/Polio dans lait), mais études cliniques : efficacité non significativement réduite. \textbf{Recommandation OMS :} Ne pas différer ni arrêter l'allaitement. Allaiter avant/après dose orale possible.
\end{exoresolu}

\begin{attention}[Erreurs de rédaction fréquentes au BEPC]
\begin{itemize}
    \item \textbf{Confondre "Vaccin" et "Sérum" :} "On vaccine un blessé avec le sérum antitétanique." $\rightarrow$ \textbf{FAUX}. On fait une \textbf{sérothérapie} (IGAT) \textbf{ET} on met à jour la \textbf{vaccination} (TTC/Td). Le sérum ne vaccinne pas.
    \item \textbf{Dire "Le vaccin donne la maladie" :} Pour vivants atténués, forme \emph{atténuée} possible (fièvre modérée, éruption discrète), ce n'est \textbf{pas la maladie sauvage}. Inactivés/sous-unitaires/ARNm \textbf{ne peuvent pas} donner la maladie.
    \item \textbf{Oublier l'adjuvant :} "Le vaccin Hépatite B contient le virus tué." $\rightarrow$ \textbf{FAUX}. C'est une protéine recombinante (HBsAg) \textbf{sans ADN viral}, \textbf{adsorbée sur sel d'aluminium (adjuvant)}.
    \item \textbf{Immunité passive = Mémoire :} "Le sérum antivenimeux protège pour la vie." $\rightarrow$ \textbf{FAUX}. Protection 3-4 semaines. Pas de lymphocytes mémoire générés.
\end{itemize}
\end{attention}

\courssub{Activités d'intégration : Synthèse et Résolution de problèmes}

\begin{exercice}[Synthèse : De la morsure à la protection durable]
Un garçon de 10 ans est mordu par un chien errant (Yaoundé). Plaie profonde, mollet. Chien non retrouvé. Carnet : BCG, Polio (complet), Pentavalent (complet), Rougeole (9 mois), Fièvre jaune (9 mois), \textbf{Pas de rappel DTC depuis 15 mois}. Pas de vaccin antirabique.
\textbf{Conduite à tenir immédiate (Urgences) :}
1. Nettoyage plaie (15 min eau/savon + antiseptique).
2. \textbf{Sérum Antirabique (SAR) -- Immunoglobulines anti-rabiques} : Infiltration autour plaie + IM site distant. (Immunité passive immédiate).
3. \textbf{Primo-vaccination Vaccin Rabique (inactivé)} : J0, J3, J7, J14, J28 (schéma Essen) ou J0, J7, J21 (Zagreb). Site distinct du SAR. (Immunité active durable).
4. \textbf{Rappel Tétanos (TTC/Td)} : Rappel $>$ 10 ans (dernier à 15 mois, âge 10 ans). Plaie sale $\rightarrow$ Indication IGAT + Rappel vaccinal.
5. Antibiothérapie prophylactique (Amoxicilline/Ac. Clavulanique) si risque infection bactérienne.
\textbf{Questions :}
\begin{enumerate}
    \item Classer les produits (SAR, Vaccin Rabique, IGAT, Rappel Tétanos) en Immunité Active ou Passive. Justifier (délai, durée).
    \item Pourquoi SAR et Vaccin Rabique sont-ils administrés \textbf{simultanément} mais en \textbf{sites distincts} ?
    \item Expliquer pourquoi l'enfant a besoin d'un rappel tétanos malgré le Pentavalent (3 doses) nourrisson.
\end{enumerate}
\end{exercice}

\begin{corrige}[Exercice : Morsure chien]
\begin{enumerate}
    \item \textbf{Immunité Passive (Immédiate, temporaire $\sim$3 sem, pas mémoire) :} SAR (Ig anti-rabiques), IGAT (Ig anti-tétaniques). \textbf{Immunité Active (Retardée 10-14j, durable années/vie, mémoire) :} Vaccin Rabique (inactivé), Rappel Tétanos (Toxoïde).
    \item \textbf{Sites distincts :} Éviter que les anticorps du SAR (passifs) ne se lient aux antigènes du vaccin rabique au site d'injection, les neutralisant (\emph{interférence antigène-anticorps}) et empêchant la présentation aux lymphocytes B/T $\rightarrow$ échec de la primo-vaccination active.
    \item \textbf{Rappel Tétanos :} Le Pentavalent nourrisson induit une mémoire, mais le taux d'anticorps antitoxine tétanique chute sous le seuil protecteur ($\sim$0,1 UI/mL) après $\sim$10 ans. La plaie sale (porte d'entrée \emph{C. tetani}) impose un \textbf{rappel vaccinal} (réactivation mémoire $\rightarrow$ pic secondaire rapide en 3-5 j) \textbf{ET} une sérothérapie (IGAT) car le délai de la réponse active (3-5 j) est incompatible avec la cinétique de la toxine tétanique (incubation 3-21 j, toxine rapide). L'IGAT neutralise la toxine \emph{immédiatement} le temps que la réponse active monte.
\end{enumerate}
\end{corrige}

\begin{aretenir}[Bilan de l'Unité : Réponse Immunitaire -- Clés pour le BEPC]
\begin{itemize}
    \item \textbf{3 lignes de défense :} 1. Barrières (Peau/Muqueuses : physique, chimique, microbiote). 2. Innée non-spécifique (Inflammation, Phagocytose, NK, Complément). 3. Adaptative Spécifique (Lymphocytes B $\rightarrow$ Ac ; Lymphocytes T CD8+ cytotoxiques, CD4+ auxiliaires).
    \item \textbf{Spécificité \& Mémoire :} Récepteurs clonaux (BCR, TCR). Sélection clonale. Lymphocytes Mémoire (B/T) = réponse secondaire plus rapide, forte, affinité $\uparrow$.
    \item \textbf{Coopération T-B :} Reconnaissance Ag (MHC II/TCR) + Cytokines (IL-4...) $\rightarrow$ Commutation classe (IgM$\rightarrow$IgG/IgA/IgE) + Maturation affinité + Mémoire. Ag T-dép (protéines) vs T-indép (polysaccharides). Vaccins conjugués = solution pour nourrissons.
    \item \textbf{Régulation :} Treg (IL-10, TGF-$\beta$), Checkpoints (CTLA-4, PD-1) = Tolérance / Arrêt réponse.
    \item \textbf{Vaccination (Immunité Active) :} Mimer infection sans maladie. Types : Vivants atténués, Inactivés, Sous-unitaires/Conjugués/Recombinants, Toxoïdes. Adjuvant = Signal danger (Innée $\rightarrow$ Adaptative). Calendrier PEV = Stratégie santé publique.
    \item \textbf{Sérothérapie (Immunité Passive) :} Ac préformés (Polyclonaux : SAV, IGAT ; Monoclonaux : thérapeutiques). Immédiat, court ($\sim$21j), pas mémoire. Urgences / Immunodéprimés.
    \item \textbf{Immunité Collective :} Couverture $> P_c = 1-1/R_0$ $\rightarrow$ Protection non vaccinés. Objectifs OMS/PEV.
    \item \textbf{Hygiène :} Lavage mains (savon = détruit enveloppe/décolle), Désinfection plaie (propre $\rightarrow$ antiseptique), Hygiène respiratoire.
\end{itemize}
\end{aretenir}

\coursec{Activités d'intégration : Synthèse et Résolution de problèmes}

Après avoir détaillé les mécanismes de la coopération immune et les applications thérapeutiques que sont la vaccination et la sérothérapie, nous devons maintenant mobiliser l'ensemble de ces acquis pour analyser des situations concrètes, interpréter des données expérimentales et cliniques, et résoudre des problèmes biologiques. C'est l'étape cruciale où l'élève devient un biologiste en herbe : il observe, hypothétise, argumente et conclut.

\courssub{Cas clinique intégrateur : Prise en charge d'une plaie souillée chez un enfant non vacciné (Risque tétanique)}

\begin{experience}[Cas clinique : Enfant de 7 ans, plaie plantaire souillée par de la terre]
\textbf{Anamnèse :} Aminou, 7 ans, marche pieds nus dans la cour. Il se blesse au pied droit sur un débris métallique rouillé enfoui dans la terre. La plaie est profonde, hémorragique, avec des corps étrangers (terre, sable). Le carnet de vaccination est introuvable ; la mère pense qu'il n'a \textbf{jamais été vacciné} contre le tétanos.

\textbf{Examen clinique :} Plaie sale, nécrotique au fond. Pas de signes de contractures musculaires ni de trismus (verrouillage de la mâchoire) pour l'instant. Température : 37,8 °C.

\textbf{Question :} Quelle est la conduite à tenir immédiate au centre de santé ? Justifier chaque geste thérapeutique.
\end{experience}

\begin{methode}[Résolution d'un cas clinique en immunologie]
1. \textbf{Identifier l'agent pathogène et son mode d'action :} \textit{Clostridium tetani} (bacille anaérobie, sporulé, ubiquitinaire dans la terre). Toxine : tétanospasmine (neurotoxine bloquant la libération de GABA/glycine $\rightarrow$ contractures spasmodiques mortelles).
2. \textbf{Analyser le statut immunitaire du patient :} Non vacciné $\rightarrow$ \textbf{aucune mémoire immune}, aucun anticorps anti-toxine (IgG) circulants. Risque majeur de tétanos.
3. \textbf{Déterminer les deux objectifs thérapeutiques urgents :}
   \begin{itemize}
       \item \textbf{Neutraliser la toxine déjà présente ou en cours de production} (urgence immédiate : heures/jours).
       \item \textbf{Induire une protection active durable} pour l'avenir (prévention à long terme : mois/années).
   \end{itemize}
4. \textbf{Prescrire les traitements adaptés :}
   \begin{enumerate}
       \item \textbf{Nettoyage chirurgical + Métronidazole (antibiothérapie) :} Éliminer le nichoir anaérobie et la bactérie (source de toxine). \textit{Ne neutralise pas la toxine déjà fixée.}
       \item \textbf{Sérum antitetanique (SAT) / Immunoglobulines anti-tétaniques (IgAT) :} \cle{Immunité passive artificielle}. Apport immédiat d'anticorps polyclonaux (IgG) humains ou équins dirigés contre la toxine. \textbf{Neutralisation immédiate} de la toxine libre. Demi-vie courte (~3 semaines). \textbf{Site d'injection :} Voie intramusculaire (fesse ou cuisse), \textbf{site différent} du vaccin.
       \item \textbf{Vaccin anatoxine tétanique (VAT) :} \cle{Immunité active artificielle}. Anatoxine (toxine inactivée, immunogène conservée). Déclenche réponse primaire lente (pic IgG à 2-4 semaines) $\rightarrow$ mémoire immune (lymphocytes B mémoire, plasmocytes de longue vie). \textbf{Site d'injection :} Autre membre (ex: deltoïde si SAT en fesse).
   \end{enumerate}
5. \textbf{Suivi :} Rappel vaccinal à 1 mois, 6 mois (protocole rattrapage OMS). Surveillance clinique 3 semaines (période d'incubation).
\end{methode}

\begin{attention}[Pièges à éviter au BEPC]
\begin{itemize}
    \item \textbf{Confondre Sérothérapie et Vaccination :} Le sérum donne des \textbf{Ac tout faits} (passif, immédiat, temporaire). Le vaccin donne l'\textbf{Ag} pour que le corps fabrique ses Ac (actif, différé, durable).
    \item \textbf{Injecter au même endroit :} Les Ac du sérum peuvent neutraliser l'anatoxine du vaccin $\rightarrow$ \textbf{échec de la vaccination}. Sites \textbf{impérativement distincts}.
    \item \textbf{Oublier le nettoyage :} L'antibiotique et le débridement traitent la \textbf{cause} (bactérie), le sérum traite la \textbf{conséquence} (toxine). Les deux sont indispensables.
\end{itemize}
\end{attention}

\begin{popfigure}
\begin{tikzpicture}[scale=0.85, every node/.style={font=\small}]
    % Patient
    \node[draw, rounded corners, fill=popcream, thick, align=center] (patient) at (0,0){\textbf{Aminou, 7 ans}\\non-vacciné\\Plaie plantaire souillée (Terre/Rouille)};
    % Diagnostic
    \node[draw, diamond, aspect=2, fill=poporangeL, thick, below=1.5cm of patient, align=center] (diag){Risque Tétanos (\textit{Cl. tetani})\\non Pas d'Ac anti-toxine circulants};
    % Two branches
    \node[draw, rounded corners, fill=popgreenL, thick, left=3cm of diag, yshift=-1.5cm, align=center] (curatif){\textbf{Traitement Curatif Immédiat}\\1. Débridement + Nettoyage\\2. Métronidazole (IV/PO)\\3. \textbf{SAT / IgAT (Immunité Passive)}\\   \textit{Ac anti-toxine tout faits}\\   \textit{Action : Minutes/Heures}\\   \textit{Site : Fesse G}};
    \node[draw, rounded corners, fill=popblueL, thick, right=3cm of diag, yshift=-1.5cm, align=center] (preventif){\textbf{Prévention Durable}\\\textbf{VAT (Anatoxine) - Immunité Active}\\   \textit{Ag inactivé}\\   \textit{Action : Semaines (Réponse primaire)}\\   \textit{Mémoire : Années (LB mémoire)}\\   \textit{Site : Deltoïde D}};
   \draw[->, thick, popdark] (patient) -- (diag);
   \draw[->, thick, popgreen] (diag) -| (curatif);
   \draw[->, thick, popblue] (diag) -| (preventif);
   % Key box
   \node[draw, rounded corners, fill=popgoldL, thick, below=3.5cm of diag, align=center, text width=12cm] (key) {\textbf{Pourquoi deux injections sites différents ?}\\Les Ac du sérum (SAT) reconnaîtraient l'anatoxine du vaccin (VAT) comme un antigène étranger et la neutraliseraient \textbf{avant} qu'elle ne stimule les lymphocytes B. $\rightarrow$ \textbf{Échec de la primovaccination}. Distance physique = Éviter la neutralisation locale.};
\end{tikzpicture}
\legende{Algorithme de prise en charge du risque tétanique chez un sujet non vacciné. Distinction claire entre l'urgence curative (immunité passive) et la prévention durable (immunité active).}
\end{popfigure}

\courssub{Interprétation de courbes sérologiques : Cinétique des IgM et IgG}

L'analyse des courbes de taux d'anticorps (sérologie) est l'outil majeur pour dater une infection, vérifier une immunité vaccinale ou diagnostiquer une réinfection.

\begin{popfigure}
\begin{tikzpicture}
\begin{axis}[
    width=\linewidth, height=9cm,
    xlabel={Temps (semaines)}, ylabel={Taux d'anticorps (UA/mL)},
    xmin=0, xmax=25, ymin=0, ymax=110,
    xtick={0,3,5,7,10,14,20,25}, ytick={0,20,40,60,80,100},
    legend style={at={(0.02,0.98)}, anchor=north west, fill=white, fill opacity=0.8, draw=popline},
    axis lines=left, clip=false,
    grid=major, grid style={popline, opacity=0.3},
]
% Primary Response IgM
\addplot[poppink, thick, mark=*, mark size=2, samples=100, domain=0:25] 
    ({x}, { (x<2) ? 0 : (x<7) ? 80*sin(90*(x-2)/5) : 80*exp(-(x-7)/6) });
\addlegendentry{IgM (Infection primaire)}
% Primary Response IgG
\addplot[popblue, thick, mark=square*, mark size=2, samples=100, domain=0:25] 
    ({x}, { (x<3) ? 0 : (x<10) ? 90*(1-exp(-(x-3)/3)) : 90*exp(-(x-10)/20) + 15 });
\addlegendentry{IgG (Infection primaire)}
% Secondary Response (Booster/Vaccination) - shifted for clarity or separate axis? Let's do a second exposure at week 20.
% Actually, better to show two separate scenarios or a combined one. Let's do a combined graph with annotations.
% Secondary IgG (Booster at week 20)
\addplot[popgreen, very thick, dashed, mark=triangle*, mark size=2, samples=100, domain=20:25] 
    ({x}, { 15 + 100*(1-exp(-(x-20)/1.5)) });
\addlegendentry{IgG (Réponse secondaire / Rappel)}
% Annotations
\node[align=center, poppink, font=\small] at (axis cs: 5, 90) {Pic IgM\\\textbf{Phase aiguë}};
\node[align=center, popblue, font=\small] at (axis cs: 12, 95) {Pic IgG\\\textbf{Convalescence}};
\node[align=center, popgreen, font=\small] at (axis cs: 22, 105) {Montée rapide\\\textbf{Réponse anamnestique}};
\node[align=center, popmuted, font=\small, rotate=90] at (axis cs: 20, 50) {\textbf{Rappel Vaccinal / Réinfection}};
\draw[popmuted, dashed, thick] (axis cs: 20, 0) -- (axis cs: 20, 110);
\end{axis}
\end{tikzpicture}
\legende{Cinétique type des IgM et IgG lors d'une primo-infection (réponse primaire) et lors d'une ré-exposition (réponse secondaire / rappel vaccinal). L'IgM est le marqueur de la phase aiguë ; l'IgG persiste et monte vite et fort lors du second contact (mémoire).}
\end{popfigure}

\begin{exoresolu}[Interprétation d'un bilan sérologique]
\textbf{Énoncé :} Trois patients consultent pour une fièvre éruptive suspecte de rougeole. Le bilan sérologique (IgM / IgG anti-rougeole) donne :
\begin{center}
\begin{tabular}{|c|c|c|}\hline
\textbf{Patient} & \textbf{IgM} & \textbf{IgG} \\ \hline
A & \textbf{+++} (Positif fort) & + (Faible) \\ \hline
B & - (Négatif) & \textbf{+++} (Positif fort) \\ \hline
C & + (Faible) & \textbf{+++} (Positif fort, taux en hausse x4 en 10 j) \\ \hline
\end{tabular}
\end{center}
\textit{Interpréter le statut immunologique de chaque patient.}

\textbf{Solution détaillée :}
\begin{itemize}
    \item \textbf{Patient A :} \cle{IgM ++++ / IgG + faible} $\rightarrow$ \textbf{Infection aiguë / Phase invasive (Primaire)}. L'IgM apparaît en premier (J+3 à J+7). L'IgG commence tout juste à monter. \textit{Conduite : Isolement, traitement symptomatique, notification.}
    \item \textbf{Patient B :} \cle{IgM - / IgG +++} $\rightarrow$ \textbf{Immunité acquise (Convalescence ancienne ou Vaccination réussie)}. Pas d'infection en cours. L'IgM a disparu (demi-vie courte ~5 j), l'IgG persiste (mémoire). \textit{Conduite : Rassurer, pas de risque.}
    \item \textbf{Patient C :} \cle{IgM + faible / IgG +++ en hausse rapide (x4)} $\rightarrow$ \textbf{Réinfection ou Échec vaccinal (Réponse secondaire / Anamnestique)}. La présence d'IgM faible peut survenir en réponse secondaire. Le signe majeur est l'\textbf{élévation significative (x4) du taux d'IgG en 10-14 jours} (séroconversion ou montée anamnestique). Preuve d'une stimulation immune récente. \textit{Conduite : Surveillance, patient peu contagieux (immunité partielle), forme souvent atténuée.}
\end{itemize}
\end{exoresolu}

\begin{aretenir}[Lecture rapide d'une sérologie (Règle des 3 « I »)]
\begin{center}
\begin{tabular}{|l|c|c|c|}\hline
\textbf{Situation} & \textbf{IgM} & \textbf{IgG} & \textbf{Interprétation} \\ \hline
Infection aiguë (Primaire) & \textbf{+++} & + / ++ & \textbf{I}ncubation terminée, \textbf{I}nvasion \\
Immunité (Guéri / Vacciné) & \textbf{-} & \textbf{+++} & \textbf{I}mmunité protectrice \\
Réinfection / Rappel & + / - & \textbf{+++} \uparrow\uparrow & \textbf{I}mmunité stimulée (Anamnèse) \\ \hline
\end{tabular}
\end{center}
\textbf{Séroconversion :} Passage de IgG négatif à positif (ou multiplication par 4 du titre entre 2 prélèvements espacés de 2-3 semaines) = Preuve d'infection récente.
\end{aretenir}

\courssub{Exploitation documentaires : Preuves expérimentales du rôle des lymphocytes}

Le programme exige de savoir exploiter des documents d'expériences historiques fondatrices (transfert adoptif, irradiation/greffe) pour \textbf{identifier le rôle des lymphocytes} dans l'immunité spécifique.

\begin{experience}[Expérience 1 : Transfert adoptif chez la souris (Mitchison, 1954)]
\textbf{Protocole :}
\begin{enumerate}
    \item Souris \textbf{donneuse} : Immunisée contre l'antigène X (ex: hématies mouton) $\rightarrow$ Rate riche en lymphocytes activés/mémoire.
    \item Souris \textbf{réceptrice} : \textbf{Irradiée (léthale)} $\rightarrow$ Destruction de sa moelle osseuse $\rightarrow$ \textbf{Aplasie médullaire} (plus de lymphocytes, plus de globules blancs/rouges). Elle meurt en 10-15 j sans greffe.
    \item \textbf{Transfert :} Injection IV de cellules de rate de la donneuse immunisée dans la réceptrice irradiée.
    \item \textbf{Test :} Quelques jours plus tard, injection de l'antigène X dans la patte de la réceptrice.
\end{enumerate}
\textbf{Résultats :}
\begin{itemize}
    \item Groupe Témoin (Réceptrice irradiée \textbf{sans} transfert + Ag X) : \textbf{Pas de réponse} (pas de lymphocytes). Mort par aplasie.
    \item Groupe Transfert (Réceptrice irradiée \textbf{avec} cellules rate immunisée + Ag X) : \textbf{Réponse immune forte et rapide} (anticorps, mémoire). La souris survit (greffe de moelle par les cellules souches hématopoïétiques contenues dans la rate).
\end{itemize}
\textbf{Conclusion :} Les \textbf{lymphoïdes de la rate (lymphocytes)} portent la \textbf{spécificité} et la \textbf{mémoire} immunitaire. Ils sont \textbf{nécessaires et suffisants} pour transférer l'immunité spécifique.
\end{experience}

\begin{popfigure}
\begin{tikzpicture}[scale=0.8, every node/.style={font=\small}]
    % Donor
    \node[draw, rounded corners, fill=popblueL, thick, align=center] (donor) at (0, 3){\textbf{Souris Donneuse}\\Immunisée Ag X\\(Rate : Lymphocytes B/T mémoire)};
    % Recipient
    \node[draw, rounded corners, fill=poppinkL, thick, align=center] (recip) at (0, -3){\textbf{Souris Réceptrice}\\Irradiée (Létale)\\Aplasie médullaire\\(0 Lymphocyte)};
    % Transfer
    \node[draw, rounded corners, fill=popgoldL, thick, right=4cm of donor, yshift=-1.5cm, align=center] (transfer){\textbf{Transfert Adoptif}\\Prélèvement cellules rate\\Injection IV (queue)};
    % Challenge
    \node[draw, rounded corners, fill=popgreenL, thick, right=4cm of transfer, align=center] (challenge){\textbf{Test (Challenge)}\\Injection Ag X (patte)\\$\rightarrow$ Réponse immune forte\\$\rightarrow$ Survie + Immunité};
    % Control
    \node[draw, rounded corners, fill=poporangeL, thick, left=4cm of transfer, align=center] (control){\textbf{Témoin Négatif}\\Irradiée, Pas de transfert\\+ Ag X\\$\rightarrow$ Pas de réponse\\$\rightarrow$ Mort};
    \draw[->, thick, popblue] (donor) -- ++(1,0) |- (transfer);
    \draw[->, thick, poppink] (recip) -- ++(1,0) |- (transfer);
    \draw[->, thick, popdark] (transfer) -- (challenge);
    \draw[->, thick, popdark] (transfer) -- (control);
    % Conclusion box
    \node[draw, rounded corners, fill=popcream, thick, below=2.5cm of transfer, text width=12cm, align=center] (concl) {\textbf{Preuve apportée :} La compétence immune (spécificité + mémoire) réside dans les \textbf{lymphocytes} (cellules de la rate). La moelle osseuse (cellules souches) assure la survie (hématopoïèse).};
\end{tikzpicture}
\legende{Schéma de l'expérience de transfert adoptif (Mitchison). Démonstration que les lymphocytes sont les cellules effectrices de l'immunité spécifique et de la mémoire.}
\end{popfigure}

\begin{experience}[Expérience 2 : Irradiation / Greffe de moelle osseuse (Épreuve de la moelle)]
\textbf{Protocole :} Souris irradiée (léthale) $\rightarrow$ Greffe de moelle osseuse.
\begin{itemize}
    \item \textbf{Groupe A :} Moelle \textbf{syngénique} (même souche, compatible MHC) $\rightarrow$ \textbf{Survie}. Reconstitution hématopoïétique + immunitaire complète.
    \item \textbf{Groupe B :} Moelle \textbf{allogénique} (souche différente, MHC différent) $\rightarrow$ \textbf{Mort par Maladie du Greffon contre l'Hôte (GVH)}. Lymphocytes T donneurs (matures dans la greffe) attaquent les tissus du receveur (reconnu comme non-soi).
    \item \textbf{Groupe C :} Moelle allogénique \textbf{épurée en lymphocytes T} (déplétion CD3+) $\rightarrow$ \textbf{Survie} (reconstitution hématopoïétique) mais \textbf{Immunodépression T prolongée} (pas de LT matures donneurs).
\end{itemize}
\textbf{Conclusion :} La moelle contient les \textbf{cellules souches hématopoïétiques} (survie hématologique) ET des \textbf{lymphocytes T matures} (responsables de la GVH). L'immunité T nécessite une éducation thymique (sélection MHC) qui ne se fait pas instantanément après greffe.
\end{experience}

\begin{attention}[Analyse documentaire : Mots-clés attendus au BEPC]
\begin{itemize}
    \item \textbf{Irradiation} = Destruction moelle = \textbf{Aplasie} (disparition lymphocytes, GR, GB, Plaquettes).
    \item \textbf{Transfert cellules rate immunisée} = Transfert \textbf{spécificité + mémoire}.
    \item \textbf{Greffe moelle} = Transfert \textbf{cellules souches} (reconstitution lignées) + risque \textbf{GVH} (si LT matures donneurs $\neq$ MHC receveur).
    \item \textbf{Souris "Nude" (nu/nu)} = Mutante sans thymus $\rightarrow$ \textbf{Pas de LT} $\rightarrow$ Pas d'immunité cellulaire, pas d'aide aux LB $\rightarrow$ Pas d'IgG (seulement IgM faible). Preuve du rôle central du \textbf{Thymus} et des \textbf{LT CD4+}.
\end{itemize}
\end{attention}

\courssub{Rédaction argumentée : Pourquoi la vaccination protège-t-elle les non-vaccinés ? (Immunité collective)}

C'est une question de synthèse majeure, mêlant immunologie, épidémiologie et citoyenneté.

\begin{methode}[Rédiger une argumentation scientifique structurée]
1. \textbf{Définir le concept :} \cle{Immunité collective (ou de groupe)} : Protection indirecte des individus non immunisés (non vaccinés, immunodéprimés, nourrissons trop jeunes) due à la forte proportion d'individus immunisés dans la population.
2. \textbf{Expliquer le mécanisme (Chaîne de transmission) :}
   \begin{itemize}
       \item Agent pathogène transmis de personne à personne (contagiosité).
       \item Sujet vacciné $\rightarrow$ \textbf{Blocage de la réplication} (immunité stérilisante) ou \textbf{réduction drastique de la charge virale/bactérienne et de la durée d'excrétion}.
       \item $\rightarrow$ \textbf{Rupture des chaînes de transmission}.
   \end{itemize}
3. \textbf{Introduire le paramètre clé : $R_0$ (Taux de reproduction de base).}
   \begin{itemize}
       \item $R_0$ = Nombre moyen de cas secondaires générés par 1 cas index dans une population \textbf{entièrement sensible}.
       \item Si couverture vaccinale $p$ (efficacité vaccin $e \approx 1$), $R_{eff} = R_0 (1 - p)$.
       \item Seuil d'élimination : $R_{eff} < 1 \Rightarrow p > 1 - \frac{1}{R_0}$.
   \end{itemize}
4. \textbf{Illustrer par des exemples chiffrés (Cameroun / Monde).}
5. \textbf{Conclure :} Vaccination = Acte de solidarité. Protège le "moi" et le "nous".
\end{methode}

\begin{exemplebox}[Calcul du seuil d'immunité collective (Exemple chiffré)]
\textbf{Formule :} $p_c = 1 - \frac{1}{R_0}$ (en supposant efficacité vaccinale = 100\%).
\begin{center}
\begin{tabular}{|l|c|c|c|}\hline
\textbf{Maladie} & \textbf{$R_0$ estimé} & \textbf{Seuil critique $p_c$} & \textbf{Couverture visée OMS} \\ \hline
Rougeole & 12 -- 18 (\textbf{très contagieux}) & 92\% -- 94\% & \textbf{95\%} (2 doses) \\ \hline
Coqueluche & 12 -- 17 & 92\% -- 94\% & 90\% -- 95\% \\ \hline
Diphtérie & 6 -- 7 & 83\% -- 86\% & 90\% \\ \hline
Polio & 5 -- 7 & 80\% -- 86\% & 90\% \\ \hline
COVID-19 (souche initiale) & 2.5 -- 3.5 & 60\% -- 71\% & Variable \\ \hline
\end{tabular}
\end{center}
\textbf{Application au Cameroun :} Pour la rougeole ($R_0 \approx 15$), il faut vacciner \textbf{au moins 93\%} des enfants pour arrêter la circulation du virus. Si la couverture tombe à 80\% (zones d'accès difficile, hésitation vaccinale), $R_{eff} = 15 \times 0.2 = 3 > 1$ : \textbf{Épidémie inévitable}. C'est ce qu'on observe lors des flambées dans le Nord ou l'Extrême-Nord.
\end{exemplebox}

\begin{savaistu}[Immunité collective et "Free-riders" (Profiteurs gratuits)]
Les personnes qui refusent la vaccination par conviction personnelle (mais pas médicale) bénéficient de la protection du groupe sans en assumer le risque (faible) ni le coût. Si leur nombre dépasse le seuil critique ($1-p_c$), l'immunité collective s'effondre : ce sont les plus vulnérables (nourrissons < 9 mois, enfants cancéreux, VIH+) qui paient le prix fort. \textbf{La vaccination obligatoire (BCG, DTC, Rougeole, Fièvre jaune au Cameroun) vise à garantir ce seuil.}
\end{savaistu}

\courssub{Résolution de problèmes : Déficits immunitaires primitifs et acquis}

Deux pathologies modèles illustrent la dissociation entre immunité humorale (LB/Ac) et immunité cellulaire (LT).

\begin{exoresolu}[Cas 1 : Agammaglobulinémie de Bruton (Déficit humorale pur)]
\textbf{Énoncé :} Nourrisson de 6 mois, garçon. Infections récidivantes : otites, sinusites, pneumonies à \textit{Streptococcus pneumoniae}, \textit{Haemophilus influenzae} (bactéries \textbf{pyogènes encapsulées}). Pas d'adénopathies, amygdales invisibles. Bilan : \textbf{IgG, IgA, IgM indosables}. Lymphocytes T normaux. Génétique : Mutation gène \textit{BTK} (Tyrosine Kinase de Bruton) sur chromosome X.

\textbf{Analyse et Résolution :}
\begin{enumerate}
    \item \textbf{Identification du défaut :} Blocage maturation LB au stade \textbf{pré-B} (BTK nécessaire pour signalisation récepteur pré-BCR). $\rightarrow$ \textbf{Absence de LB matures} $\rightarrow$ \textbf{Absence d'Ac (Agammaglobulinémie)}.
    \item \textbf{Conséquence fonctionnelle :} Pas d'opsonisation (IgG), pas de neutralisation toxines/virus, pas d'activation complément classique. \textbf{Défaillance majeure contre bactéries encapsulées} (capsule polysaccharides = Ag TI-2 nécessitant LB matures + Ac pour opsonisation).
    \item \textbf{Immunité cellulaire intacte :} LT normaux $\rightarrow$ Défense vs virus intracellulaires, champignons, tuberculose \textbf{préservée}.
    \item \textbf{Prise en charge (Standard de soins) :}
        \begin{itemize}
            \item \textbf{Immunoglobulines IV (IgIV) ou SC (sous-cutanées) :} Apport passif d'IgG polyclonales toutes les 3-4 semaines (demi-vie ~3-4 sem). \textbf{Traitement à vie}. Cible : Taux IgG résiduel > 5-7 g/L.
            \item Vaccins \textbf{vivants atténués CONTRE-INDICQUÉS} (risque maladie vaccinale non contrôlée par Ac).
            \item Antibioprophylaxie (Cotrimoxazole) souvent associée.
        \end{itemize}
\end{enumerate}
\end{exoresolu}

\begin{exoresolu}[Cas 2 : SIDA (Déficit acquis mixte : Cellulaire + Humorale secondaire)]
\textbf{Énoncé :} Jeune adulte, fièvre prolongée, amaigrissement, diarrhée chronique. Bilan : \textbf{Lymphopénie T CD4+ < 200 /µL} (Norme > 500). Charge virale VIH-1 élevée. Infections opportunistes : \textit{Pneumocystis jirovecii} (pneumonie), \textit{Toxoplasma} (encéphalite), \textit{Candida} (œsophagite), Herpès/Zona récidivants. Sérologie : Hypergammaglobulinémie polyclonale (B activés polyclonalement par VIH) mais \textbf{Ac spécifiques inefficaces} (pas d'aide LT CD4+).

\textbf{Analyse et Résolution :}
\begin{enumerate}
    \item \textbf{Cible du VIH :} Récepteur CD4 + Co-récepteurs (CCR5/CXCR4) sur \textbf{LT CD4+ (Th)}, Macrophages, Cellules dendritiques.
    \item \textbf{Mécanisme de destruction :} Infection productive $\rightarrow$ Lyse cellulaire + Apoptose "bystander" + Épuisement thymique. $\rightarrow$ \textbf{Effondrement LT CD4+}.
    \item \textbf{Conséquences systémiques :}
        \begin{itemize}
            \item \textbf{Immunité cellulaire (Th1) :} Pas d'IL-2, IFN-$\gamma$ $\rightarrow$ Pas d'activation Macrophages (Tuéris intracellulaires) + Pas de prolifération LT CD8+ cytotoxiques. $\rightarrow$ \textbf{Infections opportunistes intracellulaires (Mycobactéries, Champignons, Virus, Protozoaires)}.
            \item \textbf{Immunité humorale (Th2/Tfh) :} Pas d'aide (IL-4, IL-21, CD40L) $\rightarrow$ Pas de commutation de classe (IgG/IgA), pas d'affinité maturation, pas de LB mémoire. $\rightarrow$ \textbf{Ac de "mauvaise qualité" (faible affinité, pas de mémoire)} malgré hypergammaglobulinémie.
        \end{itemize}
    \item \textbf{Prise en charge :}
        \begin{itemize}
            \item \textbf{ARV (Antirétroviraux) :} Triple thérapie (2 INTI + 1 INNTI ou IP). Objectif : \textbf{Charge virale indétectable (< 50 cp/mL)} $\rightarrow$ Remontée CD4+ (reconstitution immune).
            \item \textbf{Prophylaxies primaires (si CD4 < seuils) :} Cotrimoxazole (Pneumocystis/Toxo), Fluconazole (Cryptocoque), Isoniazide (TB).
            \item \textbf{Vaccination :} Vivants atténués contre-indiqués si CD4 < 200. Inactivés recommandés (Grippe, Pneumo, Hep B, HPV, Tétanos, Diphtérie, Coqueluche). Réponse vaccinale sous-optimale si CD4 bas $\rightarrow$ Revacciner après reconstitution immune (CD4 > 200-500).
        \end{itemize}
\end{enumerate}
\end{exoresolu}

\begin{aretenir}[Tableau comparatif : Bruton vs SIDA (Synthèse BEPC)]
\begin{center}
\begin{tabularx}{\linewidth}{|l|X|X|}\hline
\textbf{Critère} & \textbf{Agammaglobulinémie de Bruton (XLA)} & \textbf{SIDA (VIH)} \\ \hline
\textbf{Origine} & Génétique (Récessive liée à X) & Acquise (Infection rétrovirale) \\ \hline
\textbf{Défaut principal} & \textbf{Blocage maturation LB} (Pré-B) & \textbf{Destruction LT CD4+} (Th) \\ \hline
\textbf{Immunité Humorale} & \textbf{Absente} (Agammaglobulinémie) & \textbf{Défaillante qualitative} (Pas d'aide Th) \\ \hline
\textbf{Immunité Cellulaire} & \textbf{Normale} & \textbf{Effondrée} (CD4+ bas) \\ \hline
\textbf{Infections types} & Bactéries pyogènes encapsulées (Extracellulaires) & Agents opportunistes intracellulaires + Virus + Champignons \\ \hline
\textbf{Traitement étiologique} & Aucun (Génique) & \textbf{ARV (Suppression réplication VIH)} \\ \hline
\textbf{Traitement substitutif} & \textbf{IgIV / IgSC à vie} & Non (sauf cas rares d'hypogammaglobulinémie sévère) \\ \hline
\end{tabularx}
\end{center}
\end{aretenir}

\courssub{Activité de synthèse : Carte conceptuelle à compléter}

Pour valider l'unité, vous devez être capable de structurer vos connaissances. Complétez la carte ci-dessous (version élève : cases vides) en utilisant les mots-clés du banc.

\begin{popfigure}
\begin{center}\tcbox[colback=popdark,colframe=popdark,coltext=white,arc=9pt,boxsep=4pt,left=6pt,right=6pt]{\bfseries\small RÉPONSE IMMUNITAIRE}\end{center}
\vspace{-2pt}
\begin{tcbraster}[raster columns=2,raster equal height=rows,raster column skip=6pt,raster row skip=6pt,enhanced,blankest]
\begin{tcolorbox}[enhanced,colback=white,colframe=popgreen,boxrule=0.9pt,arc=3pt,title={INNÉE (Non-spécifique)},fonttitle=\bfseries\scriptsize,coltitle=white,colbacktitle=popgreen,left=4pt,right=4pt,top=2pt,bottom=2pt,boxsep=1pt,toptitle=1.5pt,bottomtitle=1.5pt,halign=left]
\begin{itemize}[nosep,leftmargin=*,itemsep=1pt,font=\scriptsize]
\item Barrières 1ère ligne
\begin{itemize}[nosep,leftmargin=1.1em,itemsep=0pt]
\item Peau : Kératine, pH, Flore, Sébum
\item Muqueuses : Cils, Mucus, Lysozyme, pH, Flore
\end{itemize}
\item Barrières 2ème ligne
\begin{itemize}[nosep,leftmargin=1.1em,itemsep=0pt]
\item Inflammation : Rubor, Calor, Tumor, Dolor, Functio laesa
\item Phagocytose : Polynucléaires (Neutrophiles), Macrophages
\item Molécules : Complément, Interféron $\alpha/\beta$, Protéines phase aiguë (CRP)
\item Cellules : NK (Natural Killer), Mastocytes (Histamine)
\end{itemize}
\end{itemize}
\end{tcolorbox}
\begin{tcolorbox}[enhanced,colback=white,colframe=poppurple,boxrule=0.9pt,arc=3pt,title={ADAPTATIVE (Spécifique)},fonttitle=\bfseries\scriptsize,coltitle=white,colbacktitle=poppurple,left=4pt,right=4pt,top=2pt,bottom=2pt,boxsep=1pt,toptitle=1.5pt,bottomtitle=1.5pt,halign=left]
\begin{itemize}[nosep,leftmargin=*,itemsep=1pt,font=\scriptsize]
\item Lymphocytes B (Humoral)
\begin{itemize}[nosep,leftmargin=1.1em,itemsep=0pt]
\item Reconnaissance : BCR (IgM/IgD membranaire) = Ag NATIF
\item Activation : Ag + Aide LT CD4+ $\rightarrow$ Centre germinatif
\item Différenciation : Plasmocytes (Usine à Ac) + LB Mémoire
\item Effecteurs : Ac (IgG, IgA, IgM, IgE) : Neutralisation, Opsonisation, ADCC, Complément
\end{itemize}
\item Lymphocytes T (Cellulaire)
\begin{itemize}[nosep,leftmargin=1.1em,itemsep=0pt]
\item LT CD8+ Cytotoxiques : Reconnaissance Ag/MHC I $\rightarrow$ Lyse cellulaire
\item LT CD4+ Helpers : Reconnaissance Ag/MHC II $\rightarrow$ Cytokines (Aide LB, Activation Macrophages, LT CD8+)
\end{itemize}
\end{itemize}
\end{tcolorbox}
\begin{tcolorbox}[enhanced,colback=white,colframe=popblue,boxrule=0.9pt,arc=3pt,title={CARACTÉRISTIQUES COMPARÉES},fonttitle=\bfseries\scriptsize,coltitle=white,colbacktitle=popblue,left=4pt,right=4pt,top=2pt,bottom=2pt,boxsep=1pt,toptitle=1.5pt,bottomtitle=1.5pt,halign=left]
\begin{itemize}[nosep,leftmargin=*,itemsep=1pt,font=\scriptsize]
\item Délai : Innée = Minutes/Heures | Adaptative = Jours (Primaire) / Heures (Secondaire)
\item Spécificité : Innée = Large (PAMPs) | Adaptative = Fine (Épitope unique)
\item Mémoire : Innée = Faible (Entraînement) | Adaptative = Forte, Durable, Spécifique
\item Récepteurs : Innée = PRR (germline) | Adaptative = BCR/TCR (Réarrangés V(D)J)
\end{itemize}
\end{tcolorbox}
\end{tcbraster}

\legende{Carte conceptuelle synthétique de l'unité "Réponse immunitaire". À l'élève de compléter les cases vides (version élève fournie séparément) pour réviser les articulations entre composantes.}
\end{popfigure}

\begin{exercice}[Entraînement BEPC : Sujet de synthèse]
\textbf{Document 1 :} Courbe de survie de souris irradiées (8 Gy) greffées ou non avec de la moelle osseuse syngénique.
\textbf{Document 2 :} Taux d'IgG anti-tétanique chez un adulte vacciné il y a 10 ans, qui reçoit un rappel (J0) et un enfant non vacciné qui reçoit sa 1ère dose (J0).
\textbf{Document 3 :} Arbre généalogique d'une famille touchée par l'agammaglobulinémie de Bruton (maladie récessive liée à l'X).

\textbf{Questions :}
\begin{enumerate}
    \item À partir du Doc 1, expliquer le rôle de la moelle osseuse dans la reconstitution immunitaire. Pourquoi la greffe doit-elle être syngénique ?
    \item À partir du Doc 2, comparer la cinétique et l'amplitude de la réponse IgG chez l'adulte et l'enfant. Identifier le type de réponse (primaire/secondaire) et les cellules responsables de la différence.
    \item À partir du Doc 3, déterminer le mode de transmission. Calculer le risque pour un garçon né d'une mère porteuse. Justifier pourquoi les filles sont rarement malades.
    \item \textbf{Rédaction :} "En vous appuyant sur l'ensemble des documents et vos connaissances, expliquez pourquoi la vaccination est la stratégie la plus efficace et la plus économique pour contrôler une maladie infectieuse contagieuse dans une population."
\end{enumerate}
\end{exercice}

\begin{corrige}[Exercice n°1 - Éléments de corrigé type BEPC]
\begin{enumerate}
    \item \textbf{Doc 1 :} Souris irradiée non greffée $\rightarrow$ Mort (aplasie). Souris greffée moelle syngénique $\rightarrow$ Survie. \textbf{Rôle :} La moelle contient les \textbf{cellules souches hématopoïétiques} qui reconstruisent toutes les lignées (GR, GB, Plaquettes, Lymphocytes). \textbf{Syngénique :} Même MHC (H-2 chez la souris / HLA chez l'homme). Évite le rejet (LT CD8+ receveur vs greffon) et la \textbf{GVH} (LT CD4+/CD8+ donneur vs receveur).
    \item \textbf{Doc 2 :} Adulte (Rappel) : IgG base élevée $\rightarrow$ Montée \textbf{rapide} (J3-5), \textbf{forte} (taux max élevé), \textbf{affinité forte}. = \textbf{Réponse secondaire (Anamnestique)}. Cellules : \textbf{LB Mémoire + Plasmocytes de longue vie}. Enfant (Primo) : IgG base 0 $\rightarrow$ Montée \textbf{lente} (J7-10), \textbf{faible} (pic bas), affinité faible. = \textbf{Réponse primaire}. Cellules : \textbf{LB Naïfs} $\rightarrow$ Différenciation.
    \item \textbf{Doc 3 :} Transmission \textbf{récessive liée à l'X} (Gène \textit{BTK} sur X). Mère porteuse ($X^B X^b$) x Père sain ($X^B Y$). Garçon : 50\% risque ($X^b Y$ = malade). Fille : 50\% porteuse ($X^B X^b$), 50\% saine ($X^B X^B$). Filles rarement malades car il faut 2 allèles mutés ($X^b X^b$) $\rightarrow$ Père malade (rare) + Mère porteuse.
    \item \textbf{Rédaction :} Vaccination $\rightarrow$ Immunité active (LB mémoire, Plasmocytes de longue vie, LT mémoire) $\rightarrow$ Protection individuelle (réponse secondaire rapide/efficace si infection). Couverture vaccinale $> p_c = 1-1/R_0$ $\rightarrow$ $R_{eff} < 1$ $\rightarrow$ \textbf{Immunité collective} (arrêt circulation agent). Protection des non-vaccinables (nourrissons, immunodéprimés). Coût vaccin $\ll$ Coût traitement + hospitalisation + perte productivité + vies. Exemple : Éradication variole, élimination polio (Cameroun certifié libre polio sauvage 2020), contrôle rougeole/diphtérie/tétanos via PEV.
\end{enumerate}
\end{corrige}

\begin{savaistu}[Le "Test de tuberculine" (IDR) : Preuve d'immunité cellulaire mémoire]
Ce n'est \textbf{PAS} un test de diagnostic de la tuberculose active (maladie). C'est un test d'\textbf{infection latente} (immunité cellulaire mémoire).
\begin{itemize}
    \item \textbf{IDR (Intradermo-réaction) :} Injection dermique de PPD (protéines purifiées de \textit{M. tuberculosis}). Lecture à 48-72h. Induration $\ge$ 10 mm (ou 5 mm si VIH/contact) = \textbf{Réaction d'hypersensibilité retardée (Type IV)}. Médiée par \textbf{LT CD4+ Th1 mémoire} (IFN-$\gamma$, TNF-$\alpha$) recrutant macrophages. Preuve de contact passé avec le bacille (ou BCG).
    \item \textbf{Interprétation :} Positif = Infection latente (bacille dormant dans granulome). Risque réactivation ~5-10\% vie entière (plus élevé si VIH). $\rightarrow$ Prophylaxie Isoniazide 6-9 mois.
\end{itemize}
\end{savaistu}

\begin{attention}[Derniers conseils pour l'épreuve du BEPC]
\begin{itemize}
    \item \textbf{Lisez TOUS les documents} avant d'écrire. Les réponses sont DANS les docs + vos connaissances.
    \item \textbf{Vocabulaire précis :} "Antigène" $\neq$ "Anticorps" ; "Lymphocyte B" $\neq$ "Lymphocyte T" ; "Immunité passive" $\neq$ "Immunité active" ; "Primaire" $\neq$ "Secondaire".
    \item \textbf{Justifiez :} Toute affirmation doit s'appuyer sur une donnée du document ("Le document 1 montre que...") ou une connaissance ("On sait que les LT CD4+...").
    \item \textbf{Unités :} Temps en jours/semaines, Concentration en g/L ou UI/mL, Numération en /µL ou /mm$^3$ ou $10^9$/L.
    \item \textbf{Schéma :} Crayon à papier, règle, légendes fléchées, titre informatif. Ne pas faire d'art abstrait.
\end{itemize}
\end{attention}

\newpage

\coursec{Activité d'intégration}

\courssub{Objectifs de l'activité}
\begin{itemize}
    \item Mobiliser l'ensemble des connaissances de la leçon (immunité innée, immunité adaptative, vaccination, sérologie) pour analyser un dossier clinique réel.
    \item Appliquer une démarche d'investigation scientifique : observer des données cliniques et biologiques, formuler des hypothèses, interpréter des résultats, conclure et communiquer.
    \item Rédiger une note de synthèse médicale structurée et argumentée, adaptée à une situation de santé publique au Cameroun.
    \item Développer des attitudes responsables face aux maladies à prévention vaccinale et aux épidémies.
\end{itemize}

\courssub{Situation : Enquête au Centre de Santé de District de Bertoua — Le mystère de la fièvre jaune}
\begin{experience}[Dossier patient n° 2026-0457 — Centre de Santé de District de Bertoua, Région de l'Est]
\textbf{Contexte épidémiologique :} Nous sommes en mars 2026, saison sèche chaude dans la région de l'Est du Cameroun. Le district de santé signale une augmentation inhabituelle de cas de fièvre ictérique hémorragique. Le Programme Élargi de Vaccination (PEV) rapporte une couverture vaccinale anti-amarile (fièvre jaune) de 68 \% dans le district (objectif $> 80 \%$), avec des poches de non-vaccination dans les zones riveraines et les sites de déplacés internes.

\textbf{1. Anamnèse (Interrogatoire médical) :}
\begin{itemize}
    \item \textbf{Patient :} M. K., 24 ans, cultivateur, résident au village Nkolbisson (zone forestière, à 15 km de Bertoua).
    \item \textbf{Motif de consultation :} Fièvre élevée (39,5 °C) depuis 3 jours, douleurs musculaires intenses, céphalées, vomissements.
    \item \textbf{Antécédents :} Pas de pathologie connue. \textbf{Statut vaccinal :} Carte de vaccination introuvable. Le patient déclare n'avoir « jamais reçu le vaccin de la fièvre jaune » par méfiance vis-à-vis des campagnes de vaccination et par difficulté d'accès au centre de santé (éloignement, coûts de transport indirects).
    \item \textbf{Exposition :} Travaux champêtres quotidiens en lisière de forêt ; présence de moustiques \emph{Aedes} signalée. Aucun voyage récent hors du district.
\end{itemize}

\textbf{2. Examen clinique (J0 — Jour de l'admission) :}
\begin{itemize}
    \item État général altéré, prostration. Température : 39,8 °C. Fréquence cardiaque : 110 bpm. TA : 100/60 mmHg.
    \item \textbf{Ictère franc :} téguments et conjonctives jaunes.
    \item \textbf{Signes hémorragiques :} pétéchies aux membres inférieurs, gingivorragie à la palpation, hématémèse (vomissement de sang) notée ce matin.
    \item Hépatomégalie douloureuse (foie palpable 4 cm sous le rebord costal).
    \item Pas de méningisme, pas de trouble de la conscience.
\end{itemize}

\textbf{3. Bilan biologique initial (J0) :}
\begin{center}
\begin{tabularx}{\linewidth}{|l|X|X|}\hline
\textbf{Paramètre} & \textbf{Résultat} & \textbf{Valeurs de référence (adulte)} \\ \hline
Hématies & 4,2 $\times$ 10$^{12}$/L & 4,5 -- 5,5 $\times$ 10$^{12}$/L \\ \hline
Hémoglobine & 11,5 g/dL & 13 -- 17 g/dL \\ \hline
Leucocytes (Globules blancs) & \textbf{2,1 $\times$ 10$^9$/L} & 4 -- 10 $\times$ 10$^9$/L \\ \hline
\rowcolor{poporangeL} \textbf{Leucopénie} & \textbf{OUI} & -- \\ \hline
Plaquettes (Thrombocytes) & \textbf{45 $\times$ 10$^9$/L} & 150 -- 400 $\times$ 10$^9$/L \\ \hline
\rowcolor{poporangeL} \textbf{Thrombopénie} & \textbf{OUI} & -- \\ \hline
ASAT (SGOT) & \textbf{1 850 UI/L} & $< 40$ UI/L \\ \hline
ALAT (SGPT) & \textbf{2 100 UI/L} & $< 40$ UI/L \\ \hline
Bilirubine totale & \textbf{45 mg/L} & $< 17$ mg/L \\ \hline
CRP (Protéine C-Réactive) & \textbf{120 mg/L} & $< 6$ mg/L \\ \hline
Créatinine & 14 mg/L & 6 -- 12 mg/L \\ \hline
Glycémie & 0,65 g/L & 0,70 -- 1,10 g/L \\ \hline
\end{tabularx}
\end{center}

\textbf{4. Sérologie Fièvre Jaune (Technique ELISA capture IgM / IgG) — Prélèvement J0 :}
\begin{center}
\begin{tabularx}{\linewidth}{|l|X|X|}\hline
\textbf{Marqueur} & \textbf{Résultat} & \textbf{Interprétation standard} \\ \hline
IgM anti-Virus Amaril & \textbf{Positif (++++)} & Infection récente (aiguë) \\ \hline
IgG anti-Virus Amaril & \textbf{Négatif} & Pas d'infection ancienne, pas d'immunité vaccinale \\ \hline
\end{tabularx}
\end{center}

\textbf{5. Cinétique théorique des anticorps anti-Fièvre Jaune (Modèle de référence PEV/OMS) :}
\begin{popfigure}
\begin{tikzpicture}[scale=0.75, font=\small]
\begin{axis}[
    width=\linewidth, height=8cm,
    xlabel={Jours après début des symptômes}, ylabel={Taux d'anticorps (UA/mL - Unités Arbitraires)},
    xmin=0, xmax=30, ymin=0, ymax=100,
    xtick={0,5,10,15,20,25,30}, ytick={0,20,40,60,80,100},
    legend pos=north west, grid=major, grid style={popline},
    axis line style={popdark}, tick style={popdark},
]
\addplot[popblue, thick, mark=*] coordinates {
    (0,0) (3,5) (5,25) (7,55) (10,80) (14,90) (21,85) (30,70)
};
\addlegendentry{IgM (Pic J10-J14, déclin lent)}
\addplot[poporange, thick, mark=square*] coordinates {
    (0,0) (5,0) (7,5) (10,20) (14,50) (21,80) (30,95)
};
\addlegendentry{IgG (Apparition J7-J10, plateau durable)}
\draw[dashed, popmuted] (7,0) -- (7,100) node[midway, right, rotate=90] {J7};
\draw[dashed, popmuted] (14,0) -- (14,100) node[midway, right, rotate=90] {J14};
\draw[popgreen, thick, ->] (3.5, 95) -- (3.5, 80) node[above, left] {Patient (J3-J4)};
\end{axis}
\end{tikzpicture}
\legende{Cinétique théorique de la réponse humorale anti-Fièvre Jaune. Le patient est consulté vers J3-J4 (flèche verte).}
\end{popfigure}

\textbf{6. Extrait du Calendrier PEV Cameroun 2026 (Vaccin Amaril) :}
\begin{itemize}
    \item \textbf{Cible :} Enfants à 9 mois (dose unique), Rattrapage 12-23 mois, Campagnes de rattrapage massives (9 mois -- 60 ans) tous les 3-5 ans dans les zones à risque.
    \item \textbf{Stratégie :} Vaccination de routine (fixe et avancée) + Campagnes préventives de masse + Vaccination réactive en cas d'épidémie.
    \item \textbf{Stock district (Mars 2026) :} Rupture de vaccin amaril depuis 3 semaines (retard acheminement régional). Dernière campagne de masse : 2022.
\end{itemize}
\end{experience}

\courssub{Consignes}
\begin{enumerate}
    \item \textbf{Analyse de l'immunité innée :} À partir des signes cliniques (fièvre, CRP, leucopénie, thrombopénie, signes hémorragiques) et des données biologiques, expliquez les mécanismes de la réponse inflammatoire et de la phagocytose mis en jeu. Pourquoi observe-t-on une leucopénie et une thrombopénie malgré l'inflammation ?
    \item \textbf{Interprétation sérologique :} En vous référant à la courbe cinétique, justifiez le profil sérologique observé (IgM+, IgG-) au jour de la consultation (J3-J4). Que nous apprend ce profil sur le statut immunitaire préexistant du patient (vaccination, infection antérieure) ?
    \item \textbf{Analyse de la défaillance vaccinale :} Identifiez, à partir du dossier, les causes individuelles et systémiques (organisationnelles) qui expliquent l'absence de protection vaccinale chez ce patient.
    \item \textbf{Prise en charge thérapeutique :} Proposez une conduite à tenir argumentée au lit du patient (traitement étiologique, symptomatique, surveillance). Justifiez l'absence de traitement antiviral spécifique.
    \item \textbf{Mesures de santé publique :} En tant qu'agent de surveillance épidémiologique, détaillez les mesures immédiates à mettre en œuvre au niveau du district (déclaration, investigation, riposte vaccinale, lutte vectorielle, communication).
    \item \textbf{Rédaction professionnelle :} Rédigez la \emph{Note de synthèse médicale} (format standard : Identification, Motif, Antécédents, Examen, Bilan, Discussion diagnostique, Conduite à tenir, Mesures de prévention) destinée au Médecin-Chef de District.
\end{enumerate}

\courssub{Ressources à mobiliser}
\begin{aretenir}[Rappels clés pour la résolution]
\textbf{Immunité innée (Non spécifique) :}
\begin{itemize}
    \item \textbf{Barrières :} Peau (kératine, pH acide, flore), Muqueuses (cilia, mucus, lysozyme, pH gastrique).
    \item \textbf{Inflammation :} Rubor, Calor, Tumor, Dolor, Functio laesa. Médiateurs : Histamine, Prostaglandines, Cytokines (IL-1, TNF-$\alpha$). Vasodilatation + Perméabilité vasculaire $\rightarrow$ Œdème, arrivée des leucocytes.
    \item \textbf{Phagocytose :} Neutrophiles (polynucléaires neutrophiles, premiers arrivés, vie courte), Macrophages (monocytes tissulaires, présentation d'antigène). Chimiotaxie $\rightarrow$ Adhérence $\rightarrow$ Ingestion $\rightarrow$ Digestion lysosomale.
    \item \textbf{Protéines de phase aiguë :} CRP (synthèse hépatique stimulée par IL-6), marqueur sensible et précoce de l'inflammation.
\end{itemize}
\textbf{Immunité adaptative (Spécifique) :}
\begin{itemize}
    \item \textbf{Lymphocytes B :} Différenciation en Plasmocytes (usines à Ac) et Lymphocytes B mémoire. Ac = Immunoglobulines (IgM précoce, pentamère, forte avidité ; IgG tardive, monomère, haute affinité, mémoire, passage barrière placentaire).
    \item \textbf{Lymphocytes T CD4+ (Auxiliaires/Th) :} Orchestration (cytokines). Th1 (immunité cellulaire), Th2 (immunité humorale), Tfh (aide aux centres germinatifs).
    \item \textbf{Lymphocytes T CD8+ (Cytotoxiques/Tc) :} Lyse des cellules infectées (reconnaissance MHC I + peptide).
    \item \textbf{Mémoire immunologique :} Réponse secondaire plus rapide, plus intense, plus durable (IgG).
\end{itemize}
\textbf{Vaccination / Sérothérapie :}
\begin{itemize}
    \item \textbf{Vaccin amaril (Fièvre Jaune) :} Virus vivant atténué (souche 17D). Une dose $\rightarrow$ immunité à vie (OMS 2014). Induit IgM puis IgG neutralisantes + mémoire B/T.
    \item \textbf{Sérothérapie :} Injection d'Ig polyvalentes ou spécifiques (immunité passive immédiate, temporaire, pas de mémoire). Non indiquée en routine pour Fièvre Jaune.
\end{itemize}
\textbf{Fièvre Jaune (Arbovirus, Flavivirus) :}
\begin{itemize}
    \item Transmission : Moustiques \emph{Aedes} (cycle sylvatique singe-moustique-homme ; cycle urbain homme-moustique-homme).
    \item Tropisme : Hépatocytes (ictère, cytolyse), Rénaux, Myocarde, Endothélium vasculaire (hémorragies, thrombopénie par consommation/destruction).
    \item Pas de traitement antiviral spécifique homologué (soins de support vitaux).
\end{itemize}
\end{aretenir}

\courssub{Démarche et corrigé commenté}
\begin{exoresolu}[Résolution guidée de l'enquête — Dossier M. K.]
\textbf{Énoncé :} Voir les 6 questions des Consignes ci-dessus.
\tcblower
\textbf{Question 1 : Analyse de l'immunité innée — Signes cliniques et biologiques de l'inflammation.}

\textit{Observations :} Fièvre (39,8 °C), CRP 120 mg/L (+++), Leucopénie (2,1 G/L), Thrombopénie (45 G/L), Signes hémorragiques, Cytolyse hépatique massive.

\textit{Interprétation :}
\begin{enumerate}
    \item \textbf{Déclenchement :} Le virus de la fièvre jaune infecte les hépatocytes, cellules de Kupffer (macrophages hépatiques) et l'endothélium vasculaire. La lyse cellulaire et les PAMPs (motifs moléculaires associés aux pathogènes, ex: ARN viral) activent les macrophages résidents et les cellules dendritiques.
    \item \textbf{Médiateurs de l'inflammation :} Libération massive d'\textbf{histamine} (mastocytes), de \textbf{prostaglandines}, de \textbf{cytokines pro-inflammatoires} (IL-1, TNF-$\alpha$, IL-6).
    \item \textbf{Effets vasculaires (Calor, Rubor, Tumor) :} Vasodilatation artériolaire (chaleur, rougeur) + Augmentation de la perméabilité veineuse (sortie de plasma $\rightarrow$ œdème, hémoconcentration relative). La fièvre résulte de l'action de l'IL-1 et TNF-$\alpha$ sur l'hypothalamus (point de consigne élevé).
    \item \textbf{Recrutement cellulaire (Chimiotaxie) :} Les cytokines attirent les neutrophiles (polynucléaires neutrophiles) depuis le sang vers les tissus infectés (foie, rate, ganglions). \textbf{Explication de la leucopénie :} Les neutrophiles quittent massivement le compartiment sanguin (margeation, extravasation) pour migrer vers les foyers infectieux (foie surtout). La consommation périphérique dépasse la production médullaire (qui peut être freinée par le virus). On parle de « leucopénie de redistribution et de consommation ».
    \item \textbf{Phagocytose et Thrombopénie :} Les neutrophiles et macrophages phagocytent les débris cellulaires et virions. La thrombopénie s'explique par : (a) Destruction médullaire des mégacaryocytes par le virus ; (b) Consommation massive dans les microthromboses vasculaires (coagulation intravasculaire disséminée - CIVD secondaire à l'atteinte endothéliale) ; (c) Séquestration splénique (splénomégalie fréquente).
    \item \textbf{CRP :} Synthétisée par le foie sous l'effet de l'IL-6. Son taux très élevé (120 mg/L) confirme une réponse de phase aiguë intense, mais l'atteinte hépatocellulaire massive (ASAT/ALAT $> 1000$) limite paradoxalement sa synthèse potentielle maximale.
\end{enumerate}
\textbf{Conclusion Q1 :} Le tableau associe une réponse inflammatoire systémique intense (fièvre, CRP++) à une défaillance des effecteurs innés (leucopénie, thrombopénie) due à la tropisme viral pour les organes hématopoïétiques et l'endothélium, marqueur de gravité.

\textbf{Question 2 : Interprétation sérologique — Profil IgM+ / IgG- à J3-J4.}

\textit{Analyse de la courbe cinétique :}
\begin{itemize}
    \item \textbf{J0-J5 :} Taux d'IgM quasi nul puis début d'élévation. Taux d'IgG nul.
    \item \textbf{J7-J10 :} Pic de l'IgM (réponse primaire). Apparition des premières IgG (class switch).
    \item \textbf{J14+ :} Déclin de l'IgM, montée en plateau de l'IgG (mémoire).
\end{itemize}
\textit{Position du patient :} Consultation à J3-J4 après début des symptômes.
\textit{Justification :}
\begin{itemize}
    \item \textbf{IgM Positif (++++) :} Bien que précoces (J3-J4), les tests ELISA de capture IgM modernes sont très sensibles. La positivité forte témoigne d'une \textbf{infection aiguë récente} en cours. C'est le marqueur diagnostique de référence pour le diagnostic de certitude en phase aiguë (OMS).
    \item \textbf{IgG Négatif :} Deux informations majeures :
    \begin{enumerate}
        \item \textbf{Pas d'infection antérieure :} Pas de mémoire immunitaire humorale préexistante.
        \item \textbf{Pas de vaccination efficace :} Le vaccin amaril (virus vivant atténué) induit une réponse IgG neutralisante durable (titre $> 10$ UI/mL ou rapport $> 1$). L'absence d'IgG confirme que le patient \textbf{n'est pas immunisé}, corroborant l'anamnèse (non vacciné).
    \end{enumerate}
\end{itemize}
\textbf{Conclusion Q2 :} Le profil IgM+/IgG- est typique d'une \textbf{primo-infection aiguë} chez un sujet \textbf{naïf immunologiquement} (non vacciné, jamais infecté). C'est la preuve biologique de l'échec de la prévention vaccinale.

\textbf{Question 3 : Analyse de la défaillance vaccinale — Causes individuelles et systémiques.}

\begin{center}
\begin{tabularx}{\linewidth}{|l|X|}\hline
\textbf{Niveau} & \textbf{Facteurs identifiés dans le dossier} \\ \hline
\textbf{Individuel (Patient)} & \begin{itemize}\item Méfiance / Hésitation vaccinale (« méfiance vis-à-vis des campagnes »).
\item Faible littératie en santé (connaissance du risque).
\item Éloignement géographique (15 km, zone riveraine forestière) $\rightarrow$ Coût d'opportunité (temps, transport).
\item Perte/absence de carte de vaccination (traçabilité défaillante).
\end{itemize} \\ \hline
\textbf{Systémique (Système de santé)} & \begin{itemize}
\item \textbf{Rupture de stock} de vaccin amaril au district depuis 3 semaines (rupture d'approvisionnement niveau régional/central).
\item \textbf{Couverture vaccinale insuffisante} (68 \% $< 80 \%$ cible) $\rightarrow$ Immunité collective non acquise.
\item \textbf{Zone oubliée / difficile d'accès :} Village en lisière de forêt, non couvert par les stratégies avancées régulières.
\item \textbf{Délai depuis dernière campagne de masse :} 2022 (4 ans) $\rightarrow$ Accumulation de non-vaccinés (nouvelles naissances, nouveaux arrivants).
\item \textbf{Surveillance entomologique :} Présence avérée d'\emph{Aedes} non quantifiée ni contrôlée.
\end{itemize} \\ \hline
\end{tabularx}
\end{center}
\textbf{Synthèse :} C'est une défaillance \textbf{multifactorielle} : la « hésitation » du patient s'inscrit dans un contexte de \textbf{non-disponibilité de l'offre} (rupture) et d'\textbf{inaccessibilité géographique}. La responsabilité est partagée.

\textbf{Question 4 : Prise en charge thérapeutique — Conduite à tenir au lit du patient.}

\begin{methode}[Prise en charge de la Fièvre Jaune grave (OMS / Guide national Cameroun)]
\begin{enumerate}
\item \textbf{Hospitalisation immédiate en service de réanimation / médecine interne} (surveillance continue : SCORE SOFA, Glasgow, diurèse, hémostase).
\item \textbf{Traitement symptomatique et de support (VITAL) :}
\begin{itemize}
    \item \textbf{Réanimation volémique :} Solutés cristalloïdes (Sérum Salé Isotonique / Ringer Lactate) pour maintenir la perfusion rénale (créatinine 14 mg/L $\rightarrow$ risque IRA). \textbf{Attention :} Surveillance de la surcharge volémique (œdème pulmonaire) car hypoalbuminémie (foie).
    \item \textbf{Correction de l'hypoglycémie :} Perfusion de Glucosé 10 \% ou 5 \% (glycémie 0,65 g/L). Le foie ne fait plus la néoglucogenèse/glycogénolyse.
    \item \textbf{Prévention/Traitement des hémorragies :} Transfusion de \textbf{Concentrés de Plaquettes} si thrombopénie $< 20$ G/L ou saignement actif (hématémèse) + \textbf{Plasma Frais Congelé (PFC)} pour facteurs de coagulation (CIVD). Vitamine K (voie IV) si trouble de l'hémostase (facteurs II, VII, IX, X vit K-dépendants).
    \item \textbf{Protection gastrique :} IPP (Omeprazole IV) contre ulcères de stress / hémorragie digestive.
    \item \textbf{Traitement de la fièvre/douleurs :} \textbf{Paracétamol} (1 g $\times$ 3/4 IV ou PO si toléré). \textbf{INTERDIT :} AINS (Aspirine, Ibuprofène, Diclofénac) $\rightarrow$ majorent le risque hémorragique (antiagrégant) et néphrotoxiques.
    \item \textbf{Antibiotiques probabilistes :} Si suspicion de surinfection bactérienne (septicémie) : Ceftriaxone 2 g/j IV (couvre méningocoque, pneumocoque, salmonelles).
\end{itemize}
\item \textbf{Traitement étiologique :} \textbf{Aucun antiviral spécifique homologué.} (Ribavirine, Favipiravir : efficacité non prouvée en clinique humaine, non recommandés OMS). Recherche d'essais cliniques si disponibles.
\item \textbf{Surveillance biologique rapprochée (J0, J1, J2, J3...) :} NFS, Ionogramme, Urée/Créa, Coagulation (TP, TCA, Fibrinogène), Glycémie, Transaminases, Bilan hépatique.
\item \textbf{Isolement anti-vectoriel :} Moustiquaire imprégnée (PI) autour du lit 24h/24 pour éviter transmission nosocomiale (cycle urbain \emph{Aedes aegypti}).
\end{enumerate}
\end{methode}

\textbf{Question 5 : Mesures de santé publique — Riposte district (Niveau opérationnel).}

\begin{enumerate}
    \item \textbf{Déclaration obligatoire immédiate :} Fièvre Jaune = Maladie à déclaration obligatoire (MDO) au Cameroun (Classe 1). Fiche de notification IDSR (Integrated Disease Surveillance and Response) envoyée au District $\rightarrow$ Région $\rightarrow$ Centrale (DLMEP) $\rightarrow$ OMS en $< 24$h.
    \item \textbf{Investigation de cas / Recherche active :} Équipe d'enquête (Médecin, Infirmier, Agent de santé communautaire) au village Nkolbisson : recherche de cas contacts, cas suspects (fièvre + ictère), décès inexpliqués récents. Prélèvements (sang sec sur papier buvard / tube EDTA) pour confirmation PCR / IgM au Centre Pasteur du Cameroun (CPC) ou Labo National.
    \item \textbf{Vaccination réactive (Riposte) :} Cible : Population du village + zone de santé concernée (rayon 5-10 km), âge 9 mois -- 60 ans. \textbf{Objectif :} Couverture $> 95 \%$ en 7-10 jours. \textbf{Action urgente :} Demande d'acheminement d'urgence de doses vaccinales (stock régional/central/UNICEF) + chaîne de froid. Organisation de postes fixes + avancées + mobiles.
    \item \textbf{Lutte anti-vectorielle (LAV) :} \begin{itemize}
        \item \textbf{Traitement d'espaces :} Nébulisation thermique (brouillard) ou ULV (Ultra Low Volume) aux pyréthrinoïdes dans le village et périphérie (ciblage \emph{Aedes} adultes).
        \item \textbf{Destruction des gîtes larvaires :} Vidange/recouvrement des récipients (pneus, bidons, pots de fleurs), traitement larvicide (Bti - \emph{Bacillus thuringiensis israelensis}) des points d'eau non éliminables.
        \item \textbf{Protection individuelle :} Distribution de Moustiquaires Imprégnées à Longue Durée d'Action (MILDAs) + sensibilisation port vêtements longs, répulsifs.
    \end{itemize}
    \item \textbf{Communication pour le changement social et de comportement (CCSC) :} Messages radio (langues locales : Fulfulde, Français, langues locales), causeries éducatives par ASC/Relais communautaires : « La fièvre jaune tue, le vaccin protège à vie, gratuit, sûr ». Gestion des rumeurs (ex: « vaccin stérilise »).
    \item \textbf{Prise en charge gratuite :} Instruction pour gratuité des soins du cas index et des cas suspects identifiés (politique nationale urgence sanitaire).
\end{enumerate}

\textbf{Question 6 : Note de synthèse médicale pour le Médecin-Chef de District.}

\end{exoresolu}

\begin{center}
\begin{tabularx}{\linewidth}{|X|}\hline
\rowcolor{popblueL} \textbf{NOTE DE SYNTHÈSE MÉDICALE — URGENT / CONFIDENTIEL} \\ \hline
\textbf{À :} Médecin-Chef de District, Bertoua \hfill \textbf{De :} Dr [Nom], Médecin Chef de Poste, CSI Nkolbisson \\
\textbf{Date :} 15 Mars 2026 \hfill \textbf{Objet :} Cas confirmé Fièvre Jaune — M. K., 24 ans — Demande riposte vaccinale d'urgence \\ \hline
\textbf{1. IDENTIFICATION :} M. K., 24 ans, M, Cultivateur, Nkolbisson (Aire de santé CSI Nkolbisson). Non vacciné amaril (carte introuvable, anamnèse négative). \\ \hline
\textbf{2. MOTIF / HISTOIRE :} Fièvre 39,5°C $\times$ 3j, ictère, hémorragies (gingivo, hématémèse), douleurs abdominales. Contexte : zone forestière, exposition \emph{Aedes}, pas de voyage. \\ \hline
\textbf{3. EXAMEN CLINIQUE (J0) :} AG altéré, T° 39,8°C, FC 110, TA 100/60. Ictère franc. Pétéchies, gingivorragie, hématémèse. Hépatomégalie 4 cm douloureuse. Pas de signes méningés. \\ \hline
\textbf{4. BILAN BIOLOGIQUE (J0) :} \textbf{Leucopénie} (2,1 G/L), \textbf{Thrombopénie sévère} (45 G/L). \textbf{Cytolyse hépatique majeure} (ASAT 1850, ALAT 2100). Ictère (Bilit 45 mg/L). \textbf{Inflammation aiguë} (CRP 120 mg/L). Insuffisance rénale naissante (Créa 14). Hypoglycémie (0,65). \\ \hline
\textbf{5. SÉROLOGIE (J0) :} \textbf{IgM Anti-FJ : POSITIF (++++)} / \textbf{IgG Anti-FJ : NÉGATIF}. \textit{Interprétation : Infection aiguë récente chez sujet naïf (non vacciné, non immunisé).} \\ \hline
\textbf{6. DISCUSSION DIAGNOSTIQUE :} \textbf{Cas confirmé Fièvre Jaune (Critères OMS : Clinique compatible + IgM+)}. Diagnostic différentiels écartés : Paludisme grave (GT négatif x2), Hépatite virale A/E (IgM HAV/HEV non demandées mais tableau hémorragique + leucopénie évocateur FJ), Leptospirose (contexte possible mais IgM FJ+ spécifique), Fièvre hémorragique Ebola/Marburg (pas de contact connu, IgM FJ+ oriente). \\ \hline
\textbf{7. CONDUITE À TENIR (Patient) :} Hospitalisation réanimation (CSI Bertoua / Hôpital District). Réanimation volémique, correction hypoglycémie, transfusion plaquettes + PFC (saignement actif), Vit K, IPP, Paracétamol seul, AB probabiliste (Ceftriaxone). Surveillance biologique 6h/24h. Isolement moustiquaire. Pronostic réservé (Score de gravité élevé : cytolyse $> 1000$, thrombopénie $< 50$, hémorragie digestive, hypoglycémie, insuffisance rénale). \\ \hline
\textbf{8. MESURES DE SANTÉ PUBLIQUE (DEMANDE D'ACTION IMMÉDIATE) :} \\
\begin{itemize}
    \item \textbf{Déclaration IDSR} envoyée ce jour 10h00.
    \item \textbf{Enquête terrain} programmée demain 8h (Équipe district + ASC Nkolbisson) : recherche cas contacts, prélèvements confirmation CPC.
    \item \textbf{DEMANDE URGENTE DE DOSES VACCIN AMARIL} (cible $\approx$ 2 500 pers zone) pour \textbf{Vaccination Réactive} dès réception (objectif J+7). Rupture stock district critique.
    \item \textbf{LAV d'urgence} : Nébulisation village Nkolbisson + destruction gîtes larvaires (ASC + Mairie).
    \item \textbf{CCSC} : Message radio ce soir + causeries ASC demain (lever freins vaccination).
    \item \textbf{Gratuité soins} : Instruction pour cas index et cas suspects.
\end{itemize} \\
\textbf{Signature :} Dr [Nom] \hfill \textbf{Visa MCD :} \ldots\ldots\ldots \\ \hline
\end{tabularx}
\end{center}


\newpage

\coursec{Méthodes et exercices résolus}

\coursec{Réponse immunitaire}

\courssub{Les barrières de première ligne : peau et muqueuses}

\begin{methode}[Identifier les lignes de défense de l'organisme]
\textbf{Objectif :} Classer un mécanisme de défense dans la bonne ligne (1\textsuperscript{re}, 2\textsuperscript{e} ou 3\textsuperscript{e} ligne) et préciser s'il est spécifique ou non.
\begin{enumerate}
 \item \textbf{Repérer le mécanisme} décrit dans l'énoncé ou le document (ex: peau, cils vibratiles, lysozyme, inflammation, phagocytose, anticorps, lymphocytes).
 \item \textbf{Rappeler les caractéristiques} de chaque ligne :
 \begin{itemize}
 \item \textbf{1\textsuperscript{re} ligne (barrières) :} Mécanismes \textbf{non spécifiques}, \textbf{innés}, \textbf{constitutifs} (toujours présents). Ex: peau (kératine, pH acide, flore), muqueuses (mucus, cils, lysozyme, cils vibratiles, réflexes de toux/éternuement).
 \item \textbf{2\textsuperscript{e} ligne (innée interne) :} Mécanismes \textbf{non spécifiques}, \textbf{innés}, \textbf{induits} (déclenchés par l'agression). Ex: réaction inflammatoire (rubor, calor, tumor, dolor), phagocytose (polynucléaires neutrophiles, macrophages), fièvre, interférons, système du complément.
 \item \textbf{3\textsuperscript{e} ligne (adaptative/spécifique) :} Mécanismes \textbf{spécifiques}, \textbf{acquis}, \textbf{induits}, avec \textbf{mémoire}. Acteurs : Lymphocytes B (anticorps) et T (cytotoxiques, auxiliaires).
 \end{itemize}
 \item \textbf{Classer} le mécanisme en cochant les cases : Non spécifique / Spécifique ; Inné / Acquis ; Constitutif / Induit.
 \item \textbf{Justifier} par l'absence ou la présence de reconnaissance d'un antigène précis et de mémoire immunologique.
\end{enumerate}
\cle{Réflexe :} Si le mécanisme agit contre \emph{tout} agent pathogène sans distinction $\rightarrow$ Non spécifique (1\textsuperscript{re} ou 2\textsuperscript{e} ligne). S'il cible \emph{un} antigène précis et laisse une mémoire $\rightarrow$ Spécifique (3\textsuperscript{e} ligne).
\end{methode}

\begin{exoresolu}[Analyse d'un cas de contamination cutanée]
\textbf{Énoncé :} Un enfant de 10 ans se blesse au genou en jouant sur un terrain poussiéreux à Yaoundé. La plaie saigne, devient rouge, chaude, gonflée et douloureuse au bout de quelques heures. Deux jours plus tard, du pus est visible. Le médecin nettoie la plaie et prescrit un antiseptique.
\begin{enumerate}
 \item Citez les deux premières lignes de défense mises en jeu dans cette situation.
 \item Expliquez l'origine des quatre signes cliniques observés (rougeur, chaleur, gonflement, douleur).
 \item Identifiez la nature cellulaire du pus et son rôle.
 \item Pourquoi le médecin prescrit-il un antiseptique alors que l'organisme possède ses propres défenses ?
\end{enumerate}
\tcblower
\textbf{Solution détaillée :}
\begin{enumerate}
 \item \textbf{Lignes de défense mobilisées :}
 \begin{itemize}
 \item \textbf{1\textsuperscript{re} ligne :} La \textbf{peau} (barrière physique et chimique : kératine, film hydrolipidique acide, flore commensale) a été franchie par la blessure.
 \item \textbf{2\textsuperscript{e} ligne :} La \textbf{réaction inflammatoire} (signes cliniques) et la \textbf{phagocytose} (formation de pus) se déclenchent immédiatement après la contamination par les bactéries du sol/poussière.
 \end{itemize}
 \item \textbf{Origine des signes inflammatoires (vasodilatation et hyperperméabilité vasculaire) :}
 \begin{itemize}
 \item \textbf{Rougeur (Rubor) et Chaleur (Calor) :} Résultent de la \textbf{vasodilatation} des artérioles locales (sous l'effet de médiateurs comme l'histamine, la sérotonine libérés par les mastocytes et les basophiles). L'afflux sanguin massif amène chaleur et couleur rouge.
 \item \textbf{Gonflement (Tumor) :} Dû à l'\textbf{hyperperméabilité} des veines post-capillaires (écartement des cellules endothéliales). Le plasma (eau, protéines, facteurs du complément, anticorps préexistants) sort dans le tissu interstitiel $\rightarrow$ œdème.
 \item \textbf{Douleur (Dolor) :} Causée par la compression des terminaisons nerveuses par l'œdème et par l'action directe de médiateurs chimiques (bradykinine, prostaglandines) sur les nocicepteurs.
 \end{itemize}
 \item \textbf{Nature et rôle du pus :} Le pus est un liquide jaunâtre composé principalement de \textbf{polynucléaires neutrophiles morts} (phagocytes ayant ingéré des bactéries), de débris cellulaires, de bactéries (vivantes ou mortes) et de lymphe. Il témoigne de l'activité intense de la \textbf{phagocytose} : les polynucléaires sont sortis du sang (diapédèse), ont migré vers le foyer (chimiotaxisme), ont phagocyté les bactéries et sont morts en libérant leurs enzymes lysosomales.
 \item \textbf{Rôle de l'antiseptique :} Les défenses innées (inflammation, phagocytose) sont \textbf{non spécifiques} et peuvent être débordées par un nombre trop élevé de germes pathogènes (ex: \emph{Clostridium tetani} tétanos, \emph{Staphylococcus aureus}). L'antiseptique (ex: chlorhexidine, povidone-iodée) \textbf{réduit la charge bactérienne} localement, facilitant le travail des phagocytes et prévenant la généralisation de l'infection (septicémie) ou le tétanos (risque en milieu tropical poussiéreux). C'est une aide externe, non un substitut à l'immunité.
\end{enumerate}
\textbf{Conclusion :} La blessure franchit la 1\textsuperscript{re} ligne ; la 2\textsuperscript{e} ligne (inflammation + phagocytose) se déploie pour contenir l'infection ; l'antiseptique est un appoint médical indispensable en cas de forte contamination.
\end{exoresolu}

\courssub{La réaction inflammatoire et la phagocytose : la deuxième ligne innée}

\begin{methode}[Décrire les étapes de la réaction inflammatoire et de la phagocytose]
\textbf{Objectif :} Rédiger un texte explicatif séquencé (chronologique) ou légender un schéma fonctionnel.
\begin{enumerate}
 \item \textbf{Phase vasculaire (minutes) :} 
 \begin{itemize}
 \item Aggression $\rightarrow$ Libération de médiateurs (histamine, sérotonine) par mastocytes/basophiles.
 \item Vasodilatation artériolaire $\rightarrow$ Rubor, Calor.
 \item Hyperperméabilité veineuse $\rightarrow$ Tumor (œdème), Dolor (compression + médiateurs).
 \end{itemize}
 \item \textbf{Phase cellulaire (heures) :}
 \begin{enumerate}
 \item \textbf{Margination / Roulement :} Polynucléaires neutrophiles (PNN) se collent à l'endothélium (molécules d'adhésion).
 \item \textbf{Diapédèse :} PNN traversent la paroi vasculaire (jonctions inter-endothéliales élargies).
 \item \textbf{Chimiotaxisme :} Migration dirigée vers le foyer par gradient de chémokines (facteurs chimiotactiques).
 \item \textbf{Phagocytose (5 étapes) :}
 \begin{enumerate}
 \item \textbf{Adhérence / Reconnaissance :} Reconnaissance du germe (opsonisation par anticorps/complément facilite).
 \item \textbf{Ingestion :} Formation de pseudopodes $\rightarrow$ Vacuole (phagosome).
 \item \textbf{Fusion :} Phagosome + Granules lysosomaux (enzymes hydrolytiques, radicaux libres) $\rightarrow$ \textbf{Phagolysosome}.
 \item \textbf{Digestion :} Dégradation intracellulaire du germe.
 \item \textbf{Exocytose / Rejet :} Élimination des résidus indigestes.
 \end{enumerate}
 \end{enumerate}
 \item \textbf{Résolution / Réparation :} Macrophages (monocytes différenciés) nettoient les débris, produisent des facteurs de croissance $\rightarrow$ cicatrisation. Si échec $\rightarrow$ abcès chronique ou généralisation.
\end{enumerate}
\cle{Vocabulaire imposé :} Vasodilatation, hyperperméabilité, diapédèse, chimiotaxisme, opsonisation, phagosome, phagolysosome, digestion intracellulaire. Ne pas confondre \emph{ingestion} (formation phagosome) et \emph{digestion} (action enzymatique).
\end{methode}

\begin{exoresolu}[Interprétation d'une courbe cinétique de phagocytose]
\textbf{Énoncé :} On réalise une mise en culture de macrophages péritonéaux de souris avec des bactéries \emph{Escherichia coli} opsonisées. On compte, à intervalles réguliers, le nombre moyen de bactéries internalisées par macrophage. Résultats :
\begin{center}
\begin{tabular}{|c|c|c|c|c|c|c|}\hline
Temps (min) & 0 & 15 & 30 & 60 & 90 & 120 \\ \hline
Nb bactéries/macrophage & 0 & 12 & 28 & 45 & 48 & 48 \\ \hline
\end{tabular}
\end{center}
\begin{enumerate}
 \item Représentez graphiquement l'évolution du nombre de bactéries internalisées en fonction du temps.
 \item Identifiez les trois phases cinétiques et reliez-les aux étapes cellulaires de la phagocytose.
 \item Quelle est la capacité maximale de phagocytose de ces macrophages dans ces conditions ? Justifiez.
 \item Prédisez l'allure de la courbe si les bactéries \emph{n'étaient pas} opsonisées. Justifiez.
\end{enumerate}
\tcblower
\textbf{Solution détaillée :}
\begin{enumerate}
 \item \textbf{Graphique :}
 \begin{popfigure}
 \begin{tikzpicture}
 \begin{axis}[
 width=\linewidth, height=0.6\linewidth,
 xlabel={Temps (min)}, ylabel={Nb bactéries / macrophage},
 xmin=0, xmax=125, ymin=0, ymax=55,
 xtick={0,15,30,60,90,120}, ytick={0,12,28,45,48},
 grid=major, grid style={popline},
 popcurve, mark=*, mark size=2pt
 ]
 \addplot coordinates {
 (0,0) (15,12) (30,28) (60,45) (90,48) (120,48)
 };
 \end{axis}
 \end{tikzpicture}
 \legende{Cinétique de la phagocytose d'\emph{E. coli} opsonisé par des macrophages péritonéaux de souris.}
 \end{popfigure}
 \item \textbf{Phases cinétiques :}
 \begin{itemize}
 \item \textbf{Phase de latence (0--15 min) :} Pente faible. Correspond à la \textbf{reconnaissance}, l'\textbf{adhérence} des bactéries aux récepteurs membranaires et le début de la formation des pseudopodes (ingestion).
 \item \textbf{Phase exponentielle / linéaire rapide (15--60 min) :} Pente forte. Correspond à l'\textbf{ingestion massive} (formation successive de phagosomes) et à la \textbf{fusion phagosome-lysosome} (digestion active). Les macrophages sont pleinement actifs.
 \item \textbf{Phase de plateau / saturation (60--120 min) :} Pente nulle. Correspond à la \textbf{saturation} de la capacité phagocytaire (stockage des résidus, épuisement des granules lysosomaux, inhibition par rétroaction).
 \end{itemize}
 \item \textbf{Capacité maximale :} \textbf{48 bactéries par macrophage}. Justification : La courbe atteint un palier constant (plateau) à partir de 60 min, indiquant que le macrophage ne peut plus internaliser ni digérer de nouvelles proies.
 \item \textbf{Prédiction sans opsonisation :} La courbe présenterait une \textbf{phase de latence plus longue}, une \textbf{pente moins raide} en phase exponentielle et un \textbf{plateau plus bas} (capacité maximale inférieure).
 \textbf{Justification :} L'opsonisation (recouvrement par anticorps et/ou complément) crée des "poignées" reconnues par les récepteurs spécifiques du macrophage. Sans opsonisation, la reconnaissance repose sur des récepteurs de type "pattern recognition" moins efficaces pour l'ingestion $\rightarrow$ adhérence plus lente, ingestion moins rapide, capacité réduite.
\end{enumerate}
\textbf{Conclusion :} L'opsonisation est un facteur clé de l'efficacité de la phagocytose (2\textsuperscript{e} ligne), reliant d'ailleurs l'immunité humorale (anticorps, 3\textsuperscript{e} ligne) à l'immunité cellulaire innée.
\end{exoresolu}

\courssub{L'immunité spécifique (adaptative) : les lymphocytes, acteurs de la spécificité et de la mémoire}

\begin{methode}[Comparer immunité non spécifique et immunité spécifique]
\textbf{Objectif :} Remplir un tableau comparatif ou rédiger une argumentation structurée.
\begin{enumerate}
 \item \textbf{Construire un tableau} à 2 colonnes (Non spécifique / Spécifique) et lignes critères imposés.
 \item \textbf{Remplir rigoureusement} chaque case avec le vocabulaire du programme :
 \begin{center}
 \begin{tabular}{|l|>{\raggedright\arraybackslash}p{5cm}|>{\raggedright\arraybackslash}p{5cm}|}\hline
 \textbf{Critère} & \textbf{Non spécifique (Innée)} & \textbf{Spécifique (Adaptative)} \\ \hline
 \textbf{Spécificité} & Large spectre (tous pathogènes) & Étroite : 1 antigène = 1 réponse \\ \hline
 \textbf{Appartenance} & Innée (génétique, héréditaire) & Acquise (après contact Ag) \\ \hline
 \textbf{Kinétique} & Immédiate (min à heures) & Retardée (jours : phase de latence) \\ \hline
 \textbf{Mémoire} & \textbf{Aucune} (réponse identique) & \textbf{Oui} (réponse 2\textsuperscript{o} plus rapide/forte) \\ \hline
 \textbf{Acteurs cellulaires} & Phagocytes (PNN, Macrophages), NK, Mastocytes, Éosinophiles & Lymphocytes B et T (clones spécifiques) \\ \hline
 \textbf{Acteurs moléculaires} & Lysozyme, Complément, Interférons, Cytokines inflammatoires & \textbf{Anticorps (Ig)}, Récepteurs TCR/BCR, Cytokines lymphoïdiques \\ \hline
 \textbf{Reconnaissance} & Récepteurs motifs conservés (PAMPs) & BCR / TCR $\rightarrow$ Épitope précis (peptide + CMH) \\ \hline
 \end{tabular}
 \end{center}
 \item \textbf{Souligner la coopération :} L'innée \textbf{active} l'adaptative (cellules présentatrices d'antigène : Macrophages, Cellules dendritiques $\rightarrow$ Lymphocytes T CD4$^+$). L'adaptative \textbf{potencie} l'innée (Anticorps $\rightarrow$ Opsonisation ; Cytokines $\rightarrow$ Activation macrophages).
\end{enumerate}
\cle{Piège :} Ne pas dire "L'immunité innée n'a pas de mémoire". Dire : "Elle ne possède \textbf{pas de mémoire immunologique spécifique}".
\end{methode}

\begin{exoresolu}[Analyse documentée : Vaccination antivariolique historique]
\textbf{Énoncé :} Document 1 : En 1796, Edward Jenner inocule à un enfant du pus prélevé sur une pustule de vaccine (maladie bénigne de la vache, due à un \emph{Orthopoxvirus} proche de la variole). L'enfant développe une légère fièvre locale. Semaines plus tard, Jenner l'inocule avec du pus variolique (variole humaine, mortelle). L'enfant ne tombe pas malade.
Document 2 : Courbe du taux d'anticorps anti-varioliques dans le sang de cet enfant au cours du temps (schéma qualitatif : 1\textsuperscript{er} pic lent et bas après vaccine ; 2\textsuperscript{ème} pic rapide, haut, précoce après variole).
\begin{enumerate}
 \item En vous basant sur le Document 1, expliquez pourquoi l'inoculation de vaccine a protégé l'enfant contre la variole.
 \item En vous basant sur le Document 2, caractérisez la réponse immunitaire secondaire (mémoire) par rapport à la réponse primaire.
 \item Identifiez le type de vaccin utilisé par Jenner (vaccin vivant atténué, inactivé, sous-unité, ARN messager). Justifiez.
 \item Proposez une explication au fait que la vaccine (virus bovin) protège contre la variole (virus humain).
\end{enumerate}
\tcblower
\textbf{Solution détaillée :}
\begin{enumerate}
 \item \textbf{Mécanisme de protection (Vaccination) :} Le virus de la vaccine (antigène) a pénétré l'organisme. Il a déclenché une \textbf{réponse immunitaire spécifique primaire} : activation de lymphocytes B et T spécifiques des antigènes de la vaccine. Cette réponse a conduit à la production d'\textbf{anticorps} et à la mise en place de \textbf{lymphocytes mémoire} (B et T) spécifiques de ces antigènes. Comme les virus de la vaccine et de la variole sont \textbf{très proches antigéniquement} (épitopes communs ou croisés), les lymphocytes mémoire anti-vaccine reconnaissent le virus de la variole lors de la seconde inoculation.
 \item \textbf{Caractéristiques de la réponse secondaire (Doc 2) :}
 \begin{itemize}
 \item \textbf{Latence plus courte :} Pic atteint en quelques jours (vs 10-15 jours pour la primaire).
 \item \textbf{Amplitude plus élevée :} Taux d'anticorps maximal bien supérieur.
 \item \textbf{Kinétique plus rapide :} Montée brutale.
 \item \textbf{Affinité plus forte :} Anticorps de meilleure qualité (maturation de l'affinité).
 \item \textbf{Persistance plus longue.}
 \end{itemize}
 \textit{Justification :} Présence de lymphocytes mémoire (plus nombreux, réactifs, ne nécessitant pas de phase de sélection clonale longue).
 \item \textbf{Type de vaccin :} \textbf{Vaccin vivant atténué} (ou plutôt \textbf{vaccin vivant hétérologue / à virus apparenté non pathogène pour l'homme}).
 \textbf{Justification :} Jenner utilise du pus frais contenant le \textbf{virus vivant de la vaccine}, capable de se répliquer modérément dans l'organisme (fièvre locale) pour stimuler fortement l'immunité, mais non pathogène (bénin) pour l'homme. Ce n'est pas un virus inactivé (tué) ni une sous-unité purifiée.
 \item \textbf{Explication de la protection croisée :} Les \emph{Orthopoxvirus} (vaccine, variole, variole du singe) partagent une \textbf{forte homologie génomique et antigénique}. Ils possèdent des protéines de surface très conservées. Les anticorps et lymphocytes T dirigés contre des \textbf{épitopes conservés (communs)} des protéines de la vaccine reconnaissent les mêmes épitopes sur le virus de la variole $\rightarrow$ \textbf{Réactivité croisée (cross-reactivity)}.
\end{enumerate}
\textbf{Conclusion :} Cette expérience fondatrice illustre le principe de la vaccination : \textbf{mimer l'infection naturelle par un agent inoffensif mais antigéniquement proche} pour établir une mémoire protectrice sans payer le prix de la maladie.
\end{exoresolu}

\courssub{Coopération, régulation et applications concrètes : Vaccination et Sérothérapie}

\begin{methode}[Expliquer le principe et l'intérêt de la vaccination et de la sérothérapie]
\textbf{Objectif :} Différencier clairement immunité active (vaccin) et immunité passive (sérum) et savoir les situer dans le temps (prophylaxie vs urgence).
\begin{enumerate}
 \item \textbf{Vaccination (Immunité Active Artificielle) :}
 \begin{itemize}
 \item \textbf{Principe :} Introduction d'un \textbf{antigène} (agent vivant atténué, inactivé, sous-unité protéique, ARN messager, vecteur viral) $\rightarrow$ Déclenchement réponse immune \textbf{active} du sujet (lymphocytes B/T $\rightarrow$ Anticorps + Cellules mémoire).
 \item \textbf{Délai de protection :} \textbf{Long} (semaines : temps de la réponse primaire). Nécessite souvent \textbf{rappels} (boosters) pour entretenir la mémoire.
 \item \textbf{Durée :} \textbf{Longue} (années, vie entière).
 \item \textbf{Indication :} \textbf{Prophylaxie} (prévention avant exposition). Ex: PEV (Programme Élargi de Vaccination) au Cameroun : BCG, Polio, Rougeole, Fièvre jaune, Hépatite B, Tétanos, Pneumo, Rotavirus, HPV.
 \end{itemize}
 \item \textbf{Sérothérapie (Immunité Passive Artificielle) :}
 \begin{itemize}
 \item \textbf{Principe :} Injection d'\textbf{anticorps tout faits} (Immunoglobulines / Sérum hyperimmune) produits par un \textbf{autre organisme} (cheval, humain donneur, ou monoclonaux recombinants).
 \item \textbf{Délai de protection :} \textbf{Immédiat} (anticorps fonctionnels dès injection).
 \item \textbf{Durée :} \textbf{Courte} (semaines à mois : demi-vie des IgG $\approx$ 21-28 jours). \textbf{Pas de mémoire} (pas de lymphocytes du receveur activés).
 \item \textbf{Risque :} Maladie du sérum (réaction immunitaire contre les protéines étrangères du sérum, ex: protéines équines), transmission d'agents infectieux (dépistage rigoureux).
 \item \textbf{Indication :} \textbf{Urgence thérapeutique} (infection déclarée ou exposition à haut risque immédiat). Ex: Sérum antivenimeux (morsure serpent), Sérum antitétanique (blessure sale + statut vaccinal inconnu/incomplet), Immunoglobulines anti-Hépatite B (néoné de mère AgHBs+), Anti-Rage (post-exposition + vaccin).
 \end{itemize}
 \item \textbf{Association Vaccin + Sérum (Prophylaxie post-exposition) :} Ex: Rage, Tétanos, Hépatite B. Le sérum protège \textbf{immédiatement} le temps que le vaccin induise l'\textbf{immunité active durable}. Injection en \textbf{sites distincts} (interférence possible si même site).
\end{enumerate}
\cle{Mots-clés BEPC :} "Immunité active acquise", "Immunité passive acquise", "Antigène", "Anticorps", "Cellules mémoire", "Prophylaxie", "Thérapeutique", "Rappel vaccinal".
\end{methode}

\begin{exoresolu}[Cas clinique : Prise en charge d'une morsure de serpent au Cameroun]
\textbf{Énoncé :} Un paysan de 35 ans est mordu au pied par un serpent non identifié (probablement \emph{Echis ocellatus} - vipère à écailles) dans une plantation de cacao près de Dschang. Il arrive au centre de santé 2 heures plus tard. Le personnel dispose de sérum antivenimeux polyvalent (chevaux immunisés) et de vaccin antitetanique.
\begin{enumerate}
 \item Quel type d'immunité (active/passive, naturelle/artificielle) confère le sérum antivenimeux ? Justifiez.
 \item Pourquoi injecte-t-on ce sérum en urgence et non un vaccin antivenimeux ?
 \item Le médecin prescrit aussi le vaccin antitetanique (anatoxine tétanique). Quel type d'immunité ce vaccin induit-il ? Pourquoi l'associer au sérum ici ?
 \item Citez deux précautions majeures à prendre lors de l'injection du sérum antivenimeux (origine équine).
\end{enumerate}
\tcblower
\textbf{Solution détaillée :}
\begin{enumerate}
 \item \textbf{Immunité passive artificielle.}
 \textbf{Justification :} "Passive" car le receveur reçoit des \textbf{anticorps tout faits} (IgG purifiées du sang de chevaux hyperimmunisés) et ne produit pas les siens (pas d'activation de ses lymphocytes B). "Artificielle" car ces anticorps sont obtenus par une intervention humaine (immunisation du cheval, purification, injection) et non par transfert naturel (placenta/allaitement).
 \item \textbf{Urgence du sérum vs Vaccin :} L'envenimation par \emph{Echis ocellatus} provoque des troubles de la coagulation (hémorragies, hémolyse) et des nécroses \textbf{très rapides} (heures). La réponse immune \textbf{active} (vaccin) nécessite une \textbf{phase de latence de 10 à 15 jours} pour produire des anticorps protecteurs. Le patient mourrait avant. Le sérum apporte une \textbf{neutralisation immédiate} des toxines circulantes (effet "antidote").
 \item \textbf{Vaccin antitetanique :} \textbf{Immunité active artificielle.} L'anatoxine (toxine inactivée, non toxique mais immunogène) stimule les lymphocytes B du patient $\rightarrow$ production d'anticorps anti-toxine tétanique + cellules mémoire.
 \textbf{Association :} La morsure/sol est une porte d'entrée majeure pour \emph{Clostridium tetani} (tétanos). Le sérum antivenimeux ne protège \textbf{pas} contre le tétanos. Le vaccin antitetanique met trop de temps à agir pour l'urgence tétanique potentielle (incubation 4-21 jours). \textit{Note :} Si le statut vaccinal est incertain, on associe souvent \textbf{Sérum Antitetanique (SAT) = immunité passive immédiate} + \textbf{Vaccin Antitetanique (VAT) = immunité active différée} (injection en sites distincts).
 \item \textbf{Précautions sérum équin :}
 \begin{enumerate}
 \item \textbf{Test de sensibilité (Hypersensibilité de type I) :} Injection intradermique de 0,1 mL de sérum dilué 1/10 ; surveillance 15-20 min (recherche papule/érythème/démangeaison). Si positif $\rightarrow$ désensibilisation progressive ou antihistaminiques/corticoïdes avant injection IV lente.
 \item \textbf{Injection par voie intraveineuse (IV) LENTE :} Dilution dans du sérum physiologique ou glucose 5\% ; débit lent (ex: 10-20 gouttes/min au début) sous surveillance clinique (TA, FC, SpO2) pour détecter tout choc anaphylactique retardé ou maladie du sérum (fièvre, arthralgies, éruption cutanée 7-10 jours après).
 \end{enumerate}
\end{enumerate}
\textbf{Conclusion :} La sérothérapie est un traitement d'urgence "prêt à l'emploi" (passif) indispensable face aux envenimations mortelles rapides, en attendant/complétant la mise en place d'une protection active durable par vaccination (tétanos).
\end{exoresolu}

\begin{methode}[Analyser un arbre généalogique ou un pedigree immunologique (transmission d'une immunodéficience)]
\textbf{Objectif :} Déduire le mode de transmission d'une maladie génétique affectant l'immunité (ex: Agammaglobulinémie de Bruton, SCID) à partir d'un pedigree.
\begin{enumerate}
 \item \textbf{Identifier les symboles :} Carré = Homme, Cercle = Femme, Symbole plein = Malade, Symbole vide = Sain, Losange = Sexe inconnu/avortement. Ligne horizontale = Union, Ligne verticale = Descendance.
 \item \textbf{Tester les hypothèses de transmission :}
 \begin{itemize}
 \item \textbf{Récessif autosomique :} Malades de sexe différent, parents sains (porteurs), consanguinité possible, fratrie atteinte $\approx$ 1/4.
 \item \textbf{Dominant autosomique :} Malades dans chaque génération, transmission par hommes et femmes, 1 parent malade $\rightarrow$ 1/2 enfants malades.
 \item \textbf{Récessif lié à l'X (hémizygote mâle) :} \textbf{Quasi-exclusivement des hommes} malades. Femmes porteuses généralement saines. Transmission \textbf{jamais père $\rightarrow$ fils}. Père malade $\rightarrow$ toutes les filles porteuses. Mère porteuse $\rightarrow$ 1/2 fils malades, 1/2 filles porteuses.
 \item \textbf{Dominant lié à l'X :} Femmes plus touchées (2X), transmission père $\rightarrow$ toutes les filles.
 \end{itemize}
 \item \textbf{Conclure} en citant les arguments observés sur le pedigree (ex: "Atteinte only males, pas de transmission père-fils $\rightarrow$ Récessif lié à l'X").
 \item \textbf{Relier au phénotype immunitaire :} Ex: Bruton (\textit{BTK} sur X) $\rightarrow$ Pas de lymphocytes B $\rightarrow$ Pas d'Ig $\rightarrow$ Infections pyogènes à répétition. SCID (\textit{IL2RG} sur X ou \textit{ADA} autosomique) $\rightarrow$ Pas de lymphocytes T (et B/NK) $\rightarrow$ Infections graves opportunistes précoces.
\end{enumerate}
\cle{Rappel 3ème :} Le programme ne demande pas l'analyse formelle de pedigree complexes, mais la compréhension que des \textbf{défauts génétiques} (mutation) peuvent toucher les gènes codant pour des récepteurs (TCR, BCR), des enzymes (ADA), des cytokines (IL-2R$\gamma$), des protéines de signalisation (BTK) $\rightarrow$ \textbf{Immunodéficiences primitives}.
\end{methode}

\begin{exoresolu}[Étude d'une immunodéficience primitive : Agammaglobulinémie de Bruton]
\textbf{Énoncé :} Un nourrisson de 6 mois présente des infections bactériennes respiratoires et ORL à répétition (\emph{Streptococcus pneumoniae}, \emph{Haemophilus influenzae}) depuis l'âge de 3 mois. Le dosage des immunoglobulines sériques montre : IgG < 0,1 g/L (N: 5-12), IgA < 0,05 g/L (N: 0,5-2), IgM < 0,05 g/L (N: 0,5-1,5). L'hémogramme montre une numération lymphocytaire normale. Le test de stimulation mitogène (PHA) sur lymphocytes T est normal. L'étude génétique révèle une mutation du gène \textit{BTK} (Xq22).
\begin{enumerate}
 \item Quel type d'immunité est défaillant chez ce patient ? Justifiez par les résultats biologiques.
 \item Expliquez pourquoi les infections débutent vers 3-6 mois et non à la naissance.
 \item Précisez le mode de transmission de cette maladie. Justifiez par la génétique.
 \item Pourquoi les lymphocytes T fonctionnent-ils normalement ?
 \item Proposez le traitement de référence actuel pour cette pathologie.
\end{enumerate}
\tcblower
\textbf{Solution détaillée :}
\begin{enumerate}
 \item \textbf{Immunité humorale (médiée par les anticorps / Lymphocytes B) défaillante.}
 \textbf{Justification :} \textbf{Hypogammaglobulinémie profonde} (taux d'IgG, IgA, IgM quasi nuls) alors que la numération lymphocytaire est normale et la fonction T (test PHA) est conservée. Cela indique un \textbf{blocage de la différenciation des lymphocytes B} en plasmocytes sécréteurs d'anticorps. Les lymphocytes B précurseurs sont absents ou non fonctionnels.
 \item \textbf{Début à 3-6 mois :} Protection passive transitoire par les \textbf{IgG maternelles} transférées activement via le placenta (récepteur FcRn) durant le 3\textsuperscript{ème} trimestre. La demi-vie des IgG maternelles est d'environ \textbf{21-28 jours}. À 3 mois, le taux est divisé par $\approx$ 4 ; à 6 mois, il est quasi nul. L'enfant se retrouve sans anticorps propres (déficit B) et sans anticorps maternels $\rightarrow$ fenêtre de vulnérabilité aux bactéries encapsulées (pneumocoque, Haemophilus).
 \item \textbf{Transmission récessive liée à l'X.}
 \textbf{Justification :} Gène \textit{BTK} (Bruton Tyrosine Kinase) localisé sur le \textbf{chromosome X (Xq22)}. Les hommes (XY) n'ont qu'un allèle $\rightarrow$ mutation = maladie. Les femmes (XX) ont un allèle sain $\rightarrow$ porteuses saines (inactivation X aléatoire). C'est pourquoi quasi-exclusivement des garçons sont atteints.
 \item \textbf{Lymphocytes T normaux :} Le gène \textit{BTK} code une tyrosine kinase essentielle pour la transduction du signal du \textbf{BCR (Récepteur des Lymphocytes B)}. Il n'est \textbf{pas exprimé} (ou non essentiel) dans la lignée T. Le développement thymique et la signalisation via le \textbf{TCR} sont indépendants de BTK. Donc immunité cellulaire (T) intacte $\rightarrow$ défense contre virus, champignons, mycobactéries préservée.
 \item \textbf{Traitement de référence :} \textbf{Immunoglobulines polyvalentes (IgIV ou IgSC) à vie, toutes les 3-4 semaines (IV) ou 1-2 semaines (SC).}
 \begin{itemize}
 \item Apporte les anticorps manquants (immunité passive artificielle régulière).
 \item Prévenir les infections bactériennes graves.
 \item Association possible : Antibioprophylaxie (ex: Cotrimoxazole), Vaccins inactivés (oui), \textbf{Vaccins vivants atténués CONTRE-INDICQUÉS} (risque de maladie vaccinale non contrôlée par anticorps).
 \end{itemize}
 \textit{Perspective :} Thérapie génique (vecteur lentiviral corrigé \textit{BTK} dans cellules souches hématopoïétiques) en essais cliniques prometteurs.
\end{enumerate}
\textbf{Conclusion :} L'agammaglobulinémie de Bruton est le modèle type du déficit isolé de l'immunité humorale (B) d'origine génétique (X-récessif), traitée par substitution anticorps régulière.
\end{exoresolu}

\courssub{Activités d'intégration : Synthèse et Résolution de problèmes}

\begin{methode}[Rédiger une conclusion scientifique argumentée (CER : Claim - Evidence - Reasoning)]
\textbf{Objectif :} Structurer la réponse finale d'un exercice de synthèse ou d'un TP pour valider la compétence "Rédiger et justifier une démarche, une réponse ou une conclusion".
\begin{enumerate}
 \item \textbf{Claim (Affirmation / Réponse à la question posée) :} Une phrase courte, directe, qui répond \textbf{exactement} à la problématique. Ex: "La vaccination induit une immunité active durable grâce à la mise en place de lymphocytes mémoire."
 \item \textbf{Evidence (Preuves / Données) :} Citez \textbf{explicitement} les résultats d'expérience, les valeurs numériques du tableau, les observations du document, les schémas qui soutiennent l'affirmation. Ex: "Le document 2 montre que lors du 2\textsuperscript{ème} contact, le pic d'anticorps est atteint en 5 jours (vs 15 jours) et atteint 150 UI/mL (vs 20 UI/mL)."
 \item \textbf{Reasoning (Raisonnement / Mécanisme) :} Expliquez \textbf{POURQUOI} les preuves soutiennent l'affirmation, en mobilisant les connaissances théoriques (mécanismes cellulaires/moléculaires). Ex: "Ce rappel rapide et intense s'explique par la présence de lymphocytes B mémoire, déjà différenciés et en plus grand nombre, qui se transforment en plasmocytes sans passer par la phase de sélection clonale longue nécessaire lors du 1\textsuperscript{er} contact."
 \item \textbf{Rebuttal (Optionnel - Niveau excellence) :} Anticipez une objection ou précisez une limite. Ex: "Cependant, cette mémoire peut s'estomper sans rappels vaccinaux (baisse du taux d'anticorps sous seuil protecteur)."
 \item \textbf{Mise en forme :} Paragraphe unique, connecteurs logiques (\textit{donc, car, or, en effet, par conséquent, ainsi}), vocabulaire précis, pas de jargon.
\end{enumerate}
\cle{Modèle de phrase type :} "En conclusion, [Réponse]. En effet, [Preuve chiffrée/observation]. Cela s'explique par [Mécanisme biologique]."
\end{methode}

\begin{exoresolu}[Synthèse : Réponse immune face à une primo-infection vs ré-infection]
\textbf{Énoncé :} À partir de l'ensemble des connaissances acquises sur l'immunité non spécifique et spécifique, rédigez un texte argumenté d'une vingtaine de lignes comparant le déroulement et l'efficacité de la réponse immune lors d'une primo-infection par un virus grippal et lors d'une ré-infection par le même variant, chez un sujet immunocompétent.
\tcblower
\textbf{Solution détaillée (Modèle de rédaction CER) :}
\textbf{Affirmation :} La réponse immune lors d'une ré-infection par un virus grippal homologue est \textbf{plus rapide, plus intense et plus efficace} que lors de la primo-infection, empêchant souvent l'installation de la maladie clinique.

\textbf{Preuves et Raisonnement :}
Lors de la \textbf{primo-infection}, les \textbf{défenses non spécifiques (1\textsuperscript{re} et 2\textsuperscript{e} lignes)} interviennent en premier : barrières respiratoires (mucus, cils, IFN-$\alpha/\beta$), puis réaction inflammatoire locale et phagocytose par macrophages/neutrophiles, activation des cellules NK. Ces mécanismes, bien qu'immédiats, sont insuffisants pour éliminer totalement le virus. La \textbf{réponse spécifique (3\textsuperscript{e} ligne)} se met en place avec une \textbf{phase de latence de 7 à 10 jours} : présentation de l'antigène viral par les cellules dendritiques aux lymphocytes T CD4$^+$ et CD8$^+$ dans les ganglions drainants, sélection clonale, prolifération et différenciation en effecteurs (LT cytotoxiques tuant les cellules infectées, LB $\rightarrow$ plasmocytes sécrétant des IgM puis IgG/IgA). Le pic viral coïncide souvent avec les symptômes.

Lors de la \textbf{ré-infection}, les \textbf{lymphocytes mémoire B et T} (générés lors de la primo-infection) sont présents en \textbf{nombre élevé} et \textbf{réactifs immédiatement}. Ils ne nécessitent pas de sélection clonale. Dès la reconnaissance de l'antigène viral (présenté par les cellules dendritiques ou reconnu directement par les BCR mémoires), ils prolifèrent et se différencient en \textbf{quelques jours seulement (2-3 jours)}. Les \textbf{anticorps préexistants (IgG sériques, IgA muqueuses)} neutralisent une partie des virions dès l'entrée (neutralisation, opsonisation). La réponse effectrice (LTc, Plasmocytes IgG haute affinité) est massive et précoce. La charge virale est contrôlée avant d'atteindre le seuil symptomatique.

\textbf{Conclusion :} La mémoire immunologique (spécificité + rapidité + amplitude) confère une \textbf{protection stérilisante ou atténuée} lors de la ré-infection, base physiologique de la vaccination.
\end{exoresolu}

\begin{attention}[Les 5 erreurs qui coûtent des points au BEPC]
\begin{enumerate}
 \item \textbf{Confondre "Non spécifique" et "Spécifique" dans les acteurs :} Dire "Les lymphocytes font la phagocytose" ou "Les anticorps sont produits par les macrophages". \textbf{Correction :} Lymphocytes B/T = Spécifique (3\textsuperscript{e} ligne). Phagocytes (PNN, Macrophages) = Non spécifique (2\textsuperscript{e} ligne). Anticorps = Produits par Plasmocytes (LB différenciés).
 \item \textbf{Oublier le rôle de l'opsonisation / du complément :} Dire "Le macrophage mange la bactérie" sans mentionner la reconnaissance facilitée par IgG/C3b. \textbf{Correction :} Préciser "Reconnaissance via récepteurs aux Fc des IgG et au C3b (opsonisation) $\rightarrow$ ingestion facilitée".
 \item \textbf{Mélanger Vaccination (Actif) et Sérothérapie (Passif) :} Écrire "Le vaccin donne des anticorps tout de suite" ou "Le sérum donne une mémoire". \textbf{Correction :} Vaccin = Antigène $\rightarrow$ Réponse active lente + Mémoire (Prophylaxie). Sérum = Anticorps $\rightarrow$ Protection immédiate, pas de mémoire, courte durée (Urgence).
 \item \textbf{Imprécision sur l'inflammation :} Citer les 4 signes (Rubor, Calor, Tumor, Dolor) sans expliquer le \textbf{mécanisme vasculaire} (Vasodilatation $\rightarrow$ Rubor/Calor ; Hyperperméabilité $\rightarrow$ Tumor/Dolor). \textbf{Correction :} Lier systématiquement le signe clinique à sa cause physiologique (Médiateurs $\rightarrow$ Vaisseaux $\rightarrow$ Signe).
 \item \textbf{Dire "Le vaccin tue le microbe" ou "Le vaccin soigne la maladie".} \textbf{Correction :} Le vaccin \textbf{prévient} (prophylaxie). Il \textbf{stimule} le système immunitaire pour qu'il soit prêt. Il ne contient pas d'antibiotiques, ne tue pas directement, ne soigne pas une infection en cours.
\end{enumerate}
\cle{Astuce BEPC :} Relisez votre copie en remplaçant mentalement "système immunitaire" par les acteurs précis (Lymphocytes B, Macrophages, Anticorps, Complément). Si la phrase devient fausse, votre rédaction est trop vague.
\end{attention}

\newpage

\coursec{Exercices}

\courssub{Je vérifie mes connaissances}

\begin{exercice}[QCM : Barrières et inflammation]
Parmi les affirmations suivantes, une seule est exacte. La recopier en la justifiant brièvement.
\begin{enumerate}
    \item La peau est une barrière chimique uniquement.
    \item Le prurit (démangeaison) est un signe cardinal de l'inflammation aiguë.
    \item Les macrophages dérivent des monocytes sanguins et réalisent la phagocytose.
    \item Les lymphocytes B sont les principaux acteurs de l'immunité cellulaire.
\end{enumerate}
\end{exercice}

\begin{exercice}[Vrai ou Faux justifié]
Dire si chaque affirmation est vraie ou fausse, puis justifier votre réponse par une phrase précise.
\begin{enumerate}
    \item Les anticorps détruisent directement les bactéries en les lysant.
    \item La vaccination consiste à injecter des anticorps pour protéger immédiatement l'organisme.
    \item L'inflammation est toujours déclenchée par une infection bactérienne.
    \item La mémoire immunologique permet une réponse plus rapide et plus intense lors d'un second contact avec le même antigène.
\end{enumerate}
\end{exercice}

\begin{exercice}[Définitions et correspondance]
Associer chaque terme de la colonne A à sa définition exacte dans la colonne B.
\begin{center}
\begin{tabularx}{\linewidth}{|l|X|}\hline
\textbf{Colonne A} & \textbf{Colonne B} \\ \hline
1. Phagocytose & A. Protéine de défense spécifique produite par les plasmocytes. \\ \hline
2. Antigène & B. Ensemble des mécanismes de défense présents dès la naissance, non spécifiques. \\ \hline
3. Immunité innée & C. Substance étrangère capable de déclencher une réponse immunitaire spécifique. \\ \hline
4. Anticorps & D. Processus par lequel une cellule ingère et digère une particule solide. \\ \hline
\end{tabularx}
\end{center}
\end{exercice}

\begin{exercice}[Calcul direct : Cinétique anticorps]
Lors d'une primo-infection, le taux d'anticorps (Ac) atteint un maximum de \SI{20}{mg/L} au 14\textsuperscript{e} jour. Lors d'une seconde injection (rappel) du même antigène, le pic atteint \SI{1500}{mg/L} au 5\textsuperscript{e} jour.
Calculer le rapport entre le pic secondaire et le pic primaire. Que traduit ce rapport sur l'efficacité de la mémoire immunologique ?
\end{exercice}

\begin{exercice}[Schéma : Les signes de l'inflammation aiguë]
On observe une zone cutanée enflammée (rouge, chaude, gonflée, douloureuse).
\begin{enumerate}
    \item Citer les quatre signes cardinaux de l'inflammation aiguë (nomenclature de Celse) et le signe fonctionnel associé.
    \item Relier chaque signe observable à son mécanisme physiologique immédiat (vasodilatation, hyperperméabilité vasculaire, compression des nerfs, influx leucocytaire).
    \item Quel rôle jouent les médiateurs chimiques (histamine, prostaglandines) libérés par les mastocytes et les macrophages résidents dans le déclenchement de ces signes ?
\end{enumerate}
\end{exercice}

\courssub{Je m'entraîne}

\begin{exercice}[Exploitation de courbe : Cinétique de la réponse humorale]
On donne la courbe ci-dessous représentant l'évolution du taux d'anticorps sériques (en mg/L) en fonction du temps (en jours) lors d'une primo-infection (courbe continue) et d'une ré-infection (courbe pointillée) par le même antigène.
\begin{popfigure}
\begin{tikzpicture}
\begin{axis}[
    width=\linewidth, height=8cm,
    xlabel={Temps (jours)}, ylabel={Taux d'Ac (mg/L)},
    xmin=0, xmax=35, ymin=0, ymax=1600,
    xtick={0,7,14,21,28,35}, ytick={0,500,1000,1500},
    grid=major, grid style={popline},
    legend pos=north west,
    legend style={font=\small}
]
% Primary response
\addplot[popblue, thick, smooth, samples=100, domain=0:35] 
    ({x}, {100*exp(-(x-14)^2/30) + 5*exp(-(x-14)^2/10)});
\addlegendentry{Primo-infection (Primaire)}
% Secondary response
\addplot[poporange, thick, dashed, smooth, samples=100, domain=0:35] 
    ({x}, {1400*exp(-(x-5)^2/15) + 20*exp(-(x-20)^2/50)});
\addlegendentry{Ré-infection (Secondaire)}
% Annotation lines
\draw[popmuted, dashed] (axis cs:14,0) -- (axis cs:14,100);
\draw[popmuted, dashed] (axis cs:5,0) -- (axis cs:5,1400);
\end{axis}
\end{tikzpicture}
\legende{Cinétique de la réponse humorale primaire et secondaire.}
\end{popfigure}
\begin{enumerate}
    \item Déterminer graphiquement la durée de la phase de latence lors de la primo-infection et lors de la ré-infection.
    \item Quelle est la classe d'immunoglobulines majoritaire lors de la réponse secondaire (IgM ou IgG) ? Justifier.
    \item Calculer l'amplification du pic maximal (facteur d'amplification) entre la réponse secondaire et la réponse primaire.
    \item Expliquer le rôle des lymphocytes B mémoire dans ces différences observées (durée de latence, amplitude, classe d'anticorps).
\end{enumerate}
\end{exercice}

\begin{exercice}[Rédaction courte : Découverte de la phagocytose]
À partir des travaux d'Élie Metchnikoff (1882) sur les larves d'étoile de mer où il a introduit des éclats de roseau, rédiger un court paragraphe (4-5 lignes) expliquant :
\begin{enumerate}
    \item L'observation clé réalisée au microscope.
    \item L'interprétation qu'il en a tirée concernant la défense de l'organisme.
    \item Le nom donné à ce processus et à la cellule responsable.
\end{enumerate}
\end{exercice}

\begin{exercice}[Cas pratique : Prophylaxie du tétanos]
Un enfant de 4 ans, non vacciné contre le tétanos, se blesse au pied avec un clou rouillé enfoui dans la terre (risque majeur de contamination par \emph{Clostridium tetani}). Le médecin prescrit une injection de sérum antitetanique (SAT) et une injection de vaccin anatoxine tétanique (VAT) réalisée sur deux sites d'injection distincts.
\begin{enumerate}
    \item Préciser la nature des substances injectées dans le SAT et dans le VAT.
    \item Justifier l'association des deux injections (SAT + VAT) dans ce contexte d'urgence.
    \item Expliquer pourquoi les deux injections doivent être faites sur des sites distincts.
    \item Quel type d'immunité (active/passive, naturelle/artificielle) est apporté par chacune des injections ?
\end{enumerate}
\end{exercice}

\begin{exercice}[Analyse de document : Étapes de la phagocytose]
Le document ci-dessous présente les étapes successives de la phagocytose d'une bactérie par un macrophage.
\begin{popfigure}
\begin{tikzpicture}[scale=0.85, every node/.style={font=\small}]
% 1. Chemotaxie / Adhérence
\node[draw, popblue, thick, rounded corners, fill=popblue!5, minimum width=3.5cm, minimum height=2.2cm] (etape1) at (0,0) {};
\node[popblue, above=0.1cm of etape1.north] {1. Chimiotaxie \& Adhérence};
\draw[popdark, thick, fill=popdark] (-1.5,0.2) circle (0.15) node[left=0.1cm] {Bactérie};
\draw[popblue, thick, fill=popblue!30] (0.5,-0.2) ellipse (1.2 and 0.6);
\node at (0.5,-0.2) {Macrophage};
\draw[->, popred, thick] (-0.8,0.2) -- (-0.2,0.2);
% 2. Ingestion
\node[draw, poporange, thick, rounded corners, fill=poporange!5, minimum width=3.5cm, minimum height=2.2cm] (etape2) at (5,0) {};
\node[poporange, above=0.1cm of etape2.north] {2. Ingestion (Formation du phagosome)};
\draw[popblue, thick, fill=popblue!30] (5,-0.2) ellipse (1.2 and 0.6);
\draw[popdark, thick, fill=popdark] (5,-0.2) circle (0.15);
\draw[popblue, thick] (4.2,-0.5) .. controls (4.5,-0.8) and (5.5,-0.8) .. (5.8,-0.5);
\draw[popblue, thick] (4.2,0.1) .. controls (4.5,-0.2) and (5.5,-0.2) .. (5.8,0.1);
% 3. Fusion lysosomale
\node[draw, poppurple, thick, rounded corners, fill=poppurple!5, minimum width=3.5cm, minimum height=2.2cm] (etape3) at (10,0) {};
\node[poppurple, above=0.1cm of etape3.north] {3. Fusion lysosomale (Phagolysosome)};
\draw[popblue, thick, fill=popblue!30] (10,-0.2) ellipse (1.2 and 0.6);
\draw[popdark, thick, fill=popdark] (10,-0.2) circle (0.15);
\draw[poppurple, thick, fill=poppurple!50] (9.2,0.5) rectangle (9.6,0.9) node[right] {Granule lysosomal};
\draw[->, poppurple, thick] (9.4,0.7) -- (9.8,0.3);
% 4. Digestion
\node[draw, popgreen, thick, rounded corners, fill=popgreen!5, minimum width=3.5cm, minimum height=2.2cm] (etape4) at (15,0) {};
\node[popgreen, above=0.1cm of etape4.north] {4. Digestion \& Élimination};
\draw[popblue, thick, fill=popblue!30] (15,-0.2) ellipse (1.2 and 0.6);
\draw[popgreen, thick, fill=popgreen!30] (15,-0.2) circle (0.25); % Residus
\node at (15,-0.2) {Ag digérés};
\draw[->, popmuted, thick] (15.8,-0.2) -- (16.5,-0.2) node[right] {Rejet résidus};
% Arrows between steps
\draw[->, thick, popline] (etape1.east) -- (etape2.west);
\draw[->, thick, popline] (etape2.east) -- (etape3.west);
\draw[->, thick, popline] (etape3.east) -- (etape4.west);
\end{tikzpicture}
\legende{Étapes de la phagocytose par un macrophage.}
\end{popfigure}
\begin{enumerate}
    \item Ordonner les légendes suivantes selon la chronologie de la phagocytose : \textit{Fusion avec les lysosomes}, \textit{Chimiotaxie et adhérence}, \textit{Digestion enzymatique}, \textit{Ingestion et formation du phagosome}.
    \item Sur l'étape 3, nommer les organites qui fusionnent avec le phagosome et préciser le contenu de ces organites.
    \item Expliquer pourquoi la phagocytose est considérée comme un mécanisme de l'immunité innée (non spécifique).
\end{enumerate}
\end{exercice}

\courssub{Je me prépare au BEPC}

\begin{exercice}[Sujet type BEPC : La vaccination antigrippale]
La grippe saisonnière est une maladie virale contagieuse. Au Cameroun, la campagne de vaccination cible les personnes à risque (personnes âgées, malades chroniques). Le vaccin utilisé est un vaccin inactivé (virus tués) ou sous-unitaires (protéines de surface HA et NA), administré par voie intramusculaire.
\begin{enumerate}
    \item Rappeler la définition d'un vaccin. Classer ce vaccin antigrippal parmi les types de vaccins connus (vivant atténué, inactivé, sous-unitaires, ARN messager...).
    \item Décrire les grandes étapes de la mise en place de l'immunité humorale adaptative suite à cette injection, depuis la capture de l'antigène par une cellule présentatrice d'antigène (CPA) jusqu'à la production d'anticorps par les plasmocytes. Utiliser les termes : phagocytose, présentation antigénique, activation lymphocytaire, clonage, différenciation.
    \item Expliquer pourquoi il est nécessaire de se faire vacciner \emph{chaque année} contre la grippe, alors qu'une vaccination contre la rougeole (vaccin ROR) confère une protection à vie. (Indices : variabilité antigénique, dérive antigénique, mémoire immunologique).
    \item Une personne vaccinée attrape la grippe 15 jours après l'injection. Cela signifie-t-il que le vaccin a « donné » la grippe ? Justifier scientifiquement.
\end{enumerate}
\end{exercice}

\begin{exercice}[Sujet type BEPC : Sérothérapie et morsure de serpent]
Dans la région de l'Extrême-Nord, un paysan de 35 ans est mordu par un serpent venimeux (vipère). Il est évacué au centre de santé le plus proche. Le médecin administre en urgence un sérum antivenimeux (SAV) polyvalent par perfusion intraveineuse.
\begin{enumerate}
    \item Préciser l'origine de ce sérum (animal donneur, procédure d'obtention).
    \item Quel type d'immunité (active ou passive ? naturelle ou artificielle ?) est conféré par cette injection ? Justifier.
    \item Expliquer le mécanisme d'action des anticorps du sérum sur le venin (neutralisation). Pourquoi l'injection par voie intraveineuse est-elle préférée ici à la voie intramusculaire ?
    \item Ce sérum contient des immunoglobulines hétérologues (d'origine animale). Citer un risque immédiat (choc anaphylactique) et un risque retardé (maladie du sérum) liés à cette nature xénogénique. Comment le médecin prévient-il le risque immédiat avant la perfusion ?
    \item Pourquoi ce traitement ne confère-t-il pas de protection durable contre de futures morsures ?
\end{enumerate}
\end{exercice}

\newpage

\coursec{Corrigés des exercices}

% --- Je vérifie mes connaissances ---

\begin{corrige}[Exercice 1 — QCM : Barrières et inflammation]
\textbf{Bonne réponse :} \og Les macrophages dérivent des monocytes sanguins et réalisent la phagocytose. \fg

\textbf{Justifications :}
\begin{enumerate}
    \item \textbf{Faux.} La peau est une barrière \textbf{physique et mécanique} (stratum corneum kératinisé, jonctions serrées) \textbf{et chimique} (film hydrolipidique acide pH $\approx$ 5,5 ; sébum ; sueur ; peptides antimicrobiens ; flore commensale). Elle n'est pas \emph{uniquement} chimique.
    \item \textbf{Faux.} Le prurit (démangeaison) n'est pas un signe cardinal de l'inflammation aiguë. Les \textbf{quatre signes cardinaux} (nomenclature de Celse, I\textsuperscript{er} siècle) sont : \textbf{Rubor} (rougeur), \textbf{Calor} (chaleur), \textbf{Tumor} (gonflement/œdème), \textbf{Dolor} (douleur). On y adjoint souvent le \textbf{Functio laesa} (perte de fonction). Le prurit est fréquent dans l'inflammation chronique ou allergique (libération d'histamine).
    \item \textbf{Vrai.} Les \textbf{monocytes} (leucocytes sanguins circulants) migrent dans les tissus (diapédèse) et se différencient en \textbf{macrophages} (histiocytes), cellules phagocytaires professionnelles de l'immunité innée, capables d'ingérer et de digérer bactéries, débris cellulaires et particules inertes.
    \item \textbf{Faux.} Les \textbf{lymphocytes B} (et leur forme différenciée, les \textbf{plasmocytes}) sont les acteurs de l'\textbf{immunité humorale} (production d'anticorps). Les \textbf{lymphocytes T cytotoxiques (LT CD8+)} sont les principaux effecteurs de l'\textbf{immunité cellulaire} (destruction directe des cellules infectées ou tumorales).
\end{enumerate}
\end{corrige}

\begin{corrige}[Exercice 2 — Vrai ou Faux justifié]
\begin{enumerate}
    \item \textbf{Faux.} Les anticorps (Ac) ne possèdent pas d'activité lytique directe. Ils neutralisent les toxines/virus (blocage de la fixation cellulaire), opsonisent les bactéries (facilitant la phagocytose), agglutinent/précipitent les antigènes, et \textbf{activent le système du complément} (voie classique) qui forme le \textbf{Complexe d'Attaque Membranaire (MAC)} responsable de la lyse bactérienne.
    \item \textbf{Faux.} La \textbf{vaccination} consiste à administrer un \textbf{antigène} (agent pathogène vivant atténué, inactivé, sous-unitaires, ARNm, vecteur viral) pour induire une \textbf{immunité active} (production propre d'anticorps et de cellules mémoire). L'injection d'anticorps tout formés correspond à la \textbf{sérothérapie} (immunité passive artificielle), dont la protection est immédiate mais transitoire (quelques semaines).
    \item \textbf{Faux.} L'inflammation est une réponse \textbf{non spécifique} (innée) déclenchée par \textbf{toute agression tissulaire} : infection (bactérienne, virale, fongique, parasitaire), traumatisme physique (choc, brûlure, gelure), agression chimique (corrosif, venin), corps étranger, ischémie-nécrose (infarctus). Elle n'est pas réservée aux infections bactériennes.
    \item \textbf{Vrai.} La \textbf{mémoire immunologique} repose sur la persistance de \textbf{lymphocytes B et T mémoire} (clones expansés lors de la primo-réponse). Lors du second contact avec le \emph{même} antigène, ces cellules permettent une réponse \textbf{plus rapide} (phase de latence raccourcie), \textbf{plus intense} (taux d'anticorps plus élevé), \textbf{plus durable} et de \textbf{meilleure affinité} (essentiellement IgG), souvent asymptomatique.
\end{enumerate}
\end{corrige}

\begin{corrige}[Exercice 3 — Définitions et correspondance]
\begin{center}
\begin{tabularx}{\linewidth}{|l|X|l|}\hline
\textbf{Colonne A} & \textbf{Colonne B} & \textbf{Correspondance} \\ \hline
1. Phagocytose & D. Processus par lequel une cellule ingère et digère une particule solide. & \textbf{1 $\rightarrow$ D} \\ \hline
2. Antigène & C. Substance étrangère capable de déclencher une réponse immunitaire spécifique. & \textbf{2 $\rightarrow$ C} \\ \hline
3. Immunité innée & B. Ensemble des mécanismes de défense présents dès la naissance, non spécifiques. & \textbf{3 $\rightarrow$ B} \\ \hline
4. Anticorps & A. Protéine de défense spécifique produite par les plasmocytes. & \textbf{4 $\rightarrow$ A} \\ \hline
\end{tabularx}
\end{center}
\end{corrige}

\begin{corrige}[Exercice 4 — Calcul direct : Cinétique anticorps]
\textbf{Données :}
\begin{itemize}
    \item Pic primaire (primo-infection) : $Ac_1 = \SI{20}{mg/L}$ (au 14\textsuperscript{e} jour).
    \item Pic secondaire (rappel) : $Ac_2 = \SI{1500}{mg/L}$ (au 5\textsuperscript{e} jour).
\end{itemize}

\textbf{Calcul du rapport d'amplification :}
\[
\text{Rapport} = \frac{Ac_2}{Ac_1} = \frac{1500}{20} = 75
\]

\textbf{Interprétation :} Le pic de la réponse secondaire est \textbf{75 fois plus élevé} que celui de la réponse primaire. Ce facteur d'amplification majeur traduit l'efficacité de la \textbf{mémoire immunologique} : la présence de lymphocytes B mémoire (plus nombreux, à seuil d'activation plus bas, déjà commutés en IgG) permet une différenciation massive et rapide en plasmocytes sécréteurs d'anticorps à haute affinité.
\end{corrige}

\begin{corrige}[Exercice 5 — Schéma : Les signes de l'inflammation aiguë]
\begin{enumerate}
    \item \textbf{Quatre signes cardinaux (Celse) + signe fonctionnel :}
    \begin{itemize}
        \item \textbf{Rubor} (Rougeur)
        \item \textbf{Calor} (Chaleur)
        \item \textbf{Tumor} (Gonflement / Œdème)
        \item \textbf{Dolor} (Douleur)
        \item \textbf{Functio laesa} (Perte de fonction / Impuissance fonctionnelle) — signe adjoint.
    \end{itemize}
    \item \textbf{Relation Signe observable $\leftrightarrow$ Mécanisme physiologique immédiat :}
    \begin{center}
    \begin{tabularx}{\linewidth}{|l|X|}\hline
    \textbf{Signe observable} & \textbf{Mécanisme physiologique immédiat} \\ \hline
    Rubor (Rougeur) \& Calor (Chaleur) & \textbf{Vasodilatation artériolaire} locale (augmentation du flux et du calibre sanguin). \\ \hline
    Tumor (Gonflement / Œdème) & \textbf{Hyperperméabilité vasculaire} (écartement jonctions endothéliales $\rightarrow$ sortie plasma/protéines dans l'interstitiel) \textbf{et Influx leucocytaire} (diapédèse des polynucléaires neutrophiles). \\ \hline
    Dolor (Douleur) & \textbf{Compression des terminaisons nerveuses} par l'œdème interstitiel \textbf{et} action directe des \textbf{médiateurs chimiques} (bradykinine, prostaglandines E2, potassium, H+) sur les nocicepteurs. \\ \hline
    Functio laesa (Perte de fonction) & Douleur limitant le mouvement + œdème compressif + infiltration leucocytaire. \\ \hline
    \end{tabularx}
    \end{center}
    \item \textbf{Rôle des médiateurs chimiques (histamine, prostaglandines, leucotriènes, PAF, cytokines TNF-$\alpha$, IL-1) :}
    Libérés précocement par les \textbf{mastocytes} (dégranulation) et les \textbf{macrophages résidents} (synthèse \emph{de novo}) sous l'effet de l'agression (complément, toxines, lésions), ils sont les \textbf{déclencheurs directs} de la réaction vasculaire :
    \begin{itemize}
        \item \textbf{Histamine, prostaglandines (PGD2, PGE2) :} Vasodilatation artériolaire puissante ($\rightarrow$ Rubor, Calor) + augmentation de la perméabilité veineuse post-capillaire ($\rightarrow$ Tumor, exsudat riche en protéines).
        \item \textbf{Chimiotactiques (LT B4, IL-8, C5a) :} Attirent les polynucléaires neutrophiles vers le foyer ($\rightarrow$ Influx leucocytaire, phagocytose).
        \item \textbf{Bradykinine, Prostaglandines E2 :} Sensibilisation/activation des nocicepteurs ($\rightarrow$ Dolor).
    \end{itemize}
\end{enumerate}
\end{corrige}

% --- Je m'entraîne ---

\begin{corrige}[Exercice 6 — Exploitation de courbe : Cinétique de la réponse humorale]
\begin{enumerate}
    \item \textbf{Durée de la phase de latence (lecture graphique) :}
    \begin{itemize}
        \item \textbf{Primo-infection (courbe continue) :} La courbe quitte le niveau basal vers le \textbf{7\textsuperscript{e}--8\textsuperscript{e} jour}. Phase de latence $\approx$ \textbf{7 à 8 jours}.
        \item \textbf{Ré-infection (courbe pointillée) :} La courbe s'élève dès le \textbf{1\textsuperscript{er}--2\textsuperscript{e} jour}. Phase de latence $\approx$ \textbf{1 à 2 jours}.
    \end{itemize}
    \item \textbf{Classe d'immunoglobulines majoritaire en réponse secondaire : \textbf{IgG}.}
    \textbf{Justification :} Lors de la primo-infection, les lymphocytes B naïfs produisent d'abord de l'\textbf{IgM} (pas de commutation de classe initiale). Au cours de la réponse primaire, sous l'effet des cytokines (IL-4, IFN-$\gamma$), s'opère la \textbf{commutation de classe} (class switch) vers IgG, IgA ou IgE. Les \textbf{lymphocytes B mémoire} générés expriment déjà des récepteurs de classe IgG (ou IgA). Lors de la ré-infection, ils se différencient rapidement en plasmocytes sécrétant massivement des \textbf{IgG} (haute affinité, longue demi-vie, passage barrière placentaire, opsonisation forte).
    \item \textbf{Amplification du pic maximal (facteur d'amplification) :}
    Lecture graphique approximative :
    \begin{itemize}
        \item Pic primaire $\approx \SI{100}{mg/L}$ (ordonnée max courbe continue vers jour 14).
        \item Pic secondaire $\approx \SI{1400}{mg/L}$ (ordonnée max courbe pointillée vers jour 5).
    \end{itemize}
    \[
    \text{Facteur d'amplification} = \frac{\text{Pic secondaire}}{\text{Pic primaire}} \approx \frac{1400}{100} = 14
    \]
    (La réponse secondaire est environ \textbf{14 fois plus intense} en concentration sérique maximale).
    \item \textbf{Rôle des lymphocytes B mémoire dans les différences observées :}
    \begin{itemize}
        \item \textbf{Durée de latence raccourcie (1-2 j vs 7-8 j) :} Les LB mémoire sont \textbf{plus nombreux} (fréquence précurseur $\times 10^2$--$10^3$), ont un \textbf{seuil d'activation plus bas} (récepteurs BCR de haute affinité), ne nécessitent pas l'étape de maturation d'affinité ni l'aide T complexe de la primo-activation. Ils répondent quasi immédiatement à l'antigène présenté.
        \item \textbf{Amplitude accrue (facteur 14--75) :} Prolifération clonale massive et rapide (clonal burst) $\rightarrow$ nombre énorme de plasmocytes effecteurs $\rightarrow$ sécrétion massive d'Ac.
        \item \textbf{Classe d'anticorps (IgG vs IgM) :} Les LB mémoire ont déjà subi la \textbf{commutation de classe} (IgM $\rightarrow$ IgG) et la \textbf{maturation d'affinité} (mutations somatiques + sélection) dans les centres germinatifs lors de la primo-réponse. Ils produisent donc directement des \textbf{IgG de haute affinité}.
    \end{itemize}
\end{enumerate}
\end{corrige}

\begin{corrige}[Exercice 7 — Rédaction courte : Découverte de la phagocytose]
En 1882, Élie Metchnikoff observe au microscope, dans des larves d'étoile de mer transparentes, que des \textbf{cellules mobiles amiboïdes (leucocytes)} entourent et \textbf{engloutissent les éclats de roseau} introduits comme corps étrangers. Il interprète cette observation en proposant que ces cellules assurent la \textbf{défense de l'organisme} en « mangeant » les éléments étrangers (bactéries, débris), contestant ainsi la théorie humorale dominante (Pasteur, Koch). Il nomme ce processus \textbf{phagocytose} (du grec \emph{phagein}, manger, et \emph{kytos}, cellule) et la cellule responsable \textbf{phagocyte} (macrophage ou microphage/polynucléaire), fondant ainsi la théorie de l'immunité cellulaire innée.
\end{corrige}

\begin{corrige}[Exercice 8 — Cas pratique : Prophylaxie du tétanos]
\begin{enumerate}
    \item \textbf{Nature des substances injectées :}
    \begin{itemize}
        \item \textbf{SAT (Sérum Antitetanique) :} Solution stérile d'\textbf{immunoglobulines (anticorps) polyclonales} (souvent F(ab')2) \textbf{anti-toxine tétanique}, purifiées à partir du plasma de \textbf{donneurs humains} (immunoglobulines tétaniques humaines, TIG) ou, historiquement, de \textbf{chevaux hyperimmunisés}. C'est un produit \textbf{passif} (anticorps tout faits).
        \item \textbf{VAT (Vaccin Anatoxine Tétanique) :} Suspension de \textbf{toxine tétanique inactivée (anatoxine)} par le formol, adsorbée sur un adjuvant (hydroxyde d'aluminium). C'est un \textbf{antigène} inactivé, non pathogène, mais immunogène.
    \end{itemize}
    \item \textbf{Justification de l'association SAT + VAT (urgence) :}
    Le tétanos a une incubation courte (3 à 21 jours, souvent < 10 j). L'enfant n'est pas vacciné (pas de mémoire).
    \begin{itemize}
        \item \textbf{SAT (Immunité passive immédiate) :} Apporte des \textbf{anticorps neutralisants tout formés} qui diffusent immédiatement dans la circulation et les tissus pour \textbf{neutraliser la toxine tétanique (TeNT)} produite par \emph{Clostridium tetani} au site de la blessure avant qu'elle ne migre vers le SNC. Protection \textbf{immédiate} (demi-vie $\approx$ 3 semaines).
        \item \textbf{VAT (Immunité active différée) :} Stimule le système immunitaire de l'enfant pour qu'il produise \textbf{ses propres anticorps et ses lymphocytes mémoire}. Protection \textbf{apparue vers J10--J14}, durable (années).
    \end{itemize}
    L'association assure une \textbf{couverture continue} : le SAT protège pendant le délai de mise en place de la réponse active induite par le VAT.
    \item \textbf{Sites d'injection distincts (ex: fesse pour SAT, deltoïde pour VAT) :}
    Pour \textbf{éviter la neutralisation locale de l'antigène vaccinal (VAT) par les anticorps du sérum (SAT)}. Si injectés au même endroit, les anticorps du SAT se lieraient à l'anatoxine du VAT, formant des complexes immuns, empêchant sa capture par les CPA et son drainage ganglionnaire, donc \textbf{annulant la réponse active}.
    \item \textbf{Type d'immunité :}
    \begin{itemize}
        \item \textbf{SAT : Immunité \textbf{passive artificielle}} (transfert d'anticorps exogènes).
        \item \textbf{VAT : Immunité \textbf{active artificielle}} (stimulation par un antigène vaccinal).
    \end{itemize}
\end{enumerate}
\end{corrige}

\begin{corrige}[Exercice 9 — Analyse de document : Étapes de la phagocytose]
\begin{enumerate}
    \item \textbf{Ordre chronologique (correspondant au document) :}
    1. \textbf{Chimiotaxie et adhérence} (Étape 1 : migration vers le foyer, reconnaissance/adhérence via récepteurs PRR/opsonines).
    2. \textbf{Ingestion et formation du phagosome} (Étape 2 : pseudopodes, envelopement, intériorisation dans une vésicule = phagosome).
    3. \textbf{Fusion avec les lysosomes} (Étape 3 : phagosome + lysosome $\rightarrow$ phagolysosome).
    4. \textbf{Digestion enzymatique} (Étape 4 : hydrolyse par enzymes acides, élimination des résidus).
    \item \textbf{Étape 3 : Fusion lysosomale.}
    \begin{itemize}
        \item \textbf{Organites :} Les \textbf{lysosomes} (ou granules lysosomaux, vésicules primaires).
        \item \textbf{Contenu :} Un cocktail d'\textbf{enzymes hydrolytiques acides} (optimum pH 4,5--5,0) : \textbf{protéases} (cathepsines), \textbf{nucléases}, \textbf{lipases}, \textbf{glycosidases}, \textbf{phosphatases}, \textbf{lysozyme}, et des peptides antimicrobiens (défensines). Le phagolysosome s'acidifie (pompe H+ ATPase) pour activer ces enzymes.
    \end{itemize}
    \item \textbf{Pourquoi la phagocytose est un mécanisme de l'immunité innée (non spécifique) :}
    \begin{itemize}
        \item \textbf{Récepteurs non clonaux (PRR) :} Les phagocytes utilisent des \textbf{Récepteurs de Reconnaissance de Motifs (PRR)} (ex: TLR, récepteurs au mannose, récepteurs au complément CR1/CR3, récepteurs Fc$\gamma$) qui reconnaissent des \textbf{motifs moléculaires conservés (PAMPs)} communs à de larges classes de pathogènes (LPS, peptidoglycane, flagelline, ADN CpG non méthylé) ou des molécules hôtes altérées (DAMPs).
        \item \textbf{Absence de spécificité antigénique fine :} Un même macrophage phagocyte \emph{E. coli}, \emph{S. aureus}, un champignon ou une particule inerte (charbon, silice) via des récepteurs identiques. Il ne distingue pas des epitopes précis comme le fait un récepteur BCR/TCR.
        \item \textbf{Absence de mémoire :} La phagocytose d'une bactérie n'augmente pas la vitesse ni l'efficacité de la phagocytose d'une bactérie identique ultérieurement (pas de lymphocytes mémoire phagocytaires). La réponse est identique à chaque agression.
        \item \textbf{Présence dès la naissance :} Fonctionnelle immédiatement, sans apprentissage préalable.
    \end{itemize}
\end{enumerate}
\end{corrige}

% --- Je me prépare au BEPC ---

\begin{corrige}[Exercice 10 — Sujet type BEPC : La vaccination antigrippale]
\textbf{Barème indicatif : 10 points}
\begin{enumerate}
    \item \textbf{(2 pts) Définition et classification.}
    \begin{itemize}
        \item \textbf{Définition :} Un vaccin est une \textbf{préparation antigénique} (agent pathogène entier modifié ou fragments purifiés, ou acide nucléique codant l'antigène) administrée pour \textbf{induire une immunité active protectrice} (production d'anticorps neutralisants et de lymphocytes T/B mémoire) sans provoquer la maladie.
        \item \textbf{Classification :} Le vaccin antigrippal décrit (virus tués \textbf{inactivé} ou protéines de surface HA/NA purifiées \textbf{sous-unitaires}) appartient aux \textbf{vaccins inactivés (tués)} ou \textbf{sous-unitaires (sous-unitaires protéiques)}. Ce ne sont \textbf{pas} des vaccins vivants atténués (risque de réversion), ni des vaccins à ARNm (bien que existants pour la grippe, l'énoncé précise « inactivé ou sous-unitaires »).
    \end{itemize}
    \item \textbf{(4 pts) Étapes de la mise en place de l'immunité humorale adaptative.}
    \textbf{Mots-clés imposés à intégrer :} \textit{phagocytose, présentation antigénique, activation lymphocytaire, clonage, différenciation.}
    \begin{enumerate}
        \item \textbf{Phagocytose et traitement de l'antigène :} Au site d'injection (muscle), des \textbf{cellules présentatrices d'antigène professionnelles (CPA : cellules dendritiques, macrophages)} capturent l'antigène vaccinal (HA/NA) par \textbf{phagocytose} (ou endocytose). Dans les endosomes/lysosomes, l'antigène est dégradé en \textbf{peptides antigéniques}.
        \item \textbf{Présentation antigénique :} Les CPA migrent vers le \textbf{ganglion lymphatique drainant}. Elles présentent les peptides antigéniques à leur surface, insérés dans les molécules \textbf{CMH de classe II}, aux \textbf{lymphocytes T auxiliaires (LT CD4+) naïfs} spécifiques (récepteur TCR).
        \item \textbf{Activation lymphocytaire (LT CD4+) :} La reconnaissance du complexe Peptide-CMH II par le TCR, couplée aux signaux co-stimulateurs (CD80/86 - CD28) et aux cytokines (IL-12, IL-6), déclenche l'\textbf{activation} et la prolifération des LT CD4+ en cellules Th (principalement Th2 pour l'immunité humorale).
        \item \textbf{Activation et clonage des lymphocytes B :} Parallèlement, des \textbf{lymphocytes B naïfs} spécifiques de l'antigène natif (HA/NA) le capturent via leur \textbf{BCR} (IgM/IgD membranaire), l'internalisent, le traitent et présentent des peptides sur leur \textbf{CMH II}. Les \textbf{Th2 activés} reconnaissent ce complexe, délivrent des signaux co-stimulateurs (CD40L-CD40) et sécrètent des cytokines (IL-4, IL-21). Cela déclenche l'\textbf{activation}, la prolifération intense (\textbf{clonage} / expansion clonale) des LB spécifiques.
        \item \textbf{Différenciation en plasmocytes :} Les LB clonés se \textbf{différencient} majoritairement en \textbf{plasmocytes} (usines à anticorps : sécrétion massive d'IgM puis IgG/IgA spécifiques de l'HA/NA) et en \textbf{lymphocytes B mémoire} (persistance à vie, réponse secondaire rapide).
    \end{enumerate}
    \item \textbf{(2 pts) Nécessité d'une vaccination annuelle (vs rougeole à vie).}
    \begin{itemize}
        \item \textbf{Grippe : Variabilité antigénique majeure.} Le virus influenza (ARN) possède une \textbf{polymérase sans activité de relecture} $\rightarrow$ taux de mutation élevé $\rightarrow$ \textbf{dérive antigénique} (mutations ponctuelles progressives des glycoprotéines de surface \textbf{Hémagglutinine HA} et \textbf{Neuraminidase NA}). Les anticorps mémoire de l'année précédente (dirigés contre l'ancienne HA/NA) ne reconnaissent plus efficacement les nouvelles souches circulantes $\rightarrow$ \textbf{échappement immunitaire}. De plus, le virus subit des \textbf{réassortiments génomiques (shift antigénique)} lors de co-infections, créant de nouveaux sous-types pandémiques. La composition vaccinale est mise à jour chaque année (OMS) pour matcher les souches circulantes.
        \item \textbf{Rougeole : Stabilité antigénique.} Le virus de la rougeole (Paramyxovirus, ARN mais plus stable, un seul sérotype) ne présente pas de dérive antigénique significative. L'infection naturelle ou le vaccin vivant atténué (ROR) induit une \textbf{immunité active solide et durable} avec des lymphocytes mémoire persistant \textbf{toute la vie} (protection stérilisante).
    \end{itemize}
    \item \textbf{(2 pts) Grippe 15 jours post-vaccination : échec ou maladie vaccinale ?}
    \textbf{Non, le vaccin n'a pas « donné » la grippe.}
    \textbf{Justification scientifique :}
    \begin{itemize}
        \item Le vaccin utilisé est \textbf{inactivé (virus tués) ou sous-unitaires (protéines HA/NA purifiées)}. Il ne contient \textbf{aucun virus vivant, ni réplicatif, ni infectieux}. Il est \textbf{incapable de se multiplier} ou de causer la maladie.
        \item \textbf{Délai de protection :} L'immunité active protectrice (taux d'anticorps neutralisants suffisant) met \textbf{10 à 14 jours} (parfois 3 semaines) à s'installer après l'injection (phase de latence, clonage, différenciation).
        \item \textbf{Explication :} La personne était soit \textbf{déjà incubée} (infectée avant la vaccination, période d'incubation grippe = 1--4 jours), soit \textbf{contaminée pendant la fenêtre de vulnérabilité} (les 15 premiers jours post-vaccination où l'immunité n'est pas encore effective). C'est un \textbf{échec de la prophylaxie par délai trop court}, pas une maladie induite par le vaccin.
    \end{itemize}
\end{enumerate}
\textbf{Points de vigilance :} Ne pas confondre vaccin inactivé/sous-unitaires (non vivant) et vaccin vivant atténué (ROR, BCG, Fièvre jaune) qui, eux, se répliquent faiblement. Bien distinguer \emph{dérive} (mutations progressives) et \emph{shift} (réassortiment brutal). Préciser que l'absence de virus vivant exclut toute réversion virulente.
\end{corrige}

\begin{corrige}[Exercice 11 — Sujet type BEPC : Sérothérapie et morsure de serpent]
\textbf{Barème indicatif : 10 points}
\begin{enumerate}
    \item \textbf{(2 pts) Origine du sérum antivenimeux (SAV) polyvalent.}
    \begin{itemize}
        \item \textbf{Animal donneur :} Principalement le \textbf{cheval} (parfois mouton, chèvre, lapin) pour la grande quantité de plasma disponible.
        \item \textbf{Procédure :} \textbf{Hyperimmunisation} de l'animal par injections répétées (sur plusieurs mois) de \textbf{doses croissantes de venins} (mélange de venins de plusieurs espèces pour le polyvalent : vipères, cobras, mambas, etc. de la région). L'animal développe un taux d'anticorps (IgG) très élevé (titre élevé). Prélèvement sanguin (plasmaphérèse ou saignée) $\rightarrow$ séparation du plasma $\rightarrow$ \textbf{purification des immunoglobulines} (précipitation salin, chromatographie) $\rightarrow$ souvent \textbf{digestion enzymatique (pépine)} pour obtenir des fragments \textbf{F(ab')2} (élimination du fragment Fc responsable de la plupart des réactions anaphylactoïdes et de la maladie du sérum) $\rightarrow$ formulation stérile, contrôle de puissance (neutralisation LD50), conditionnement.
    \end{itemize}
    \item \textbf{(1,5 pt) Type d'immunité conférée.}
    \textbf{Immunité \textbf{passive artificielle}.}
    \textbf{Justification :} \textbf{Passive} car le receveur reçoit des \textbf{anticorps tout formés} (effecteurs immédiats) sans stimulation de son propre système immunitaire (pas d'activation lymphocytaire, pas de clonage, pas de différenciation, \textbf{pas de mémoire}). \textbf{Artificielle} car ces anticorps proviennent d'une \textbf{autre espèce (hétérologues)} et sont administrés par injection (perfusion), non par voie naturelle (transfert maternel).
    \item \textbf{(2,5 pts) Mécanisme d'action et voie IV.}
    \begin{itemize}
        \item \textbf{Mécanisme : Neutralisation.} Les anticorps (F(ab')2) du SAV se lient spécifiquement aux \textbf{toxines du venin} (neurotoxines $\alpha$, phospholipases A2, métalloprotéinases, etc.) par leurs sites de liaison à l'antigène (paratopes). Cela \textbf{bloque leur fixation sur leurs récepteurs cibles} (ex: récepteur nicotinique de la plaque motrice pour les neurotoxines $\alpha$) ou \textbf{inhibe leur activité enzymatique}. Les complexes \textbf{Anticorps-Venin} sont ensuite éliminés par \textbf{phagocytose} (opsonisation via récepteurs Fc sur macrophages) et clairance hépatique/splénique.
        \item \textbf{Voie intraveineuse (IV) privilégiée :} \textbf{Urgent vital} (risque de paralysie respiratoire, hémorragie, choc). La voie IV assure une \textbf{biodisponibilité de 100\% immédiate} (concentration plasmatique thérapeutique instantanée). La voie intramusculaire (IM) est \textbf{proscrite en urgence} : absorption lente, imprévisible (vasoconstriction du choc), douloureuse, volume limité, risque de nécrose locale ; elle retarde l'action neutralisante critique.
    \end{itemize}
    \item \textbf{(2 pts) Risques liés à la nature xénogénique (hétérologue) et prévention.}
    \begin{itemize}
        \item \textbf{Risque immédiat : \textbf{Choc anaphylactoïde / anaphylactique} (Type I, IgE-dépendant ou activation directe du complément/mastocytes).} Survient en minutes. Manifestations : hypotension, bronchospasme, œdème de Quincke, urticaire, arrêt cardiaque. Dû à des IgE préexistantes anti-Ig cheval (exposition antérieure environnementale ou sérum précédent) ou activation non spécifique du complément par les agrégats d'IgG cheval.
        \item \textbf{Risque retardé : \textbf{Maladie du sérum} (Type III, complexes immuns circulants).} Survient \textbf{5 à 10 jours} après. Manifestations : fièvre, arthralgies/arthrite, éruption cutanée, glomérulonéphrite (dépôts complexes Ag-Ac dans glomérules), vascularite. Dû à la formation de complexes immuns entre \textbf{anticorps humains anti-Ig cheval (anti-idiotypiques/anti-allotypiques)} et les \textbf{IgG cheval résiduelles} (ou fragments Fc si IgG entières injectées).
        \item \textbf{Prévention du risque immédiat (protocole standard) :}
        \begin{enumerate}
            \item \textbf{Test cutané (intradermo-réaction - IDR) :} Injection ID de 0,02--0,05 mL de SAV dilué 1/10 ou 1/100 sur l'avant-bras. Lecture à 15--20 min. Papule $\ge$ 5--10 mm + érythème $\rightarrow$ \textbf{Test POSITIF} (risque élevé). Négatif $\rightarrow$ risque faible (mais pas nul).
            \item \textbf{Prémédication systématique :} Antihistaminique H1 (ex: chlorphéniramine) + Corticoïde (ex: dexaméthasone) IV 15--30 min avant.
            \item \textbf{Perfusion lente et fractionnée :} Début très lent (quelques gouttes/min) sous surveillance continue (PA, FC, SpO2, FR) pendant 15--30 min, puis accélération si tolérance.
            \item \textbf{Disponibilité immédiate kit urgence :} Adrénaline IM (0,01 mg/kg), solutés de remplissage, O2, intubation.
        \end{enumerate}
        \textit{Note : L'utilisation de fragments F(ab')2 (sans Fc) réduit drastiquement ces deux risques vs les IgG entières.}
    \end{itemize}
    \item \textbf{(2 pts) Absence de protection durable.}
    L'immunité conférée est \textbf{passive}. Les anticorps injectés (IgG cheval/mouton ou F(ab')2) ont une \textbf{demi-vie plasmatique limitée} ($\approx$ 3--4 semaines pour IgG entières, $\approx$ 3--5 jours pour F(ab')2). Ils sont \textbf{catabolisés} (dégradés) par l'organisme du receveur \textbf{sans stimuler sa propre réponse immune} (pas d'antigène vaccinal, pas d'activation de lymphocytes B naïfs, \textbf{pas de formation de lymphocytes B mémoire ni de plasmocytes à longue durée de vie}). Une fois les anticorps exogènes éliminés (quelques semaines), le taux d'anticorps neutralisants chute sous le seuil protecteur. Le patient redevient \textbf{totalement sensible} à une nouvelle envenimation. Seule une \textbf{vaccination active (anatoxines venimeuses)} — inexistante en routine pour les venins de serpent — conférerait une protection durable.
\end{enumerate}
\textbf{Points de vigilance :} Distinguer clairement \emph{passive} (Ac tout faits, pas de mémoire, transitoire) et \emph{active} (Ag $\rightarrow$ propre réponse, mémoire, durable). Préciser que le SAV neutralise le venin \emph{libre}, il ne répare pas les lésions déjà installées (nécrose, paralysie installée) $\rightarrow$ urgence d'administration. Bien citer le test cutané (IDR) et la prémédication pour la prévention du choc immédiat.
\end{corrige}

\newpage

\coursec{Fiche bilan}

\begin{popfigure}
\begin{center}\tcbox[colback=popdark,colframe=popdark,coltext=white,arc=9pt,boxsep=4pt,left=6pt,right=6pt]{\bfseries\small Réponse immunitaire}\end{center}
\vspace{-2pt}
\begin{tcbraster}[raster columns=3,raster equal height=rows,raster column skip=6pt,raster row skip=6pt,enhanced,blankest]
\begin{tcolorbox}[enhanced,colback=white,colframe=popblue,boxrule=0.9pt,arc=3pt,title={Peau et muqueuses},fonttitle=\bfseries\scriptsize,coltitle=white,colbacktitle=popblue,left=4pt,right=4pt,top=2pt,bottom=2pt,boxsep=1pt,toptitle=1.5pt,bottomtitle=1.5pt,halign=left]
\begin{itemize}[nosep,leftmargin=*,itemsep=1pt,font=\scriptsize]
\item Barrière mécanique
\item Sécrétions (sueur, larmes)
\end{itemize}
\end{tcolorbox}
\begin{tcolorbox}[enhanced,colback=white,colframe=poporange,boxrule=0.9pt,arc=3pt,title={Réaction inflammatoire},fonttitle=\bfseries\scriptsize,coltitle=white,colbacktitle=poporange,left=4pt,right=4pt,top=2pt,bottom=2pt,boxsep=1pt,toptitle=1.5pt,bottomtitle=1.5pt,halign=left]
\begin{itemize}[nosep,leftmargin=*,itemsep=1pt,font=\scriptsize]
\item Rougeur, chaleur, gonflement
\item Afflux de leucocytes
\end{itemize}
\end{tcolorbox}
\begin{tcolorbox}[enhanced,colback=white,colframe=popgreen,boxrule=0.9pt,arc=3pt,title={Phagocytose},fonttitle=\bfseries\scriptsize,coltitle=white,colbacktitle=popgreen,left=4pt,right=4pt,top=2pt,bottom=2pt,boxsep=1pt,toptitle=1.5pt,bottomtitle=1.5pt,halign=left]
\begin{itemize}[nosep,leftmargin=*,itemsep=1pt,font=\scriptsize]
\item Phagocytes
\item Digestion du microbe
\end{itemize}
\end{tcolorbox}
\begin{tcolorbox}[enhanced,colback=white,colframe=poppurple,boxrule=0.9pt,arc=3pt,title={Immunité spécifique},fonttitle=\bfseries\scriptsize,coltitle=white,colbacktitle=poppurple,left=4pt,right=4pt,top=2pt,bottom=2pt,boxsep=1pt,toptitle=1.5pt,bottomtitle=1.5pt,halign=left]
\begin{itemize}[nosep,leftmargin=*,itemsep=1pt,font=\scriptsize]
\item Lymphocytes B : anticorps
\item Lymphocytes T : destruction
\end{itemize}
\end{tcolorbox}
\begin{tcolorbox}[enhanced,colback=white,colframe=poppink,boxrule=0.9pt,arc=3pt,title={Mémoire immunitaire},fonttitle=\bfseries\scriptsize,coltitle=white,colbacktitle=poppink,left=4pt,right=4pt,top=2pt,bottom=2pt,boxsep=1pt,toptitle=1.5pt,bottomtitle=1.5pt,halign=left]
\begin{itemize}[nosep,leftmargin=*,itemsep=1pt,font=\scriptsize]
\item Réponse plus rapide au 2\ieme{} contact
\end{itemize}
\end{tcolorbox}
\begin{tcolorbox}[enhanced,colback=white,colframe=popgold,boxrule=0.9pt,arc=3pt,title={Applications},fonttitle=\bfseries\scriptsize,coltitle=white,colbacktitle=popgold,left=4pt,right=4pt,top=2pt,bottom=2pt,boxsep=1pt,toptitle=1.5pt,bottomtitle=1.5pt,halign=left]
\begin{itemize}[nosep,leftmargin=*,itemsep=1pt,font=\scriptsize]
\item Vaccination
\item Sérothérapie
\end{itemize}
\end{tcolorbox}
\end{tcbraster}

\legende{Carte mentale de la leçon : les trois lignes de défense de l'organisme (barrières naturelles, réaction inflammatoire et phagocytose, immunité spécifique) et leurs applications médicales.}
\end{popfigure}

\begin{bilan}
\textbf{1. Définitions clés}
\begin{itemize}
\item \cle{Immunité} : ensemble des moyens de défense de l'organisme contre les agents étrangers (microbes, virus\ldots).
\item \cle{Antigène} : substance étrangère reconnue par l'organisme, qui déclenche une réponse immunitaire.
\item \cle{Anticorps} : protéine fabriquée par les lymphocytes B, capable de se fixer spécifiquement sur un antigène pour le neutraliser.
\item \cle{Phagocytose} : absorption puis digestion d'un microbe par une cellule appelée \cle{phagocyte}.
\item \cle{Réaction inflammatoire} : réaction locale (rougeur, chaleur, gonflement, douleur) qui accompagne la défense contre une agression.
\item \cle{Lymphocyte} : globule blanc responsable de l'immunité \emph{spécifique}.
\end{itemize}

\textbf{2. Les trois lignes de défense (à connaître par c\oe ur)}
\begin{center}
\begin{tabularx}{\linewidth}{|l|X|X|}
\hline
\rowcolor{popblueL} \textbf{Ligne de défense} & \textbf{Acteurs} & \textbf{Caractère} \\
\hline
1\iere{} ligne : barrières naturelles & Peau, muqueuses, sécrétions (sueur, larmes, mucus) & Non spécifique \\
\hline
2\ieme{} ligne : réaction inflammatoire & Phagocytes (phagocytose) & Non spécifique \\
\hline
3\ieme{} ligne : immunité spécifique & Lymphocytes B (anticorps) et lymphocytes T (destruction des cellules infectées) & Spécifique + mémoire \\
\hline
\end{tabularx}
\end{center}

\textbf{3. Comparaison des deux types de réponse}
\begin{center}
\begin{tabularx}{\linewidth}{|l|X|X|}
\hline
\rowcolor{popgreenL} \textbf{Critère} & \textbf{Réponse non spécifique} & \textbf{Réponse spécifique} \\
\hline
Rapidité & Immédiate & Plus lente (quelques jours) \\
\hline
Cible & Tout agent étranger & Un antigène précis \\
\hline
Mémoire & Non & Oui (2\ieme{} contact plus rapide) \\
\hline
\end{tabularx}
\end{center}

\textbf{4. Applications médicales}
\begin{itemize}
\item \cle{Vaccination} : injection d'un antigène inoffensif $\Rightarrow$ l'organisme fabrique des anticorps et garde la \emph{mémoire} du microbe (ex. vaccins contre la rougeole, le tétanos).
\item \cle{Sérothérapie} : injection d'anticorps déjà fabriqués (sérum) $\Rightarrow$ protection \emph{immédiate mais temporaire} (ex. sérum antivenimeux après morsure de serpent).
\end{itemize}

\textbf{5. Méthode express pour analyser un document}
\begin{enumerate}
\item Identifier la ligne de défense mise en jeu (barrière, phagocytose ou lymphocytes).
\item Dire si la réponse est spécifique ou non, et justifier.
\item Conclure en reliant au rôle protecteur pour l'organisme.
\end{enumerate}
\end{bilan}

\begin{aretenir}[Checklist avant le BEPC]
\begin{itemize}
\item[$\square$] Je sais définir : antigène, anticorps, phagocytose, réaction inflammatoire.
\item[$\square$] Je cite les deux barrières naturelles : la peau et les muqueuses.
\item[$\square$] Je connais les quatre signes de l'inflammation : rougeur, chaleur, gonflement, douleur.
\item[$\square$] Je décris les étapes de la phagocytose (absorption puis digestion du microbe).
\item[$\square$] Je distingue lymphocytes B (anticorps) et lymphocytes T (destruction des cellules infectées).
\item[$\square$] Je sais que l'immunité spécifique possède une mémoire : le 2\ieme{} contact est plus rapide et plus efficace.
\item[$\square$] Je différencie vaccination (protection durable) et sérothérapie (protection immédiate mais temporaire).
\item[$\square$] Je sais classer une défense en « spécifique » ou « non spécifique » avec justification.
\item[$\square$] Je sais rédiger une conclusion claire à partir d'un document ou d'une expérience.
\item[$\square$] Je relie la vaccination aux campagnes de santé publique au Cameroun (rougeole, poliomyélite, tétanos).
\end{itemize}
\end{aretenir}

\findecours

\end{document}
